\chapter{Opérations intrinsèques corrélatives du même continu}
\label{chap:pdfv2-operaciones-intrinsecas-correlativas}

\begin{thesisbox}
Les symboles \(\mathrm{sp},\mathrm{sol},\mathrm{gau},\mathrm{coh}\) et \(\mathrm{det}\) qui apparaissent dans ce chapitre marquent cinq opérations corrélatives sur un unique état APP--TRIT--TPK\@. Chaque opération reçoit la même extension survivante, la même orientation, la même retenue, la même mémoire, la même frontière et le même parcours. Les sections suivantes constituent des étapes mathématiques successives d'une seule construction, non des exposés autonomes.
\end{thesisbox}

\section{Arête génératrice et état généalogique commun}

Nous ne répétons ici ni l'état, ni l'arbre, ni la mesure construits dans le chapitre précédent. Nous prenons \(\widehat x=\widehat x_k(\gamma)\in
\widehat{\mathcal X}^{\rm tot}_k\) de \eqref{eq:continuo-comun-estado} et toutes ses émissions admissibles \(\alpha\in\operatorname{Out}(\widehat x)\). Si \(\widehat y=\alpha_*\widehat x\), nous écrivons
\begin{equation}
\label{eq:pdfv2-cor-hijos-supervivientes}
 \operatorname{Ch}_k(\widehat x):=
 \{(\alpha,\alpha_*\widehat x):
      \alpha\in\operatorname{Out}(\widehat x)\},
 \qquad
 d_k(\widehat x):=\#\operatorname{Ch}_k(\widehat x).
\end{equation}
La finitude et la non-vacuité de cet ensemble sont déjà \eqref{eq:continuo-comun-out-finito} ; la survie à tous les horizons, la multisection complète et sa mesure sont \cref{thm:continuo-comun-total-superviviente,thm:continuo-comun-multiseccion-no-vacia,thm:continuo-comun-medida-canonica}. Aucune seconde non-vacuité n'est supposée ni démontrée à nouveau.

Chaque paire \((\alpha,\widehat y)\) de \eqref{eq:pdfv2-cor-hijos-supervivientes} conserve le registre unique
\begin{equation}
\label{eq:pdfv2-cor-registro-extension}
 \begin{aligned}
 \mathfrak e_\alpha:=\bigl(&
 \alpha,\operatorname{Ext}_{\gamma_\alpha};
 \operatorname{Sh}_{s(\alpha)},\operatorname{Sh}_{t(\alpha)};\\
 &
 \operatorname{res}^{\Sigma,\Pi}(\alpha),
 \operatorname{quo}^{\Sigma,\Pi}(\alpha);\\
 &\operatorname{or}(\alpha),\kappa(\alpha),\mathsf C_\alpha;
 \mathsf Q(s(\alpha)),\mathsf Q(t(\alpha)),g(s(\alpha)),g(t(\alpha)),
 q^{(9)}(s(\alpha)),q^{(9)}(t(\alpha));\\
 &
 m(\gamma\alpha;m_0),\partial\alpha,\gamma\alpha,
 \operatorname{Led}(\gamma\alpha);\\
 &\{\operatorname{Term}_{k+1,N}(\widehat y)\}_{N\geq k+1},
 \operatorname{Msect}(\widehat y),\mu_{\widehat y}\bigr).
 \end{aligned}
\end{equation}
Ici, \(\operatorname{Ext}_{\gamma_\alpha}\) désigne la relation complète déjà satisfaite par l'émission admissible ; elle n'introduit pas un autre filtre. Les opérations ultérieures sont des applications coordonnées de \(\mathfrak e_\alpha\) et ne reconstruisent pas le registre à partir de leurs sorties.

\section{Actions intrinsèques corrélatives sur le générateur commun}

Les marques internes sont fixées sans changer de source : \(\mathrm{sp}\) désigne le transport et la graduation des deux feuillets ; \(\mathrm{sol}\), l'accroissement signé, le bilan et la mémoire ; \(\mathrm{gau}\), l'incidence et la courbure de retenue ; \(\mathrm{coh}\), les ports, la différence rectangulaire et le complexe cellulaire ; et \(\mathrm{det}\), le complexe muni de bases, le remplissage corrélatif et la droite déterminant. Les formules suivantes montrent leurs couplages sur le même \(\mathfrak e_\alpha\) ; leurs actions partagent littéralement source, cible, parcours et registre.

La double évaluation APP fournit, sur le même support, les étiquettes \(\Sigma\) et \(\Pi\). Pour une extension \(\alpha\), leurs multiplicités se calculent à partir du registre complet, sans compteur extérieur :
\begin{equation}
\label{eq:pdfv2-cor-multiplicidades-hojas}
 N_\diamond(\alpha):=
 \sum_{\epsilon\ {\rm en}\ \gamma\alpha}
 \mathbf1_{\{\operatorname{res}^{\diamond}(\epsilon)\neq0\}},
 \qquad \diamond\in\{\Sigma,\Pi\}.
\end{equation}
Soit en outre le coefficient orienté total
\[
 \Delta_{\gamma_\alpha}:=
 \sum_{\epsilon\ {\rm en}\ \gamma\alpha}\operatorname{or}(\epsilon)
 \in\mathbb Z.
\]
L'action signée est
\begin{equation}
\label{eq:pdfv2-cor-b-alpha}
 b_\alpha:=\Delta_{\gamma_\alpha}
 \bigl(N_\Sigma(\alpha)-N_\Pi(\alpha)\bigr),
 \qquad
 b_\alpha^\pm:=\frac{|b_\alpha|\pm b_\alpha}{2}.
\end{equation}
Par construction,
\begin{equation}
\label{eq:pdfv2-cor-b-identidades}
 b_\alpha=b_\alpha^+-b_\alpha^-,\qquad
 |b_\alpha|=b_\alpha^++b_\alpha^-,\qquad
 b_\alpha^+b_\alpha^-=0.
\end{equation}

Le même ensemble d'extensions produit simultanément les ports
\begin{equation}
\label{eq:pdfv2-cor-puertos-comunes}
 I_k(\widehat x_k):=
 \operatorname{Ch}_k(\widehat x_k)\times\{\Sigma\},
 \qquad
 J_k(\widehat x_k):=
 \operatorname{Ch}_k(\widehat x_k)\times\{\Pi\}.
\end{equation}
On ne suppose pas qu'un feuillet visible contienne davantage de prolongements que l'autre : chaque extension survivante est évaluée dans les deux feuillets du même support. Par \eqref{eq:continuo-comun-out-finito}, les deux ensembles de ports sont non vides.

Pour \(A_N=\mathbb Z/3^N\mathbb Z\), l'action cellulaire sur les générateurs est
\begin{equation}
\label{eq:pdfv2-cor-kappa-celda}
 \begin{aligned}
 \kappa_{k,N}:A_N^{I_k}\oplus A_N^{J_k}
 &\longrightarrow A_N^{I_k\times J_k},\\
 \kappa_{k,N}(u,v)_{(\alpha,\Sigma),(\beta,\Pi)}
 &:=u_{(\alpha,\Sigma)}-v_{(\beta,\Pi)}.
 \end{aligned}
\end{equation}
Après choix des premiers ports suivant l'ordre TPK, la différence rectangulaire
\begin{equation}
\label{eq:pdfv2-cor-delta-celda}
 \delta(w)_{ij}:=w_{ij}-w_{ij_0}-w_{i_0j}+w_{i_0j_0}
\end{equation}
réalise la suite scindée
\begin{equation}
\label{eq:pdfv2-cor-sucesion-celda}
 0\longrightarrow A_N\xrightarrow{\Delta}
 A_N^{I_k}\oplus A_N^{J_k}
 \xrightarrow{\kappa_{k,N}}A_N^{I_k\times J_k}
 \xrightarrow{\delta}A_N^{(|I_k|-1)(|J_k|-1)}
 \longrightarrow0.
\end{equation}
C'est une identité de la même fibre APP : elle n'ajoute aucune cellule étrangère au générateur \(\mathfrak e_\alpha\).

L'action cellulaire locale utilise les deux coordonnées entières accumulées par ce même registre. Pour \(\widehat x=\widehat x_k(\gamma)\) et \(\widehat y=\alpha_*\widehat x=\widehat x_{k+1}(\gamma\alpha)\), nous posons
\begin{equation}
\label{eq:pdfv2-cor-c-v}
 \begin{aligned}
 c_{\widehat x}&:=\kappa(\gamma), &
 c_\alpha:=c_{\widehat y}&:=\kappa(\gamma\alpha),\\
 v_{\widehat x}&:=\sum_{\epsilon\ {\rm en}\ \gamma}
                  \operatorname{or}(\epsilon), &
 v_{\widehat y}&:=\sum_{\epsilon\ {\rm en}\ \gamma\alpha}
                  \operatorname{or}(\epsilon),\\
 v_\alpha&:=v_{\widehat y}-v_{\widehat x}
           =\operatorname{or}(\alpha).&&
 \end{aligned}
\end{equation}
Sur les générateurs \(f,r,e,b,g\), on définit
\begin{equation}
\label{eq:pdfv2-cor-complejo-local}
 \partial f=3c_\alpha e+b,\qquad
 \partial r=0,\qquad
 \partial e=g,\qquad
 \partial b=-3c_\alpha g,\qquad
 \partial g=0.
\end{equation}
L'annulation
\begin{equation}
\label{eq:pdfv2-cor-dos-cero}
 \partial^2f=3c_\alpha g-3c_\alpha g=0
\end{equation}
est littérale et utilise la même retenue avec les deux signes imposés par l'orientation.

La graine d'incidence s'obtient à partir du module libre des deux feuillets,
\begin{equation}
\label{eq:pdfv2-cor-v-tpk}
 V_{\rm TPK}:=\mathbf F_3e_\Sigma\oplus\mathbf F_3e_\Pi,
 \qquad
 \mathscr I:=\mathbf P^1(V_{\rm TPK}),
 \qquad |\mathscr I|=4,
\end{equation}
et de ses paires non ordonnées
\begin{equation}
\label{eq:pdfv2-cor-incidencias}
 \mathscr E:=\binom{\mathscr I}{2},
 \qquad |\mathscr E|=6,
 \qquad
 \mathscr B:=\mathscr I\times\{+,-\},
 \qquad |\mathscr B|=8.
\end{equation}
L'application d'incidence est déterminée sur la base par
\begin{equation}
\label{eq:pdfv2-cor-b-incidencia-generadores}
 B[e_L]=\sum_{L'\neq L}e_{\{L,L'\}},
 \qquad L\in\mathscr I,
\end{equation}
et, par linéarité,
\begin{equation}
\label{eq:pdfv2-cor-b-incidencia-coordenadas}
 (Ba)_{\{L,L'\}}=a_L+a_{L'}.
\end{equation}
Si \(s(q)=\sum_{E\in\mathscr E}q_E\), chaque \(L\) apparaît trois fois et
\begin{equation}
\label{eq:pdfv2-cor-sb-cero}
 s(Ba)=0.
\end{equation}
Par conséquent, la fonction
\begin{equation}
\label{eq:pdfv2-cor-wc}
 W_c(q):=\mathbf1_{\{s(q)=0\}}-\frac13
\end{equation}
satisfait \(W_c(q+Ba)=W_c(q)\). Les identités \eqref{eq:pdfv2-cor-v-tpk}--\eqref{eq:pdfv2-cor-wc} sont internes et exactes. Le relèvement concret de \(q\) à partir des cellules et retenues communes est traité dans la couche d'assemblage ; cette identité d'incidence ne se confond pas avec la construction de ce relèvement.

\section{Transport unitaire, polarisation et bilan sur la fibre commune}

La non-vacuité de \eqref{eq:pdfv2-cor-hijos-supervivientes} permet de fixer, sans sélectionner un prolongement,
\begin{equation}
\label{eq:pdfv2-cor-peso-conteo}
 p_k(\alpha,\widehat y\mid\widehat x)
 :=\frac1{d_k(\widehat x)}.
\end{equation}
Soit l'espace des deux feuillets du même état
\begin{equation}
\label{eq:pdfv2-cor-hilbert-hojas}
 \mathscr H_k:=
 \bigoplus_{\widehat x\in\widehat{\mathcal X}^{\rm tot}_k}
 \bigl(\mathsf H_{\widehat x}^{\Sigma}
       \oplus\mathsf H_{\widehat x}^{\Pi}\bigr),
\end{equation}
et soit \(\iota_{\widehat x}\) l'inclusion de la fibre indiquée. Le transport TPK de \eqref{eq:continuo-comun-transporte-hojas} donne l'action exacte
\begin{equation}
\label{eq:pdfv2-cor-u-comun}
 U_k\iota_{\widehat x}v:=
 \sum_{(\alpha,\widehat y)\in\operatorname{Ch}_k(\widehat x)}
 p_k(\alpha,\widehat y\mid\widehat x)^{1/2}
 \iota_{\widehat y}U_\alpha v.
\end{equation}
Les successeurs de deux prédécesseurs distincts sont disjoints puisque la troncature est une fonction. En outre, deux témoins \(\alpha\) distincts ne se confondent pas dans le même facteur direct : le parcours complet \(\gamma\), y compris son dernier pas, est une coordonnée de \(\widehat y\). En conséquence, l'identité de Gram est
\begin{equation}
\label{eq:pdfv2-cor-u-gram}
 \left.U_k^*U_k\right|_{\mathsf H_{\widehat x}^{\Sigma}
       \oplus\mathsf H_{\widehat x}^{\Pi}}
 =
 \sum_{(\alpha,\widehat y)\in\operatorname{Ch}_k(\widehat x)}
 p_k(\alpha,\widehat y\mid\widehat x)\,U_\alpha^*U_\alpha.
\end{equation}
Par conséquent, \(U_k^*U_k=I\) exactement lorsque, dans chaque fibre initiale,
\begin{equation}
\label{eq:pdfv2-cor-gram-promedio-normalizado}
 \sum_{(\alpha,\widehat y)\in\operatorname{Ch}_k(\widehat x)}
 p_k(\alpha,\widehat y\mid\widehat x)\,U_\alpha^*U_\alpha=I.
\end{equation}
L'isométrie de chaque transport de feuillet \(U_\alpha\) est une condition suffisante de cette égalité, conjointement avec la normalisation des poids ; la condition nécessaire et suffisante est la normalisation du Gram moyen exposée dans \eqref{eq:pdfv2-cor-gram-promedio-normalizado}. Pour \(\ell>k\), l'histoire complète est conservée au moyen de
\begin{equation}
\label{eq:pdfv2-cor-u-iterado}
 U_{\ell,k}:=U_{\ell-1}\cdots U_k,
 \qquad U_{k,k}:=I.
\end{equation}

La graduation agit sur les deux coordonnées déjà présentes dans chaque fibre de \eqref{eq:pdfv2-cor-hilbert-hojas} :
\begin{equation}
\label{eq:pdfv2-cor-sigma}
 \sigma_k\iota_{\widehat x}(v_\Sigma,v_\Pi)
 :=\iota_{\widehat x}(v_\Sigma,-v_\Pi),
\qquad
 P_k:=\frac{I+\sigma_k}{2},\quad
 Q_k:=\frac{I-\sigma_k}{2}.
\end{equation}
Le seuil TRIT appartient à l'orientation et à la retenue de \(\widehat x\) ; il n'efface aucun des deux feuillets APP et ne nécessite pas de troisième cas dans \eqref{eq:pdfv2-cor-sigma}.

La partie qui conserve le feuillet et celle qui le change sont deux composantes du même \(U_k\) :
\begin{align}
 A_k&:=P_{k+1}U_kP_k+Q_{k+1}U_kQ_k,
 \label{eq:pdfv2-cor-a-comun}\\
 \eta_k&:=Q_{k+1}U_kP_k+P_{k+1}U_kQ_k.
 \label{eq:pdfv2-cor-eta-comun}
\end{align}
Par orthogonalité de \(P_{k+1}\) et de \(Q_{k+1}\),
\begin{equation}
\label{eq:pdfv2-cor-balance-hojas}
 \begin{aligned}
 U_k&=A_k+\eta_k,\\
 A_k^*A_k+\eta_k^*\eta_k
 &=P_kU_k^*U_kP_k+Q_kU_k^*U_kQ_k.
 \end{aligned}
\end{equation}
\begin{proof}[Démonstration du bilan de Gram par feuillets]
En développant les produits, l'orthogonalité des projecteurs d'arrivée donne
\[
\begin{aligned}
 A_k^*A_k
 &=P_kU_k^*P_{k+1}U_kP_k+Q_kU_k^*Q_{k+1}U_kQ_k,\\
 \eta_k^*\eta_k
 &=P_kU_k^*Q_{k+1}U_kP_k+Q_kU_k^*P_{k+1}U_kQ_k.
\end{aligned}
\]
Leur somme et \(P_{k+1}+Q_{k+1}=I\) prouvent \eqref{eq:pdfv2-cor-balance-hojas}. La différence entre \(U_k^*U_k\) et ce bilan est \(P_kU_k^*U_kQ_k+Q_kU_k^*U_kP_k\) ; elle s'annule exactement lorsque les deux blocs croisés sont nuls. En particulier, la condition \eqref{eq:pdfv2-cor-gram-promedio-normalizado} donne \(U_k^*U_k=I\) et réduit le bilan à \(I\).
\end{proof}
La Proposition~\ref{x_reciprocidad:iv:cont:gram} développe cette même identité générale dans la réalisation ultérieure sur les histoires complètes, dont l'isométrie se démontre en conservant l'étiquette de chaque prolongement. Cette identité et l'accroissement \eqref{eq:pdfv2-cor-b-alpha} ne sont pas deux transports : \(A_k,\eta_k\) agissent sur la base des histoires et \(b_\alpha\) lit l'arête de cette même base.

\subsection{Polarisation signée, télescopie non autonome et retour sans réinitialisation}
On n'introduit pas ici de branche spectrale indépendante. Les notations suivantes sont restreintes au même espace \(\mathscr H_k\) et au même transport \(U_k\) déjà construits
\[
 \mathcal H_k^{\mathrm{sp}}:=\mathscr H_k,
 \qquad
 \Lambda(\gamma_{\widehat x})
 :=\bigl(\operatorname{or},\operatorname{Carr},\operatorname{Mem},
          \partial,\gamma\bigr)(\widehat x),
\]
et \(\Sigma_{\rm reg}\Lambda(\gamma_{\widehat x})\) désigne leur système de troncatures compatibles. Ainsi, le lecteur signé, le défaut et le terminal suivants sont des actions du constructeur corrélatif commun, non une sélection rétrospective d'histoires.
\begingroup
\let\section\subsection
% Extracción quirúrgica del propietario V2 05a: líneas 293--394.
% El resto permanece en este archivo como procedencia, pero no se publica.
\iffalse
\chapter{Macrostructure spectrale--polaire interne du continu}
\label{chap:continuo-macro-sp-interna}

\section{Portée strictement interne}

La construction de ce chapitre se termine avant toute reconnaissance
externe. Ses seules données sont le support APP à deux feuillets, l'orientation
TRIT, les extensions et la retenue TPK, le système projectif d'histoires, la
mesure cylindrique et le ledger signé. Le résultat est une macrostructure
graduée avec transport, expression de changement de feuillet, lecteur signé et
télescopage non autonome. Aucune fonction extérieure n'est utilisée pour définir
l'une quelconque de ses applications.

\section{Domaine survivant de la branche}

Pour \(x_k\in\mathcal A_k\), soit
\begin{equation}
\label{eq:continuo-sp-msect-n}
 \operatorname{Msect}^{\mathrm{sp},(N)}_k(x_k)
 =\left\{(x_j)_{j=k}^{N}:
 \tau_{j+1,j}x_{j+1}=x_j,
 \ (x_j,x_{j+9})\in\Gamma_{9,j},
 \ \text{et le faisceau est conservé }\{\Sigma,\Pi\}\right\}.
\end{equation}
On définit
\begin{equation}
\label{eq:continuo-sp-dominio-superviviente}
 \mathcal C_k^{\mathrm{sp}}
 =\left\{x_k:
 \operatorname{Msect}^{\mathrm{sp},(N)}_k(x_k)\neq\varnothing
 \text{ pour tout }N\geq k\right\},
 \qquad
 \mathcal C_{\rm cont}^{\mathrm{sp}}
 =\varprojlim_k\mathcal C_k^{\mathrm{sp}}.
\end{equation}
Le seuil TRIT n'est pas exclu de ce domaine. Le faisceau conserve les deux feuillets
arithmétiques même si une carte visible n'en publie aucun lors d'un
événement concret.

\section{Paquet de treize coordonnées}

Pour \(x_k\in\mathcal C_k^{\mathrm{sp}}\), le paquet dépendant est
\begin{equation}
\label{eq:continuo-sp-d13}
\begin{aligned}
 D_{\mathrm{sp},k}(x_k)=\bigl(&
 \mathcal F_{\mathrm{sp},k}(x_k),
 \operatorname{Ext}_{\Gamma,\mathrm{sp},k}(x_k),
 \operatorname{Rel}_{\mathrm{sp},k}(x_k),
 \mathcal N_{\mathrm{sp},k}(x_k),
 \operatorname{or}_{\mathrm{sp},k}(x_k),\\
 &\operatorname{Carr}_{\mathrm{sp},k}(x_k),
 \operatorname{Mem}_{\mathrm{sp},k}(x_k),
 \partial_{\mathrm{sp},k}(x_k),
 \gamma_{\mathrm{sp},k}(x_k),
 \operatorname{Msect}_{\mathrm{sp},k}(x_k),\\
 &\operatorname{Surv}_{\mathrm{sp},k}(x_k),
 \operatorname{Term}_{\mathrm{sp},k}(x_k),
 \mathcal L^{\rm sgn}_{\mathrm{sp},k}(x_k)\bigr).
\end{aligned}
\end{equation}
Ses composantes sont déterminées corrélativement comme suit.

La fibre \(\mathcal F_{\mathrm{sp},k}\) contient l'état TPK et le faisceau
ordonné \(\{\Sigma,\Pi\}\). L'extension est la restriction complète de
\(\operatorname{Ext}_\Gamma\) aux histoires qui satisfont
\eqref{eq:continuo-sp-dominio-superviviente}. Les relations contiennent les
lois d'identité et de composition, le relèvement résidu--quotient, le cocycle de
retenue, la compatibilité de frontière et la naturalité de la troncature. La
coordonnée numérique contient profondeur, phase, résidus, quotients et signature
interne. L'orientation stocke séparément feuillet et orientation TRIT. La
mémoire contient le quotient, le ledger et l'incrément de tour. La
frontière et le chemin restent non contractés. La multisection est le système
complet des prolongements de \eqref{eq:continuo-sp-msect-n}. La
survie est l'appartenance stable de
\eqref{eq:continuo-sp-dominio-superviviente}. Le terminal et le lecteur sont
construits dans les sections suivantes.

Les troncatures de toutes les coordonnées forment une application
\(u^{\mathrm{sp}}_{\ell,k}\) qui satisfait
\begin{equation}
\label{eq:continuo-sp-d-naturalidad}
 u^{\mathrm{sp}}_{\ell,k}D_{\mathrm{sp},\ell}
 =D_{\mathrm{sp},k}\tau_{\ell,k}.
\end{equation}

\section{Graphe de la restriction et section dépendante}

On définit
\begin{equation}
\label{eq:continuo-sp-g-graph}
 \mathcal G_{\mathrm{sp},k}
 =\operatorname{Graph}(D_{\mathrm{sp},k})
 =\{(x_k,D_{\mathrm{sp},k}x_k):
     x_k\in\mathcal C_k^{\mathrm{sp}}\}.
\end{equation}
La restriction totale est
\begin{equation}
\label{eq:continuo-sp-res-graph}
 \operatorname{Res}_{\mathrm{sp}}=s_{\mathrm{sp}}:
 \mathcal C_{\rm cont}^{\mathrm{sp}}
 \longrightarrow\mathcal G_{\mathrm{sp}},
 \qquad
 s_{\mathrm{sp}}(x)=(x,D_{\mathrm{sp}}x).
\end{equation}
La projection \(\pi_1(x,d)=x\) vérifie
\begin{equation}
\label{eq:continuo-sp-res-inversa}
 \pi_1s_{\mathrm{sp}}=I,
 \qquad
 s_{\mathrm{sp}}\pi_1=I_{\mathcal G_{\mathrm{sp}}}.
\end{equation}

Soit \(r_{\mathrm{sp},k}(x_k)\) la multisection des feuillets et des chemins qui appartient
déjà à \(D_{\mathrm{sp},k}(x_k)\). On définit
\begin{equation}
\label{eq:continuo-sp-x-dependiente}
 X_{\chi,\mathrm{sp},k}
 =\{(r_{\mathrm{sp},k}x_k,x_k,D_{\mathrm{sp},k}x_k):
     x_k\in\mathcal C_k^{\mathrm{sp}}\},
\end{equation}
\begin{equation}
\label{eq:continuo-sp-pr-x}
 \operatorname{pr}_{X,\mathrm{sp}}(x,D_{\mathrm{sp}}x)
 =(r_{\mathrm{sp}}x,x,D_{\mathrm{sp}}x).
\end{equation}
Oublier la première coordonnée de
\eqref{eq:continuo-sp-x-dependiente} est l'inverse de
\(\operatorname{pr}_{X,\mathrm{sp}}\).

\section{Faisceau de feuillets et mesure projective}

Soit
\begin{equation}
\label{eq:continuo-sp-haz-hojas}
 \widetilde{\mathcal C}_k^{\mathrm{sp}}
 =\{(x_k,h):x_k\in\mathcal C_k^{\mathrm{sp}},
              h\in\{\Sigma,\Pi\}\}.
\end{equation}
La mesure projective du continu induit des probabilités
\(\widetilde\mu_k\) sur ce faisceau. On se restreint au support positif ; pour
\(z\) dans ce support,
\begin{equation}
\label{eq:continuo-sp-consistencia-medida}
 \widetilde\mu_k(z)
 =\sum_{\widetilde\tau_{k+1,k}(z')=z}
  \widetilde\mu_{k+1}(z').
\end{equation}
La linéarisation de niveau est
\begin{equation}
\label{eq:continuo-sp-hilbert-nivel}
 \mathcal H_k^{\mathrm{sp}}
 =\ell^2(\operatorname{supp}\widetilde\mu_k),
\end{equation}
de base orthonormale \(e_z\).

\section{Transport de toutes les extensions et retenue}

Pour \(z'\) qui prolonge \(z\), nous définissons le poids conditionnel positif
\begin{equation}
\label{eq:continuo-sp-peso-condicional}
 w_k(z'\mid z)
 =\left(\frac{\widetilde\mu_{k+1}(z')}
               {\widetilde\mu_k(z)}\right)^{1/2}.
\end{equation}
Le transport est
\begin{equation}
\label{eq:continuo-sp-u-accion}
 U_ke_z
 =\sum_{\substack{z':\widetilde\tau_{k+1,k}(z')=z\\
            (z,z')\in\operatorname{Ext}_{\Gamma,\mathrm{sp},k}}}
 w_k(z'\mid z)e_{z'}.
\end{equation}
La somme contient toutes les extensions admissibles. Elle ne dépend pas d'un sélecteur.

Par \eqref{eq:continuo-sp-consistencia-medida},
\[
 \sum_{z':\widetilde\tau(z')=z}|w_k(z'\mid z)|^2=1.
\]
Les ensembles de descendants de deux parents distincts sont disjoints. Par
conséquent,
\begin{equation}
\label{eq:continuo-sp-u-isometria}
 U_k^*U_k=I_{\mathcal H_k^{\mathrm{sp}}}.
\end{equation}

Chaque base \(e_{z'}\) inclut la retenue de l'extension correspondante. Pour
des chemins composables, on conserve
\begin{align}
 \kappa(\gamma_2\gamma_1)
 &=\kappa(\gamma_2)+
   \mathsf C(\gamma_2)\kappa(\gamma_1),
 \label{eq:continuo-sp-carry-nativo}\\
 \operatorname{Carr}(\gamma_2\gamma_1)
 &=\operatorname{Carr}(\gamma_2)H(\gamma_1)
  +Y(\gamma_2)\operatorname{Carr}(\gamma_1).
 \label{eq:continuo-sp-carry-operador}
\end{align}
Les poids conditionnels se multiplient le long de l'unique chaîne de
troncatures d'un descendant. Les équations
\eqref{eq:continuo-sp-carry-nativo}--
\eqref{eq:continuo-sp-carry-operador} garantissent que les étiquettes de retenue
coïncident. Par conséquent,
\begin{equation}
\label{eq:continuo-sp-u-composicion}
 U_{\ell,k}=U_{\ell-1}\cdots U_k.
\end{equation}

\section{Graduation, projecteurs et décomposition du transport}

La graduation de feuillet est l'involution
\begin{equation}
\label{eq:continuo-sp-sigma}
 \sigma_ke_{(x,\Sigma)}=+e_{(x,\Sigma)},
 \qquad
 \sigma_ke_{(x,\Pi)}=-e_{(x,\Pi)}.
\end{equation}
Ainsi,
\[
 \sigma_k^*=\sigma_k,
 \qquad
 \sigma_k^2=I.
\]
Les projecteurs internes sont
\begin{equation}
\label{eq:continuo-sp-p-q}
 P_k=\frac{I+\sigma_k}{2},
 \qquad
 Q_k=\frac{I-\sigma_k}{2}.
\end{equation}
Elles satisfont
\begin{equation}
\label{eq:continuo-sp-p-q-relaciones}
 P_k^2=P_k,
 \quad Q_k^2=Q_k,
 \quad P_kQ_k=0,
 \quad P_k+Q_k=I.
\end{equation}

Nous séparons la partie du transport qui conserve le feuillet de celle qui le change :
\begin{align}
 A_k&=P_{k+1}U_kP_k+Q_{k+1}U_kQ_k,
 \label{eq:continuo-sp-a}\\
 \eta_k&=Q_{k+1}U_kP_k+P_{k+1}U_kQ_k.
 \label{eq:continuo-sp-eta}
\end{align}
Alors
\begin{equation}
\label{eq:continuo-sp-u-a-eta}
 U_k=A_k+\eta_k,
\end{equation}
et la graduation vérifie
\begin{equation}
\label{eq:continuo-sp-paridad-a-eta}
 \sigma_{k+1}A_k=A_k\sigma_k,
 \qquad
 \sigma_{k+1}\eta_k=-\eta_k\sigma_k.
\end{equation}

\begin{proposition}[Bilan local des feuillets]
\label{prop:continuo-sp-balance-local}
Pour chaque niveau,
\begin{equation}
\label{eq:continuo-sp-balance-local}
 A_k^*A_k+\eta_k^*\eta_k=I.
\end{equation}
\end{proposition}

\begin{proof}
L'orthogonalité de \(P_{k+1}\) et \(Q_{k+1}\) donne
\begin{align*}
 A_k^*A_k
 &=P_kU_k^*P_{k+1}U_kP_k
   +Q_kU_k^*Q_{k+1}U_kQ_k,\\
 \eta_k^*\eta_k
 &=P_kU_k^*Q_{k+1}U_kP_k
   +Q_kU_k^*P_{k+1}U_kQ_k.
\end{align*}
En sommant et en utilisant \(P_{k+1}+Q_{k+1}=I\),
\[
 A_k^*A_k+\eta_k^*\eta_k
 =P_kU_k^*U_kP_k+Q_kU_k^*U_kQ_k
 =P_k+Q_k=I.
\]
\end{proof}

L'opérateur \(\eta_k\) enregistre l'abandon du canal qui conserve le feuillet. Il
n'est pas la retenue entière : les coordonnées
\eqref{eq:continuo-sp-carry-nativo}--
\eqref{eq:continuo-sp-carry-operador} restent dans le paquet
\eqref{eq:continuo-sp-d13}.

\fi
\section{Lecteur signé avec registre généalogique}

Pour \(v\in\mathcal H_k^{\mathrm{sp}}\), la polarisation signée est
\begin{equation}
\label{eq:continuo-sp-polarizacion-firmada}
 \operatorname{pol}_k(v)
 =\langle v,\sigma_kv\rangle
 =\|P_kv\|^2-\|Q_kv\|^2.
\end{equation}
Le lecteur spécialisé conserve en outre la source :
\begin{equation}
\label{eq:continuo-sp-lector-firmado}
 \mathcal L^{\rm sgn}_{\mathrm{sp},k}(x_k;v)
 =\bigl(
   \Lambda(\gamma_{x_k}),
   \Sigma_{\rm reg}(\Lambda(\gamma_{x_k})),
   \|P_kv\|^2,
   \|Q_kv\|^2,
   \operatorname{pol}_k(v)\bigr).
\end{equation}
L'orientation TRIT est conservée dans \(\Lambda(\gamma_{x_k})\) ; elle ne
s'identifie pas au signe de feuillet. À l'horizon \(N\), le registre garde la
famille des lecteurs pour tous les niveaux \(k\leq N\). Tronquer signifie
supprimer uniquement les entrées ultérieures.

\section{Produit non autonome et terminal}

On définit
\begin{equation}
\label{eq:continuo-sp-f-producto}
 F_{0,0}=I_{\mathcal H_0^{\mathrm{sp}}},
 \qquad
 F_{0,k}=A_{k-1}\cdots A_1A_0
 \quad(k\geq1).
\end{equation}
Par conséquent,
\begin{equation}
\label{eq:continuo-sp-f-recursion}
 F_{0,k+1}=A_kF_{0,k}.
\end{equation}
Le terme publié lors de l'abandon du canal à l'étape \(k\) et le terminal à
l'horizon \(N\) sont
\begin{equation}
\label{eq:continuo-sp-defecto-terminal}
 \mathcal D_k=F_{0,k}^*\eta_k^*\eta_kF_{0,k},
 \qquad
 \mathcal T_N=F_{0,N}^*F_{0,N}.
\end{equation}

\begin{theorem}[Télescopage spectral--polaire non autonome]
\label{thm:continuo-sp-telescopia-no-autonoma}
Pour tout horizon fini \(N\),
\begin{equation}
\label{eq:continuo-sp-telescopia-no-autonoma}
 \boxed{
 I=\sum_{k=0}^{N-1}
 F_{0,k}^*\eta_k^*\eta_kF_{0,k}
 +F_{0,N}^*F_{0,N}.}
\end{equation}
De plus,
\[
 0\preceq\mathcal T_{N+1}
 \preceq\mathcal T_N\preceq I,
\]
de sorte qu'il existe la limite forte positive
\[
 \mathcal T_\infty
 =\operatorname{s-lim}_{N\to\infty}\mathcal T_N
\]
et
\begin{equation}
\label{eq:continuo-sp-telescopia-infinita}
 I=\sum_{k=0}^{\infty}
 F_{0,k}^*\eta_k^*\eta_kF_{0,k}
 +\mathcal T_\infty
\end{equation}
en topologie forte. Le théorème n'affirme pas que
\(\mathcal T_\infty\) soit nul.
\end{theorem}

\begin{proof}
En multipliant \eqref{eq:pdfv2-cor-balance-hojas} par
\(F_{0,k}^*\) et \(F_{0,k}\), et en utilisant
\eqref{eq:continuo-sp-f-recursion}, on obtient
\[
 F_{0,k}^*F_{0,k}
 =F_{0,k+1}^*F_{0,k+1}
  +F_{0,k}^*\eta_k^*\eta_kF_{0,k}.
\]
La somme de \(k=0\) à \(N-1\) annule tous les termes
intermédiaires et laisse précisément
\eqref{eq:continuo-sp-telescopia-no-autonoma}. Comme
\(A_k^*A_k\preceq I\),
\[
 \mathcal T_{N+1}
 =F_{0,N}^*A_N^*A_NF_{0,N}
 \preceq F_{0,N}^*F_{0,N}=\mathcal T_N.
\]
La monotonie et la borne positive fournissent la limite forte. Le passage à la
limite dans l'identité finie donne
\eqref{eq:continuo-sp-telescopia-infinita}.
\end{proof}

\iffalse
\section{Assembleur comme graphe}

Pour
\[
 z_N=(r_{\mathrm{sp},N}x_N,x_N,D_{\mathrm{sp},N}x_N)
 \in X_{\chi,\mathrm{sp},N},
\]
on définit l'assembleur matériel
\begin{equation}
\label{eq:continuo-sp-a-accion}
\begin{aligned}
 a_{\mathrm{sp},N}(z_N)=\bigl(&
 (\mathcal H_k^{\mathrm{sp}},\sigma_k,P_k,Q_k)_{k\leq N},
 (U_k,A_k,\eta_k)_{k<N},\\
 &(F_{0,k},\mathcal D_k,\mathcal T_k,
   \mathcal L^{\rm sgn}_{\mathrm{sp},k})_{k\leq N}\bigr).
\end{aligned}
\end{equation}
La macrostructure est
\begin{equation}
\label{eq:continuo-sp-m-graph}
 \mathfrak M_{\mathrm{sp},N}
 =\operatorname{Graph}(a_{\mathrm{sp},N}),
 \qquad
 \mathcal A_{\mathrm{sp},N}(z_N)
 =(z_N,a_{\mathrm{sp},N}z_N).
\end{equation}
La projection sur \(z_N\) est l'inverse de
\(\mathcal A_{\mathrm{sp},N}\). La troncature supprime les entrées dont
l'indice dépasse le nouvel horizon et vérifie
\begin{equation}
\label{eq:continuo-sp-a-naturalidad}
 m^{\mathrm{sp}}_{\ell,k}a_{\mathrm{sp},\ell}
 =a_{\mathrm{sp},k}\xi^{\mathrm{sp}}_{\ell,k}.
\end{equation}

\section{Publicateur comme graphe}

Le publicateur ne contracte pas la macrostructure en un scalaire. Il renvoie le
certificat complet
\begin{equation}
\label{eq:continuo-sp-p-accion}
 p_{\mathrm{sp},N}(m_N)
 =\left(
 \mathcal L^{\rm sgn}_{\mathrm{sp},k},
 \mathcal D_k,
 \mathcal T_k,
 I=\sum_{j=0}^{k-1}\mathcal D_j+\mathcal T_k
 \right)_{0\leq k\leq N}.
\end{equation}
On définit
\begin{equation}
\label{eq:continuo-sp-c-graph}
 \mathcal C_{\mathrm{sp},N}
 =\operatorname{Graph}(p_{\mathrm{sp},N}),
 \qquad
 \mathcal P_{\mathrm{sp},N}(m_N)
 =(m_N,p_{\mathrm{sp},N}m_N).
\end{equation}
La projection sur \(m_N\) est son inverse. Tronquer le certificat supprime seulement
les identités ultérieures, de sorte que
\begin{equation}
\label{eq:continuo-sp-p-naturalidad}
 c^{\mathrm{sp}}_{\ell,k}p_{\mathrm{sp},\ell}
 =p_{\mathrm{sp},k}m^{\mathrm{sp}}_{\ell,k}.
\end{equation}

\section{Théorème de la chaîne interne complète}

\begin{theorem}[Macrostructure spectrale--polaire engendrée par le continu]
\label{thm:continuo-sp-cadena-interna}
Les actions
\eqref{eq:continuo-sp-res-graph},
\eqref{eq:continuo-sp-pr-x},
\eqref{eq:continuo-sp-m-graph} et
\eqref{eq:continuo-sp-c-graph} forment la chaîne naturelle
\begin{equation}
\label{eq:continuo-sp-cadena-interna}
 \mathcal C_{\rm cont}^{\rm disc}
 \supset\mathcal C_{\rm cont}^{\mathrm{sp}}
 \xrightarrow[\cong]{\operatorname{Res}_{\mathrm{sp}}}
 \mathcal G_{\mathrm{sp}}
 \xrightarrow[\cong]{\operatorname{pr}_{X,\mathrm{sp}}}
 X_{\chi,\mathrm{sp}}
 \xrightarrow[\cong]{\mathcal A_{\mathrm{sp}}}
 \mathfrak M_{\mathrm{sp}}
 \xrightarrow[\cong]{\mathcal P_{\mathrm{sp}}}
 \mathcal C_{\mathrm{sp}}.
\end{equation}
Son contenu publié est le faisceau à deux feuillets, la graduation
\(\sigma_k\), les projecteurs \(P_k,Q_k\), le transport total \(U_k\),
la séparation \(U_k=A_k+\eta_k\), la retenue et la mémoire, le lecteur signé
et le télescopage
\eqref{eq:continuo-sp-telescopia-no-autonoma} avec terminal explicite. Chaque
objet conserve une projection canonique vers tous ses antécédents.
\end{theorem}

\begin{proof}
Les identités
\eqref{eq:continuo-sp-res-inversa} prouvent la première équivalence. La
définition dépendante \eqref{eq:continuo-sp-x-dependiente} prouve la
deuxième. Les graphes
\eqref{eq:continuo-sp-m-graph} et
\eqref{eq:continuo-sp-c-graph} prouvent les deux restantes. La naturalité des
quatre étapes est
\eqref{eq:continuo-sp-d-naturalidad},
\eqref{eq:continuo-sp-a-naturalidad} et
\eqref{eq:continuo-sp-p-naturalidad}, avec la naturalité de
\(r_{\mathrm{sp},k}\). Le contenu opératoriel est prouvé par
\eqref{eq:continuo-sp-u-isometria}, le
\cref{prop:continuo-sp-balance-local} et le
\cref{thm:continuo-sp-telescopia-no-autonoma}. Toutes les extensions apparaissent
dans \eqref{eq:continuo-sp-u-accion} ; aucun choix de
chemin n'intervient donc.
\end{proof}

\section{Retour de phase et évolution de l'état}

La branche conserve expressément
\begin{equation}
\label{eq:continuo-sp-no-reset}
 g(k+9)=g(k),
 \qquad
 x_{k+9}\neq x_k\quad\text{en général}.
\end{equation}
Le second membre enregistre l'accroissement de mémoire et les changements possibles
de feuillet, de retenue, de cylindre et de frontière. Ni \(\sigma_k\), ni les projecteurs, ni
le lecteur signé ne contractent cette différence.

\section{Provenance}

Les deux feuillets APP, l'état TRIT enrichi, les fibres et extensions TPK,
la retenue, le retour sans réinitialisation, la survie et le ledger sont une
\texttt{ARCHITECTURE\_AUTORALE\_PRÉEXISTANTE}. L'obligation de conserver le
terminal fini est un \texttt{RÉSULTAT\_RÉCUPÉRÉ}. Le faisceau explicite de feuillets,
l'involution \(\sigma_k\), les projecteurs, la linéarisation qui somme toutes
les extensions, la décomposition \(U_k=A_k+\eta_k\), le produit non
autonome \(F_{0,k}\) et les graphes dépendants sont une
\texttt{FORMALISATION\_NOUVELLE}. Le
\cref{thm:continuo-sp-telescopia-no-autonoma,thm:continuo-sp-cadena-interna}
constitue le \texttt{CERTIFICAT\_NOUVEAU} de cette formalisation.

\section{Falsadores internos}

\begin{proposition}[Falsificateurs de la macrostructure]
\label{prop:continuo-sp-falsadores}
La conclusion du
\cref{thm:continuo-sp-cadena-interna} n'est pas transférable si l'une quelconque
des substitutions suivantes se produit :
\begin{enumerate}
 \item choisir un prolongement dans
       \eqref{eq:continuo-sp-u-accion} et effacer les autres ;
 \item remplacer \(\mathcal G_{\mathrm{sp}}\) par l'image extérieure de
       \(D_{\mathrm{sp}}\) ;
 \item effacer \((x,D_{\mathrm{sp}}x)\) lors de la formation de
       \(X_{\chi,\mathrm{sp}}\) ;
 \item identifier le feuillet à l'orientation TRIT ou effacer le seuil ;
 \item rompre la cohérence
       \eqref{eq:continuo-sp-consistencia-medida} ;
 \item effacer la retenue ou enfreindre
       \eqref{eq:continuo-sp-carry-nativo}--
       \eqref{eq:continuo-sp-carry-operador} ;
 \item identifier \(\eta_k\) à toute la retenue ;
 \item remplacer la famille \((A_k)_k\) par un seul opérateur sans démontrer
       la stationnarité ;
 \item supprimer \(F_{0,N}^*F_{0,N}\) à un horizon fini ;
 \item affirmer \(\mathcal T_\infty=0\) sans preuve supplémentaire ;
 \item reconstruire une histoire à partir de la polarisation signée ;
 \item utiliser \(g(k+9)=g(k)\) pour affirmer \(x_{k+9}=x_k\) ;
 \item introduire un domaine, une fonction ou un critère extérieur pour
       définir l'une quelconque des flèches internes ;
 \item introduire une sortie ultérieure comme antécédent de la branche.
\end{enumerate}
\end{proposition}
\fi

% Extracción quirúrgica del propietario V2 05a: líneas 512--523.
% El resto permanece en este archivo como procedencia, pero no se publica.
\iffalse
\chapter{Macrostructure spectrale--polaire interne du continu}
\label{chap:continuo-macro-sp-interna}

\section{Portée strictement interne}

La construction de ce chapitre se termine avant toute reconnaissance
externe. Ses seules données sont le support APP à deux feuillets, l'orientation
TRIT, les extensions et la retenue TPK, le système projectif d'histoires, la
mesure cylindrique et le ledger signé. Le résultat est une macrostructure
graduée avec transport, expression de changement de feuillet, lecteur signé et
télescopage non autonome. Aucune fonction extérieure n'est utilisée pour définir
l'une quelconque de ses applications.

\section{Domaine survivant de la branche}

Pour \(x_k\in\mathcal A_k\), soit
\begin{equation}
\label{eq:continuo-sp-msect-n}
 \operatorname{Msect}^{\mathrm{sp},(N)}_k(x_k)
 =\left\{(x_j)_{j=k}^{N}:
 \tau_{j+1,j}x_{j+1}=x_j,
 \ (x_j,x_{j+9})\in\Gamma_{9,j},
 \ \text{et le faisceau est conservé }\{\Sigma,\Pi\}\right\}.
\end{equation}
On définit
\begin{equation}
\label{eq:continuo-sp-dominio-superviviente}
 \mathcal C_k^{\mathrm{sp}}
 =\left\{x_k:
 \operatorname{Msect}^{\mathrm{sp},(N)}_k(x_k)\neq\varnothing
 \text{ pour tout }N\geq k\right\},
 \qquad
 \mathcal C_{\rm cont}^{\mathrm{sp}}
 =\varprojlim_k\mathcal C_k^{\mathrm{sp}}.
\end{equation}
Le seuil TRIT n'est pas exclu de ce domaine. Le faisceau conserve les deux feuillets
arithmétiques même si une carte visible n'en publie aucun lors d'un
événement concret.

\section{Paquet de treize coordonnées}

Pour \(x_k\in\mathcal C_k^{\mathrm{sp}}\), le paquet dépendant est
\begin{equation}
\label{eq:continuo-sp-d13}
\begin{aligned}
 D_{\mathrm{sp},k}(x_k)=\bigl(&
 \mathcal F_{\mathrm{sp},k}(x_k),
 \operatorname{Ext}_{\Gamma,\mathrm{sp},k}(x_k),
 \operatorname{Rel}_{\mathrm{sp},k}(x_k),
 \mathcal N_{\mathrm{sp},k}(x_k),
 \operatorname{or}_{\mathrm{sp},k}(x_k),\\
 &\operatorname{Carr}_{\mathrm{sp},k}(x_k),
 \operatorname{Mem}_{\mathrm{sp},k}(x_k),
 \partial_{\mathrm{sp},k}(x_k),
 \gamma_{\mathrm{sp},k}(x_k),
 \operatorname{Msect}_{\mathrm{sp},k}(x_k),\\
 &\operatorname{Surv}_{\mathrm{sp},k}(x_k),
 \operatorname{Term}_{\mathrm{sp},k}(x_k),
 \mathcal L^{\rm sgn}_{\mathrm{sp},k}(x_k)\bigr).
\end{aligned}
\end{equation}
Ses composantes sont déterminées corrélativement comme suit.

La fibre \(\mathcal F_{\mathrm{sp},k}\) contient l'état TPK et le faisceau
ordonné \(\{\Sigma,\Pi\}\). L'extension est la restriction complète de
\(\operatorname{Ext}_\Gamma\) aux histoires qui satisfont
\eqref{eq:continuo-sp-dominio-superviviente}. Les relations contiennent les
lois d'identité et de composition, le relèvement résidu--quotient, le cocycle de
retenue, la compatibilité de frontière et la naturalité de la troncature. La
coordonnée numérique contient profondeur, phase, résidus, quotients et signature
interne. L'orientation stocke séparément feuillet et orientation TRIT. La
mémoire contient le quotient, le ledger et l'incrément de tour. La
frontière et le chemin restent non contractés. La multisection est le système
complet des prolongements de \eqref{eq:continuo-sp-msect-n}. La
survie est l'appartenance stable de
\eqref{eq:continuo-sp-dominio-superviviente}. Le terminal et le lecteur sont
construits dans les sections suivantes.

Les troncatures de toutes les coordonnées forment une application
\(u^{\mathrm{sp}}_{\ell,k}\) qui satisfait
\begin{equation}
\label{eq:continuo-sp-d-naturalidad}
 u^{\mathrm{sp}}_{\ell,k}D_{\mathrm{sp},\ell}
 =D_{\mathrm{sp},k}\tau_{\ell,k}.
\end{equation}

\section{Graphe de la restriction et section dépendante}

On définit
\begin{equation}
\label{eq:continuo-sp-g-graph}
 \mathcal G_{\mathrm{sp},k}
 =\operatorname{Graph}(D_{\mathrm{sp},k})
 =\{(x_k,D_{\mathrm{sp},k}x_k):
     x_k\in\mathcal C_k^{\mathrm{sp}}\}.
\end{equation}
La restriction totale est
\begin{equation}
\label{eq:continuo-sp-res-graph}
 \operatorname{Res}_{\mathrm{sp}}=s_{\mathrm{sp}}:
 \mathcal C_{\rm cont}^{\mathrm{sp}}
 \longrightarrow\mathcal G_{\mathrm{sp}},
 \qquad
 s_{\mathrm{sp}}(x)=(x,D_{\mathrm{sp}}x).
\end{equation}
La projection \(\pi_1(x,d)=x\) vérifie
\begin{equation}
\label{eq:continuo-sp-res-inversa}
 \pi_1s_{\mathrm{sp}}=I,
 \qquad
 s_{\mathrm{sp}}\pi_1=I_{\mathcal G_{\mathrm{sp}}}.
\end{equation}

Soit \(r_{\mathrm{sp},k}(x_k)\) la multisection des feuillets et des chemins qui appartient
déjà à \(D_{\mathrm{sp},k}(x_k)\). On définit
\begin{equation}
\label{eq:continuo-sp-x-dependiente}
 X_{\chi,\mathrm{sp},k}
 =\{(r_{\mathrm{sp},k}x_k,x_k,D_{\mathrm{sp},k}x_k):
     x_k\in\mathcal C_k^{\mathrm{sp}}\},
\end{equation}
\begin{equation}
\label{eq:continuo-sp-pr-x}
 \operatorname{pr}_{X,\mathrm{sp}}(x,D_{\mathrm{sp}}x)
 =(r_{\mathrm{sp}}x,x,D_{\mathrm{sp}}x).
\end{equation}
Oublier la première coordonnée de
\eqref{eq:continuo-sp-x-dependiente} est l'inverse de
\(\operatorname{pr}_{X,\mathrm{sp}}\).

\section{Faisceau de feuillets et mesure projective}

Soit
\begin{equation}
\label{eq:continuo-sp-haz-hojas}
 \widetilde{\mathcal C}_k^{\mathrm{sp}}
 =\{(x_k,h):x_k\in\mathcal C_k^{\mathrm{sp}},
              h\in\{\Sigma,\Pi\}\}.
\end{equation}
La mesure projective du continu induit des probabilités
\(\widetilde\mu_k\) sur ce faisceau. On se restreint au support positif ; pour
\(z\) dans ce support,
\begin{equation}
\label{eq:continuo-sp-consistencia-medida}
 \widetilde\mu_k(z)
 =\sum_{\widetilde\tau_{k+1,k}(z')=z}
  \widetilde\mu_{k+1}(z').
\end{equation}
La linéarisation de niveau est
\begin{equation}
\label{eq:continuo-sp-hilbert-nivel}
 \mathcal H_k^{\mathrm{sp}}
 =\ell^2(\operatorname{supp}\widetilde\mu_k),
\end{equation}
de base orthonormale \(e_z\).

\section{Transport de toutes les extensions et retenue}

Pour \(z'\) qui prolonge \(z\), nous définissons le poids conditionnel positif
\begin{equation}
\label{eq:continuo-sp-peso-condicional}
 w_k(z'\mid z)
 =\left(\frac{\widetilde\mu_{k+1}(z')}
               {\widetilde\mu_k(z)}\right)^{1/2}.
\end{equation}
Le transport est
\begin{equation}
\label{eq:continuo-sp-u-accion}
 U_ke_z
 =\sum_{\substack{z':\widetilde\tau_{k+1,k}(z')=z\\
            (z,z')\in\operatorname{Ext}_{\Gamma,\mathrm{sp},k}}}
 w_k(z'\mid z)e_{z'}.
\end{equation}
La somme contient toutes les extensions admissibles. Elle ne dépend pas d'un sélecteur.

Par \eqref{eq:continuo-sp-consistencia-medida},
\[
 \sum_{z':\widetilde\tau(z')=z}|w_k(z'\mid z)|^2=1.
\]
Les ensembles de descendants de deux parents distincts sont disjoints. Par
conséquent,
\begin{equation}
\label{eq:continuo-sp-u-isometria}
 U_k^*U_k=I_{\mathcal H_k^{\mathrm{sp}}}.
\end{equation}

Chaque base \(e_{z'}\) inclut la retenue de l'extension correspondante. Pour
des chemins composables, on conserve
\begin{align}
 \kappa(\gamma_2\gamma_1)
 &=\kappa(\gamma_2)+
   \mathsf C(\gamma_2)\kappa(\gamma_1),
 \label{eq:continuo-sp-carry-nativo}\\
 \operatorname{Carr}(\gamma_2\gamma_1)
 &=\operatorname{Carr}(\gamma_2)H(\gamma_1)
  +Y(\gamma_2)\operatorname{Carr}(\gamma_1).
 \label{eq:continuo-sp-carry-operador}
\end{align}
Les poids conditionnels se multiplient le long de l'unique chaîne de
troncatures d'un descendant. Les équations
\eqref{eq:continuo-sp-carry-nativo}--
\eqref{eq:continuo-sp-carry-operador} garantissent que les étiquettes de retenue
coïncident. Par conséquent,
\begin{equation}
\label{eq:continuo-sp-u-composicion}
 U_{\ell,k}=U_{\ell-1}\cdots U_k.
\end{equation}

\section{Graduation, projecteurs et décomposition du transport}

La graduation de feuillet est l'involution
\begin{equation}
\label{eq:continuo-sp-sigma}
 \sigma_ke_{(x,\Sigma)}=+e_{(x,\Sigma)},
 \qquad
 \sigma_ke_{(x,\Pi)}=-e_{(x,\Pi)}.
\end{equation}
Ainsi,
\[
 \sigma_k^*=\sigma_k,
 \qquad
 \sigma_k^2=I.
\]
Les projecteurs internes sont
\begin{equation}
\label{eq:continuo-sp-p-q}
 P_k=\frac{I+\sigma_k}{2},
 \qquad
 Q_k=\frac{I-\sigma_k}{2}.
\end{equation}
Elles satisfont
\begin{equation}
\label{eq:continuo-sp-p-q-relaciones}
 P_k^2=P_k,
 \quad Q_k^2=Q_k,
 \quad P_kQ_k=0,
 \quad P_k+Q_k=I.
\end{equation}

Nous séparons la partie du transport qui conserve le feuillet de celle qui le change :
\begin{align}
 A_k&=P_{k+1}U_kP_k+Q_{k+1}U_kQ_k,
 \label{eq:continuo-sp-a}\\
 \eta_k&=Q_{k+1}U_kP_k+P_{k+1}U_kQ_k.
 \label{eq:continuo-sp-eta}
\end{align}
Alors
\begin{equation}
\label{eq:continuo-sp-u-a-eta}
 U_k=A_k+\eta_k,
\end{equation}
et la graduation vérifie
\begin{equation}
\label{eq:continuo-sp-paridad-a-eta}
 \sigma_{k+1}A_k=A_k\sigma_k,
 \qquad
 \sigma_{k+1}\eta_k=-\eta_k\sigma_k.
\end{equation}

\begin{proposition}[Bilan local des feuillets]
\label{prop:continuo-sp-balance-local}
Pour chaque niveau,
\begin{equation}
\label{eq:continuo-sp-balance-local}
 A_k^*A_k+\eta_k^*\eta_k=I.
\end{equation}
\end{proposition}

\begin{proof}
L'orthogonalité de \(P_{k+1}\) et \(Q_{k+1}\) donne
\begin{align*}
 A_k^*A_k
 &=P_kU_k^*P_{k+1}U_kP_k
   +Q_kU_k^*Q_{k+1}U_kQ_k,\\
 \eta_k^*\eta_k
 &=P_kU_k^*Q_{k+1}U_kP_k
   +Q_kU_k^*P_{k+1}U_kQ_k.
\end{align*}
En sommant et en utilisant \(P_{k+1}+Q_{k+1}=I\),
\[
 A_k^*A_k+\eta_k^*\eta_k
 =P_kU_k^*U_kP_k+Q_kU_k^*U_kQ_k
 =P_k+Q_k=I.
\]
\end{proof}

L'opérateur \(\eta_k\) enregistre l'abandon du canal qui conserve le feuillet. Il
n'est pas la retenue entière : les coordonnées
\eqref{eq:continuo-sp-carry-nativo}--
\eqref{eq:continuo-sp-carry-operador} restent dans le paquet
\eqref{eq:continuo-sp-d13}.

\section{Lecteur signé avec ledger}

Pour \(v\in\mathcal H_k^{\mathrm{sp}}\), la polarisation signée est
\begin{equation}
\label{eq:continuo-sp-polarizacion-firmada}
 \operatorname{pol}_k(v)
 =\langle v,\sigma_kv\rangle
 =\|P_kv\|^2-\|Q_kv\|^2.
\end{equation}
Le lecteur spécialisé conserve en outre la source :
\begin{equation}
\label{eq:continuo-sp-lector-firmado}
 \mathcal L^{\rm sgn}_{\mathrm{sp},k}(x_k;v)
 =\bigl(
   \Lambda(\gamma_{x_k}),
   \Sigma_{\rm reg}(\Lambda(\gamma_{x_k})),
   \|P_kv\|^2,
   \|Q_kv\|^2,
   \operatorname{pol}_k(v)\bigr).
\end{equation}
L'orientation TRIT est conservée dans \(\Lambda(\gamma_{x_k})\) ; elle ne
s'identifie pas au signe de feuillet. À l'horizon \(N\), le registre garde la
famille des lecteurs pour tous les niveaux \(k\leq N\). Tronquer signifie
supprimer uniquement les entrées ultérieures.

\section{Produit non autonome et terminal}

On définit
\begin{equation}
\label{eq:continuo-sp-f-producto}
 F_{0,0}=I_{\mathcal H_0^{\mathrm{sp}}},
 \qquad
 F_{0,k}=A_{k-1}\cdots A_1A_0
 \quad(k\geq1).
\end{equation}
Par conséquent,
\begin{equation}
\label{eq:continuo-sp-f-recursion}
 F_{0,k+1}=A_kF_{0,k}.
\end{equation}
Le terme publié lors de l'abandon du canal à l'étape \(k\) et le terminal à
l'horizon \(N\) sont
\begin{equation}
\label{eq:continuo-sp-defecto-terminal}
 \mathcal D_k=F_{0,k}^*\eta_k^*\eta_kF_{0,k},
 \qquad
 \mathcal T_N=F_{0,N}^*F_{0,N}.
\end{equation}

\begin{theorem}[Télescopage spectral--polaire non autonome]
\label{thm:continuo-sp-telescopia-no-autonoma}
Pour tout horizon fini \(N\),
\begin{equation}
\label{eq:continuo-sp-telescopia-no-autonoma}
 \boxed{
 I=\sum_{k=0}^{N-1}
 F_{0,k}^*\eta_k^*\eta_kF_{0,k}
 +F_{0,N}^*F_{0,N}.}
\end{equation}
De plus,
\[
 0\preceq\mathcal T_{N+1}
 \preceq\mathcal T_N\preceq I,
\]
de sorte qu'il existe la limite forte positive
\[
 \mathcal T_\infty
 =\operatorname{s-lim}_{N\to\infty}\mathcal T_N
\]
et
\begin{equation}
\label{eq:continuo-sp-telescopia-infinita}
 I=\sum_{k=0}^{\infty}
 F_{0,k}^*\eta_k^*\eta_kF_{0,k}
 +\mathcal T_\infty
\end{equation}
en topologie forte. Le théorème n'affirme pas que
\(\mathcal T_\infty\) soit nul.
\end{theorem}

\begin{proof}
En multipliant \eqref{eq:continuo-sp-balance-local} par
\(F_{0,k}^*\) et \(F_{0,k}\), et en utilisant
\eqref{eq:continuo-sp-f-recursion}, on obtient
\[
 F_{0,k}^*F_{0,k}
 =F_{0,k+1}^*F_{0,k+1}
  +F_{0,k}^*\eta_k^*\eta_kF_{0,k}.
\]
La somme de \(k=0\) à \(N-1\) annule tous les termes
intermédiaires et laisse précisément
\eqref{eq:continuo-sp-telescopia-no-autonoma}. Comme
\(A_k^*A_k\preceq I\),
\[
 \mathcal T_{N+1}
 =F_{0,N}^*A_N^*A_NF_{0,N}
 \preceq F_{0,N}^*F_{0,N}=\mathcal T_N.
\]
La monotonie et la borne positive fournissent la limite forte. Le passage à la
limite dans l'identité finie donne
\eqref{eq:continuo-sp-telescopia-infinita}.
\end{proof}

\section{Assembleur comme graphe}

Pour
\[
 z_N=(r_{\mathrm{sp},N}x_N,x_N,D_{\mathrm{sp},N}x_N)
 \in X_{\chi,\mathrm{sp},N},
\]
on définit l'assembleur matériel
\begin{equation}
\label{eq:continuo-sp-a-accion}
\begin{aligned}
 a_{\mathrm{sp},N}(z_N)=\bigl(&
 (\mathcal H_k^{\mathrm{sp}},\sigma_k,P_k,Q_k)_{k\leq N},
 (U_k,A_k,\eta_k)_{k<N},\\
 &(F_{0,k},\mathcal D_k,\mathcal T_k,
   \mathcal L^{\rm sgn}_{\mathrm{sp},k})_{k\leq N}\bigr).
\end{aligned}
\end{equation}
La macrostructure est
\begin{equation}
\label{eq:continuo-sp-m-graph}
 \mathfrak M_{\mathrm{sp},N}
 =\operatorname{Graph}(a_{\mathrm{sp},N}),
 \qquad
 \mathcal A_{\mathrm{sp},N}(z_N)
 =(z_N,a_{\mathrm{sp},N}z_N).
\end{equation}
La projection sur \(z_N\) est l'inverse de
\(\mathcal A_{\mathrm{sp},N}\). La troncature supprime les entrées dont
l'indice dépasse le nouvel horizon et vérifie
\begin{equation}
\label{eq:continuo-sp-a-naturalidad}
 m^{\mathrm{sp}}_{\ell,k}a_{\mathrm{sp},\ell}
 =a_{\mathrm{sp},k}\xi^{\mathrm{sp}}_{\ell,k}.
\end{equation}

\section{Publicateur comme graphe}

Le publicateur ne contracte pas la macrostructure en un scalaire. Il renvoie le
certificat complet
\begin{equation}
\label{eq:continuo-sp-p-accion}
 p_{\mathrm{sp},N}(m_N)
 =\left(
 \mathcal L^{\rm sgn}_{\mathrm{sp},k},
 \mathcal D_k,
 \mathcal T_k,
 I=\sum_{j=0}^{k-1}\mathcal D_j+\mathcal T_k
 \right)_{0\leq k\leq N}.
\end{equation}
On définit
\begin{equation}
\label{eq:continuo-sp-c-graph}
 \mathcal C_{\mathrm{sp},N}
 =\operatorname{Graph}(p_{\mathrm{sp},N}),
 \qquad
 \mathcal P_{\mathrm{sp},N}(m_N)
 =(m_N,p_{\mathrm{sp},N}m_N).
\end{equation}
La projection sur \(m_N\) est son inverse. Tronquer le certificat supprime seulement
les identités ultérieures, de sorte que
\begin{equation}
\label{eq:continuo-sp-p-naturalidad}
 c^{\mathrm{sp}}_{\ell,k}p_{\mathrm{sp},\ell}
 =p_{\mathrm{sp},k}m^{\mathrm{sp}}_{\ell,k}.
\end{equation}

\section{Théorème de la chaîne interne complète}

\begin{theorem}[Macrostructure spectrale--polaire engendrée par le continu]
\label{thm:continuo-sp-cadena-interna}
Les actions
\eqref{eq:continuo-sp-res-graph},
\eqref{eq:continuo-sp-pr-x},
\eqref{eq:continuo-sp-m-graph} et
\eqref{eq:continuo-sp-c-graph} forment la chaîne naturelle
\begin{equation}
\label{eq:continuo-sp-cadena-interna}
 \mathcal C_{\rm cont}^{\rm disc}
 \supset\mathcal C_{\rm cont}^{\mathrm{sp}}
 \xrightarrow[\cong]{\operatorname{Res}_{\mathrm{sp}}}
 \mathcal G_{\mathrm{sp}}
 \xrightarrow[\cong]{\operatorname{pr}_{X,\mathrm{sp}}}
 X_{\chi,\mathrm{sp}}
 \xrightarrow[\cong]{\mathcal A_{\mathrm{sp}}}
 \mathfrak M_{\mathrm{sp}}
 \xrightarrow[\cong]{\mathcal P_{\mathrm{sp}}}
 \mathcal C_{\mathrm{sp}}.
\end{equation}
Son contenu publié est le faisceau à deux feuillets, la graduation
\(\sigma_k\), les projecteurs \(P_k,Q_k\), le transport total \(U_k\),
la séparation \(U_k=A_k+\eta_k\), la retenue et la mémoire, le lecteur signé
et le télescopage
\eqref{eq:continuo-sp-telescopia-no-autonoma} avec terminal explicite. Chaque
objet conserve une projection canonique vers tous ses antécédents.
\end{theorem}

\begin{proof}
Les identités
\eqref{eq:continuo-sp-res-inversa} prouvent la première équivalence. La
définition dépendante \eqref{eq:continuo-sp-x-dependiente} prouve la
deuxième. Les graphes
\eqref{eq:continuo-sp-m-graph} et
\eqref{eq:continuo-sp-c-graph} prouvent les deux restantes. La naturalité des
quatre étapes est
\eqref{eq:continuo-sp-d-naturalidad},
\eqref{eq:continuo-sp-a-naturalidad} et
\eqref{eq:continuo-sp-p-naturalidad}, avec la naturalité de
\(r_{\mathrm{sp},k}\). Le contenu opératoriel est prouvé par
\eqref{eq:continuo-sp-u-isometria}, le
\cref{prop:continuo-sp-balance-local} et le
\cref{thm:continuo-sp-telescopia-no-autonoma}. Toutes les extensions apparaissent
dans \eqref{eq:continuo-sp-u-accion} ; aucun choix de
chemin n'intervient donc.
\end{proof}

\fi
\section{Retour de phase et évolution non réinitialisée de l'état commun}

La branche conserve expressément
\begin{equation}
\label{eq:continuo-sp-no-reset}
 g(k+9)=g(k),
 \qquad
 x_{k+9}\neq x_k\quad\text{en général}.
\end{equation}
Le second membre enregistre l'accroissement de mémoire et les changements possibles
de feuillet, de retenue, de cylindre et de frontière. Ni \(\sigma_k\), ni les projecteurs, ni
le lecteur signé ne contractent cette différence.

\iffalse
\section{Provenance}

Les deux feuillets APP, l'état TRIT enrichi, les fibres et extensions TPK,
la retenue, le retour sans réinitialisation, la survie et le ledger sont une
\texttt{ARCHITECTURE\_AUTORALE\_PRÉEXISTANTE}. L'obligation de conserver le
terminal fini est un \texttt{RÉSULTAT\_RÉCUPÉRÉ}. Le faisceau explicite de feuillets,
l'involution \(\sigma_k\), les projecteurs, la linéarisation qui somme toutes
les extensions, la décomposition \(U_k=A_k+\eta_k\), le produit non
autonome \(F_{0,k}\) et les graphes dépendants sont une
\texttt{FORMALISATION\_NOUVELLE}. Le
\cref{thm:continuo-sp-telescopia-no-autonoma,thm:continuo-sp-cadena-interna}
constitue le \texttt{CERTIFICAT\_NOUVEAU} de cette formalisation.

\section{Falsadores internos}

\begin{proposition}[Falsificateurs de la macrostructure]
\label{prop:continuo-sp-falsadores}
La conclusion du
\cref{thm:continuo-sp-cadena-interna} n'est pas transférable si l'une quelconque
des substitutions suivantes se produit :
\begin{enumerate}
 \item choisir un prolongement dans
       \eqref{eq:continuo-sp-u-accion} et effacer les autres ;
 \item remplacer \(\mathcal G_{\mathrm{sp}}\) par l'image extérieure de
       \(D_{\mathrm{sp}}\) ;
 \item effacer \((x,D_{\mathrm{sp}}x)\) lors de la formation de
       \(X_{\chi,\mathrm{sp}}\) ;
 \item identifier le feuillet à l'orientation TRIT ou effacer le seuil ;
 \item rompre la cohérence
       \eqref{eq:continuo-sp-consistencia-medida} ;
 \item effacer la retenue ou enfreindre
       \eqref{eq:continuo-sp-carry-nativo}--
       \eqref{eq:continuo-sp-carry-operador} ;
 \item identifier \(\eta_k\) à toute la retenue ;
 \item remplacer la famille \((A_k)_k\) par un seul opérateur sans démontrer
       la stationnarité ;
 \item supprimer \(F_{0,N}^*F_{0,N}\) à un horizon fini ;
 \item affirmer \(\mathcal T_\infty=0\) sans preuve supplémentaire ;
 \item reconstruire une histoire à partir de la polarisation signée ;
 \item utiliser \(g(k+9)=g(k)\) pour affirmer \(x_{k+9}=x_k\) ;
 \item introduire un domaine, une fonction ou un critère extérieur pour
       définir l'une quelconque des flèches internes ;
 \item introduire une sortie ultérieure comme antécédent de la branche.
\end{enumerate}
\end{proposition}
\fi

\endgroup

L'espérance conditionnelle associée au même poids est
\begin{equation}
\label{eq:pdfv2-cor-esperanza-comun}
 (E_kF)(\widehat x):=
 \sum_{(\alpha,\widehat y)\in\operatorname{Ch}_k(\widehat x)}
 p_k(\alpha,\widehat y\mid\widehat x)F(\alpha,\widehat y).
\end{equation}
Pour l'accroissement fin \(b_f\), on pose
\begin{equation}
\label{eq:pdfv2-cor-kappa-jensen}
 b_c:=E_kb_f,\qquad
 \kappa_k^{\rm can}:=\frac{E_k|b_f|-|b_c|}{2}\geq0.
\end{equation}
Alors
\begin{equation}
\label{eq:pdfv2-cor-jensen-hojas}
 E_kb_f^+=b_c^++\kappa_k^{\rm can},\qquad
 E_kb_f^-=b_c^-+\kappa_k^{\rm can}.
\end{equation}
L'exposant \({\rm can}\) distingue cette retenue d'annulation de la retenue entière \(\kappa_\alpha\) ; toutes deux proviennent du même transport, mais sont de types distincts.

Les mémoires signées ne nécessitent pas de condition initiale extérieure : pour \(\gamma=(\epsilon_1,\ldots,\epsilon_k)\), elles sont définies par la somme du registre lui-même,
\begin{equation}
\label{eq:pdfv2-cor-memorias}
 \begin{aligned}
 M_0^\pm&:=0,\\
 M_k^\pm(\gamma)&:=\sum_{j=1}^k b_{\epsilon_j}^\pm,\\
 M_{k+1}^\pm(\gamma\alpha)&:=M_k^\pm(\gamma)+b_\alpha^\pm,\\
 D_k^0&:=M_k^+-M_k^-, &
 H_k^0&:=M_k^++M_k^-.
 \end{aligned}
\end{equation}
Ainsi, le bilan signé, la mémoire totale et la polarisation \(\langle v,\sigma_kv\rangle\) sont trois lectures coordonnées du même préfixe, non trois états reconstruits séparément.

\subsection{Raffinement nonadique, pertes orthogonales et mémoire terminale}
L'indice \(j\) du développement suivant est le niveau du même arbre des prolongements, et son registre à treize coordonnées est la restriction du registre commun de \eqref{eq:pdfv2-cor-registro-extension}. On fixe donc \(D_{{\rm sol},j}:=D_j^{(13)}\) sans postuler de domaine supplémentaire.
\begingroup
\let\section\subsection
% Extracción quirúrgica del propietario V2 05b: líneas 100--359.
% El resto permanece en este archivo como procedencia, pero no se publica.
\iffalse
% Macroestructura solenoidal estrictamente interna del continuo discreto.
% Este módulo no introduce ninguna ecuación clásica ni un problema exterior.

\chapter{Macrostructure interne de bilan, mémoire et télescopage}
\label{chap:pdfv2-sol-interna}

\begin{statusbox}
L'unique entrée de ce chapitre est la structure discrète du continu.
La branche \(\mathrm{sol}\) se construit avant toute évaluation scalaire
extérieure et avant toute formulation différentielle classique. Sa chaîne
complète est
\[
 \mathcal C_{\rm cont}^{\rm disc}
 \supset\mathcal C_{\rm cont}^{\rm sol}
 \xrightarrow{\operatorname{Res}_{\rm sol}}
 \mathcal G_{\rm sol}
 \xrightarrow{\operatorname{pr}_{X,{\rm sol}}}
 X_{\chi,{\rm sol}}
 \xrightarrow{\mathcal A_{\rm sol}}
 \mathfrak M_{\rm sol}
 \xrightarrow{\mathcal P_{\rm sol}}
 \mathcal C_{\rm sol}^{\rm HMT}.
\]
Chaque flèche conserve comme coordonnées l'histoire et l'appareil qui la
produit. Aucune image de cette chaîne n'est un quotient qui oublie sa source.
\end{statusbox}

\section{Domaine propre avant toute évaluation ultérieure}
\label{sec:pdfv2-sol-dominio}

Soit
\[
 x=(x_0\xrightarrow{e_0}x_1\xrightarrow{e_1}\cdots)
 \in\mathcal C_{\rm cont}^{\rm disc}
\]
une histoire survivante. Chaque arête \(e_j\) conserve ses deux feuillets,
son orientation, son résidu, son quotient et sa retenue entière. Soient
\(N_j^+(e_j)\) et \(N_j^-(e_j)\) les multiplicités des deux feuillets et soit
\(\Delta_{\gamma_j}\in\mathbb Z\) le coefficient orienté du chemin, avec la
retenue de la composition déjà incorporée. L'accroissement signé interne est
\begin{equation}\label{eq:pdfv2-sol-b}
 b_j=b(e_j):=
 \Delta_{\gamma_j}\bigl(N_j^+(e_j)-N_j^-(e_j)\bigr)\in\mathbb Z.
\end{equation}
Il ne s'obtient pas en évaluant une trajectoire extérieure : c'est une coordonnée du
ledger de l'arête elle-même.

Ses parties positive et négative se construisent par
\begin{equation}\label{eq:pdfv2-sol-bpm}
 b_j^+:=\frac{|b_j|+b_j}{2},
 \qquad
 b_j^-:=\frac{|b_j|-b_j}{2}.
\end{equation}
Par conséquent,
\begin{equation}\label{eq:pdfv2-sol-bpm-identidades}
 b_j=b_j^+-b_j^-,
 \qquad
 |b_j|=b_j^++b_j^-,
 \qquad
 b_j^+b_j^-=0.
\end{equation}

Les deux mémoires primitives et leurs deux lecteurs sont
\begin{align}
 M_0^+=M_0^-&=0,
 &M_j^+&=\sum_{0\leq i<j}b_i^+,
 &M_j^-&=\sum_{0\leq i<j}b_i^-;
 \label{eq:pdfv2-sol-memorias}\\
 D_j^0&:=M_j^+-M_j^-,
 &H_j^0&:=M_j^++M_j^-.
 \label{eq:pdfv2-sol-dh}
\end{align}
En conséquence,
\begin{equation}\label{eq:pdfv2-sol-dh-diferencias}
 D_{j+1}^0-D_j^0=b_j,
 \qquad
 H_{j+1}^0-H_j^0=|b_j|.
\end{equation}
La mémoire signée \(D^0\) et la mémoire totale \(H^0\) restent des
objets distincts ; conserver seulement leur différence finale détruirait la
généalogie.

La sous-structure \(\mathcal C_{\rm cont}^{\rm sol}\) se compose des histoires
pour lesquelles, à tout préfixe et à tout raffinement nonadique :
\begin{enumerate}
 \item les accroissements \eqref{eq:pdfv2-sol-b} et les deux mémoires
       \eqref{eq:pdfv2-sol-memorias} existent ;
 \item les cylindres fins se regroupent en fibres de neuf extensions avec la
       même loi de probabilité projective ;
 \item les projections de profondeur et de résolution définies ci-dessous sont
       compatibles avec les extensions d'histoire ;
 \item chaque ligne de pertes et chaque colonne terminale est finie à horizon
       fini ;
 \item tout préfixe admet au moins une extension compatible à tout
       horizon ultérieur.
\end{enumerate}
Ce sont des conditions d'appartenance à la branche, non des prémisses tacites
attribuées à une histoire arbitraire.

\fi
\section{Raffinement nonadique commun : inclusion, espérance et mémoire de bilan}
\label{sec:pdfv2-sol-e9j9}

Soit \(Y_j(x)\) l'ensemble fini des cylindres compatibles avec le préfixe
\(x_{\leq j}\), muni de la mesure projective \(\mu_j\). La troncature
\[
 \tau_{j+1,j}:Y_{j+1}(x)\longrightarrow Y_j(x)
\]
a neuf extensions admissibles au-dessus de chaque cylindre. On considère les
espaces finis
\[
 \mathscr H_j(x):=\ell^2\bigl(Y_j(x),\mu_j\bigr).
\]
L'inclusion et l'espérance nonadiques sont
\begin{align}
 J_{9,j}:\mathscr H_j&\longrightarrow\mathscr H_{j+1},
 &(J_{9,j}f)(y)&=f(\tau_{j+1,j}y),
 \label{eq:pdfv2-sol-j9}\\
 E_{9,j}:\mathscr H_{j+1}&\longrightarrow\mathscr H_j,
 &(E_{9,j}g)(z)&=\frac1{9}
   \sum_{\tau_{j+1,j}y=z}g(y).
 \label{eq:pdfv2-sol-e9}
\end{align}
La compatibilité projective de \(\mu_{j+1}\) et \(\mu_j\) donne
\begin{equation}\label{eq:pdfv2-sol-ej-projector}
 E_{9,j}J_{9,j}=I_{\mathscr H_j},
 \qquad
 J_{9,j}^*=E_{9,j},
 \qquad
 P_{9,j}:=J_{9,j}E_{9,j}=P_{9,j}^*=P_{9,j}^2.
\end{equation}
Le projecteur microscopique complémentaire est
\begin{equation}\label{eq:pdfv2-sol-qmicro-preliminar}
 Q_{9,j}:=I-P_{9,j}.
\end{equation}

Si \(b_f\) est l'accroissement sur la fibre fine et
\begin{equation}\label{eq:pdfv2-sol-bcoarse}
 b_c:=E_{9,j}b_f,
\end{equation}
la convexité élémentaire de \(|\cdot|\) construit, et ne postule pas, la retenue de
compensation
\begin{equation}\label{eq:pdfv2-sol-kappa}
 \kappa_j
 :=\frac{E_{9,j}|b_f|-|b_c|}{2}\geq0.
\end{equation}
En effet, en utilisant \(t^\pm=(|t|\pm t)/2\), on obtient
\begin{align}
 E_{9,j}b_f^+
 &=\frac{E_{9,j}|b_f|+E_{9,j}b_f}{2}
   =b_c^++\kappa_j,
 \label{eq:pdfv2-sol-jensen-plus}\\
 E_{9,j}b_f^-
 &=\frac{E_{9,j}|b_f|-E_{9,j}b_f}{2}
   =b_c^-+\kappa_j.
 \label{eq:pdfv2-sol-jensen-minus}
\end{align}
Les deux feuillets perdent la même quantité lors de la compression ; c'est pourquoi la différence
signée est conservée tandis que la mémoire totale enregistre la compensation :
\begin{equation}\label{eq:pdfv2-sol-jensen-dh}
 E_{9,j}(b_f^+-b_f^-)=b_c,
 \qquad
 E_{9,j}(b_f^++b_f^-)=|b_c|+2\kappa_j.
\end{equation}

\section{Factorisation orthogonale des degrés de profondeur et des couches de résolution dans \(\mathscr H_j^{\mathrm{depth}}\widehat\otimes\mathscr H_j^{\mathrm{res}}\)}
\label{sec:pdfv2-sol-micro-shell}

Pour empêcher qu'une même perte soit comptée à la fois comme profondeur et
comme résolution, l'espace de chaque niveau est typé comme
\begin{equation}\label{eq:pdfv2-sol-factor-hilbert}
 \mathscr H_j
 =\mathscr H_j^{\rm depth}\widehat\otimes
  \mathscr H_j^{\rm res}.
\end{equation}
Pour \(j\geq1\), le premier facteur porte
\[
 P_j^{\rm depth}:=J_{9,j-1}E_{9,j-1}
 \quad\text{sur }\mathscr H_j^{\rm depth},
\]
et l'on fixe \(P_0^{\rm depth}=I\). Le second facteur porte la partition
finie des shells héritée de la frontière des cylindres. Si
\(\sigma_j:Y_j^{\rm res}\to S_j\) est cette partition, ses opérateurs sont
\begin{align}
 J_j^{\rm res}f&=f\circ\sigma_j,
 &E_j^{\rm res}g(s)&=
 \frac1{\#\sigma_j^{-1}(s)}
 \sum_{\sigma_j(y)=s}g(y),
 \label{eq:pdfv2-sol-shell-je}\\
 P_j^{\rm res}&:=J_j^{\rm res}E_j^{\rm res}.
 \label{eq:pdfv2-sol-shell-p}
\end{align}

On définit les trois projecteurs mutuellement orthogonaux
\begin{align}
 P_j^{\rm term}
 &:=P_j^{\rm depth}\otimes P_j^{\rm res},
 \label{eq:pdfv2-sol-pterm}\\
 Q_j^{\rm micro}
 &:=(I-P_j^{\rm depth})\otimes I,
 \label{eq:pdfv2-sol-qmicro}\\
 Q_j^{\rm shell}
 &:=P_j^{\rm depth}\otimes(I-P_j^{\rm res}).
 \label{eq:pdfv2-sol-qshell}
\end{align}
Alors
\begin{equation}\label{eq:pdfv2-sol-ortogonalidad}
 Q_j^{\rm micro}Q_j^{\rm shell}=0,
 \qquad
 I=P_j^{\rm term}+Q_j^{\rm micro}+Q_j^{\rm shell},
\end{equation}
et, pour tout \(\xi\in\mathscr H_j\),
\begin{equation}\label{eq:pdfv2-sol-pitagoras}
 \|\xi\|^2
 =\|P_j^{\rm term}\xi\|^2
  +\|Q_j^{\rm micro}\xi\|^2
  +\|Q_j^{\rm shell}\xi\|^2.
\end{equation}
Le terme
\((I-P_j^{\rm depth})\otimes(I-P_j^{\rm res})\) n'apparaît pas une seconde fois : il est déjà
contenu exactement une fois dans \(Q_j^{\rm micro}\).

\section{Extensions de cylindre et ligne complète de pertes}
\label{sec:pdfv2-sol-gamma-loss}

Une extension de cylindres du niveau \(j\) au niveau \(j+1\) induit l'application
\begin{equation}\label{eq:pdfv2-sol-gamma}
 \Gamma_j:\mathscr H_j\longrightarrow\mathscr H_{j+1},
 \qquad
 \Gamma_j f=f\circ\tau_{j+1,j},
\end{equation}
avec toutes les étiquettes de feuillet, d'orientation, de retenue, de mémoire et de chemin
transportées dans la base des cylindres. Pour \(n>j\), on écrit
\begin{equation}\label{eq:pdfv2-sol-gamma-composition}
 \Gamma_{j:n}:=\Gamma_{n-1}\cdots\Gamma_j,
 \qquad
 \Gamma_{j:j}:=I.
\end{equation}

Sur \(\mathscr H_j\), les accroissements non négatifs définissent des opérateurs de
multiplication
\begin{equation}\label{eq:pdfv2-sol-loss-multipliers}
 B_j^+:=M_{\sqrt{b_j^+}},
 \qquad
 B_j^-:=M_{\sqrt{b_j^-}},
 \qquad
 K_j:=M_{\sqrt{2\kappa_j}}.
\end{equation}
Chaque racine est prise point par point sur l'ensemble fini des cylindres ; si une
coordonnée est nulle, son opérateur est nul. L'espace des pertes est la somme
orthogonale externe
\[
 \mathscr E_j
 :=\mathscr H_j^{(+)}\oplus\mathscr H_j^{(-)}
   \oplus\mathscr H_j^{(\kappa)}
   \oplus\mathscr H_j^{({\rm micro})}
   \oplus\mathscr H_j^{({\rm shell})},
\]
et la ligne complète est
\begin{equation}\label{eq:pdfv2-sol-loss-row}
 L_j:\mathscr H_j\longrightarrow\mathscr E_j,
 \qquad
 L_j\xi
 =\bigl(
 B_j^+\xi,B_j^-\xi,K_j\xi,
 Q_j^{\rm micro}\xi,Q_j^{\rm shell}\xi
 \bigr).
\end{equation}
Par conséquent,
\begin{equation}\label{eq:pdfv2-sol-loss-norm}
 \|L_j\xi\|^2
 =\|B_j^+\xi\|^2+\|B_j^-\xi\|^2+\|K_j\xi\|^2
  +\|Q_j^{\rm micro}\xi\|^2
  +\|Q_j^{\rm shell}\xi\|^2.
\end{equation}
Les parties positive, négative, retenue de compensation, microstructure et shell
apparaissent chacune une fois. Elles ne fusionnent pas en une constante anonyme et ne se
répètent pas dans deux termes de la somme.

\section{Terminal d'histoire complète}
\label{sec:pdfv2-sol-terminal}

Soit \(\operatorname{Hist}_{J}^{\rm sol}\) l'ensemble des histoires
survivantes jusqu'à l'horizon \(J\), chacune avec son paquet de treize
coordonnées. L'espace terminal est défini comme l'espace libre de ces
registres complets :
\begin{equation}\label{eq:pdfv2-sol-terminal-space}
 \mathscr T_J^{\rm sol}
 :=\ell^2\!\left(
 \left\{(h,D_{{\rm sol},J}(h)):
 h\in\operatorname{Hist}_{J}^{\rm sol}\right\}
 \right).
\end{equation}
L'évaluation terminale
\begin{equation}\label{eq:pdfv2-sol-terminal-evaluation}
 E_J^{\rm term}:\mathscr H_J\longrightarrow\mathscr T_J^{\rm sol}
\end{equation}
envoie chaque vecteur de base de cylindre au registre complet de l'histoire qui
l'engendre et se prolonge linéairement. En particulier, le terminal retient
fibre, relations, nombres, orientation, retenue, mémoire, frontière, chemin,
multisection, survie et lecteur. Il n'est pas remplacé par la dernière valeur
visible du chemin.

\section{Colonne d'observation, Gram et télescopage}
\label{sec:pdfv2-sol-gram}

Pour \(0\leq j\leq J\), la colonne future complète est
\begin{equation}\label{eq:pdfv2-sol-g-column}
 \begin{aligned}
 G_{j;J}:\mathscr H_j&\longrightarrow
 \mathscr T_J^{\rm sol}\oplus
 \bigoplus_{n=j}^{J-1}\mathscr E_n,\\
 G_{j;J}\xi
 &:={}
 E_J^{\rm term}\Gamma_{j:J}\xi
 \ \oplus\ 
 \bigoplus_{n=j}^{J-1}L_n\Gamma_{j:n}\xi.
 \end{aligned}
\end{equation}
Pour \(j=J\), la somme des pertes est vide et
\(G_{J;J}=E_J^{\rm term}\). Le Gram terminal est
\begin{equation}\label{eq:pdfv2-sol-u-gram}
 U_{j;J}:=G_{j;J}^*G_{j;J}\succeq0.
\end{equation}

En réordonnant la somme orthogonale du codomaine, on obtient l'identité littérale
\begin{equation}\label{eq:pdfv2-sol-g-recursion}
 G_{j;J}\xi
 \simeq
 \bigl(L_j\xi,\,G_{j+1;J}\Gamma_j\xi\bigr).
\end{equation}
En conséquence,
\begin{equation}\label{eq:pdfv2-sol-stein}
 \boxed{U_{j;J}
 =L_j^*L_j+\Gamma_j^*U_{j+1;J}\Gamma_j}
\end{equation}
et
\begin{equation}\label{eq:pdfv2-sol-stein-difference}
 U_{j;J}-\Gamma_j^*U_{j+1;J}\Gamma_j
 =L_j^*L_j\succeq0.
\end{equation}
Il n'y a pas de mutation de signe : le terme de perte est du côté positif
de l'identité parce qu'il provient d'une somme de carrés.

En itérant \eqref{eq:pdfv2-sol-stein}, on obtient le télescopage exact
\begin{equation}\label{eq:pdfv2-sol-telescopia-operatorial}
 U_{j;J}
 =\Gamma_{j:J}^*(E_J^{\rm term})^*E_J^{\rm term}\Gamma_{j:J}
  +\sum_{n=j}^{J-1}
   \Gamma_{j:n}^*L_n^*L_n\Gamma_{j:n},
\end{equation}
ou, sous forme quadratique,
\begin{equation}\label{eq:pdfv2-sol-telescopia-normas}
 \langle U_{j;J}\xi,\xi\rangle
 =\|E_J^{\rm term}\Gamma_{j:J}\xi\|^2
  +\sum_{n=j}^{J-1}\|L_n\Gamma_{j:n}\xi\|^2.
\end{equation}
Le premier terme conserve le futur terminal complet ; la somme conserve,
niveau par niveau, toutes les pertes et toutes les retenues qui y ont conduit.

\iffalse
\section{Les treize coordonnées de la restriction}
\label{sec:pdfv2-sol-trece}

Pour un préfixe \(x_{\leq j}\), on définit
\begin{equation}\label{eq:pdfv2-sol-d-package}
 \begin{aligned}
 D_{{\rm sol},j}(x):=\bigl(&
 \mathcal F_{{\rm sol},j},
 \operatorname{Ext}_{\Gamma,{\rm sol},j},
 \operatorname{Rel}_{{\rm sol},j},
 \mathcal N_{{\rm sol},j},
 \operatorname{or}_{{\rm sol},j},
 \operatorname{Carr}_{{\rm sol},j},
 \operatorname{Mem}_{{\rm sol},j},\\
 &\partial_{{\rm sol},j},
 \gamma_{{\rm sol},j},
 \operatorname{Msect}_{{\rm sol},j},
 \operatorname{Surv}_{{\rm sol},j},
 \operatorname{Term}_{{\rm sol},j},
 \mathcal L_{{\rm sol},j}\bigr).
 \end{aligned}
\end{equation}
Ses composantes sont les suivantes.
\begin{enumerate}
 \item \(\mathcal F_{{\rm sol},j}\) contient les cylindres \(Y_j\), leur mesure
       projective, \(\mathscr H_j^{\rm depth}\),
       \(\mathscr H_j^{\rm res}\), leurs produits tensoriels et les cinq
       espaces de pertes.
 \item \(\operatorname{Ext}_{\Gamma,{\rm sol},j}\) contient
       \(J_{9,j},E_{9,j},\Gamma_j,\Gamma_{j:n}\), les inclusions de shell et
       toutes les applications de troncature et de raffinement.
 \item \(\operatorname{Rel}_{{\rm sol},j}\) contient
       \eqref{eq:pdfv2-sol-bpm-identidades},
       \eqref{eq:pdfv2-sol-ej-projector},
       \eqref{eq:pdfv2-sol-jensen-plus}--
       \eqref{eq:pdfv2-sol-jensen-minus},
       \eqref{eq:pdfv2-sol-ortogonalidad} et
       \eqref{eq:pdfv2-sol-stein}.
 \item \(\mathcal N_{{\rm sol},j}\) conserve
       \(j,J,9,b_j,b_j^+,b_j^-,\kappa_j\), les multiplicités de feuillet, les tailles
       des fibres et les dimensions des shells.
 \item \(\operatorname{or}_{{\rm sol},j}\) conserve
       \(\Delta_{\gamma_j}\), l'ordre du chemin et l'action d'inversion
       \(b_j\mapsto-b_j\), qui échange \(b_j^+\) et \(b_j^-\) sans changer
       \(|b_j|\) ni \(\kappa_j\).
 \item \(\operatorname{Carr}_{{\rm sol},j}\) conserve la retenue TPK entière,
       la retenue de composition et la retenue de compensation \(\kappa_j\) comme
       coordonnées différentes.
 \item \(\operatorname{Mem}_{{\rm sol},j}\) est
       \((M_j^+,M_j^-,D_j^0,H_j^0)\) avec le ledger ordonné de tous les
       accroissements antérieurs.
 \item \(\partial_{{\rm sol},j}\) contient la frontière orientée des cylindres,
       les accroissements \(b_j\), la ligne \(L_j\) et les différences de Stein.
 \item \(\gamma_{{\rm sol},j}\) est le chemin complet
       \(e_0e_1\cdots e_{j-1}\), et non seulement ses extrémités.
 \item \(\operatorname{Msect}_{{\rm sol},j}\) contient toutes les extensions
       nonadiques et de shell compatibles ; il ne choisit pas un prolongement.
 \item \(\operatorname{Surv}_{{\rm sol},j}\) enregistre les horizons
       \(J\geq j\) pour lesquels il existe \(G_{j;J}\), \(U_{j;J}\) et
       des terminaux compatibles.
 \item \(\operatorname{Term}_{{\rm sol},j}\) contient
       \(E_J^{\rm term}\Gamma_{j:J}\), le registre complet d'histoire et
       les treize coordonnées à chaque horizon.
 \item \(\mathcal L_{{\rm sol},j}\) publie en interne
       \((b^\pm,M^\pm,D^0,H^0,\kappa,Q^{\rm micro},Q^{\rm shell},
       L,G,U)\) et ses identités.
\end{enumerate}

Pour une extension \(u:x_j\to x_{j'}\), le transport coordonné
\(u_{{\rm sol},j',j}\) satisfait
\begin{equation}\label{eq:pdfv2-sol-naturality-d}
 u_{{\rm sol},j',j}D_{{\rm sol},j'}(x_{j'})
 =D_{{\rm sol},j}(\tau_{j',j}x_{j'}).
\end{equation}
L'égalité inclut les cinq composantes de \(L\), le terminal complet et
la composition de \(\Gamma\) ; ce n'est pas une égalité des seuls cardinaux.

\section{Graphe de restriction et objet dépendant}
\label{sec:pdfv2-sol-graphs}

On définit
\begin{equation}\label{eq:pdfv2-sol-g-graph}
 \mathcal G_{\rm sol}:=\operatorname{Graph}(D_{\rm sol}),
 \qquad
 \operatorname{Res}_{\rm sol}(x):=(x,D_{\rm sol}x).
\end{equation}
La première projection est inverse sur le graphe :
\begin{equation}\label{eq:pdfv2-sol-res-inverse}
 \pi_1\operatorname{Res}_{\rm sol}=I,
 \qquad
 \operatorname{Res}_{\rm sol}\pi_1=I_{\mathcal G_{\rm sol}}.
\end{equation}

Soit \(\chi_{\rm sol}(x)\) la multisection complète formée de tous les
cylindres, raffinements, shells, lignes de pertes et terminaux compatibles
avec \(x\). Alors
\begin{equation}\label{eq:pdfv2-sol-xchi}
 X_{\chi,{\rm sol}}
 :=\{(\chi_{\rm sol}(x),x,D_{\rm sol}x):
 x\in\mathcal C_{\rm cont}^{\rm sol}\},
\end{equation}
et
\begin{equation}\label{eq:pdfv2-sol-prx}
 \operatorname{pr}_{X,{\rm sol}}(x,D_{\rm sol}x)
 :=(\chi_{\rm sol}(x),x,D_{\rm sol}x).
\end{equation}
De nouveau, la projection qui conserve \((x,D_{\rm sol}x)\) est inverse sur
cette image dépendante.

\section{Assembleur et publicateur internes}
\label{sec:pdfv2-sol-assembler-publisher}

L'espace ambiant \(\mathscr M_{\rm sol}^{\rm amb}\) se compose de systèmes
compatibles
\[
 \bigl(
 \{b,b^\pm,M^\pm,D^0,H^0,\kappa\},
 \{E_9,J_9,P^{\rm term},Q^{\rm micro},Q^{\rm shell}\},
 \{\Gamma,L,E^{\rm term},G,U\}
 \bigr)
\]
avec toutes les identités précédentes. L'assembleur brut est
\begin{equation}\label{eq:pdfv2-sol-assembler}
 \begin{aligned}
 a_{\rm sol}:X_{\chi,{\rm sol}}&\longrightarrow
 \mathscr M_{\rm sol}^{\rm amb},\\
 a_{\rm sol}(\chi(x),x,Dx)&:=
 \bigl(
 \{b,b^\pm,M^\pm,D^0,H^0,\kappa\}_x,
 \{E_9,J_9,P,Q\}_x,
 \{\Gamma,L,E^{\rm term},G,U\}_x
 \bigr).
 \end{aligned}
\end{equation}
La macrostructure conserve son antécédent par
\begin{equation}\label{eq:pdfv2-sol-macro-graph}
 \mathfrak M_{\rm sol}:=\operatorname{Graph}(a_{\rm sol}),
 \qquad
 \mathcal A_{\rm sol}(z):=(z,a_{\rm sol}z).
\end{equation}

Soit \(\mathscr Y_{\rm sol}^{\rm int}\) le type des certificats formés par
les identités
\begin{equation}\label{eq:pdfv2-sol-certificate-list}
 \begin{gathered}
 b=b^+-b^-,\qquad |b|=b^++b^-,\\
 E_9J_9=I,\qquad J_9E_9=P_9,\\
 E_9b_f^\pm=b_c^\pm+\kappa,\qquad\kappa\geq0,\\
 I=P^{\rm term}+Q^{\rm micro}+Q^{\rm shell},
 \qquad Q^{\rm micro}Q^{\rm shell}=0,\\
 U_{j;J}=L_j^*L_j+\Gamma_j^*U_{j+1;J}\Gamma_j,\\
 \langle U_{j;J}\xi,\xi\rangle
 =\|E_J^{\rm term}\Gamma_{j:J}\xi\|^2
  +\sum_{n=j}^{J-1}\|L_n\Gamma_{j:n}\xi\|^2,
 \end{gathered}
\end{equation}
avec la naturalité de toutes les flèches. Le publicateur brut
\begin{equation}\label{eq:pdfv2-sol-publisher}
 p_{\rm sol}^{\rm raw}:\mathfrak M_{\rm sol}
 \longrightarrow\mathscr Y_{\rm sol}^{\rm int}
\end{equation}
publie ce certificat sans supprimer le macro-objet. Enfin,
\begin{equation}\label{eq:pdfv2-sol-output-graph}
 \mathcal C_{\rm sol}^{\rm HMT}
 :=\operatorname{Graph}(p_{\rm sol}^{\rm raw}),
 \qquad
 \mathcal P_{\rm sol}(m):=(m,p_{\rm sol}^{\rm raw}m).
\end{equation}

\section{Théorème interne de génération solénoïdale}
\label{sec:pdfv2-sol-theorem}

\begin{theorem}[Génération interne de la macrostructure de bilan]
\label{thm:pdfv2-sol-internal-generation}
Pour tout \(x\in\mathcal C_{\rm cont}^{\rm sol}\), la chaîne
\[
 \mathcal C_{\rm cont}^{\rm sol}
 \xrightarrow{\operatorname{Res}_{\rm sol}}
 \mathcal G_{\rm sol}
 \xrightarrow{\operatorname{pr}_{X,{\rm sol}}}
 X_{\chi,{\rm sol}}
 \xrightarrow{\mathcal A_{\rm sol}}
 \mathfrak M_{\rm sol}
 \xrightarrow{\mathcal P_{\rm sol}}
 \mathcal C_{\rm sol}^{\rm HMT}
\]
construit naturellement une macrostructure qui contient les deux mémoires,
la retenue de compensation, la séparation orthogonale micro/shell, la ligne complète
de pertes, la colonne future, le terminal d'histoire complète, le Gram et le
télescopage. Les treize coordonnées du continu peuvent être récupérées par
projections successives depuis n'importe quel point de la chaîne.
\end{theorem}

\begin{proof}
Les identités \eqref{eq:pdfv2-sol-bpm-identidades} et
\eqref{eq:pdfv2-sol-dh-diferencias} s'obtiennent directement à partir de la
décomposition positive et négative de l'accroissement entier ; les
mémoires ne sont donc pas des données ajoutées. Les formules
\eqref{eq:pdfv2-sol-j9}--\eqref{eq:pdfv2-sol-e9}, avec la mesure
projective, prouvent \(E_9J_9=I\), \(J_9^*=E_9\) et que \(P_9\) est un
projecteur orthogonal.

L'inégalité triangulaire donne
\(E_9|b_f|\geq|E_9b_f|=|b_c|\), de sorte que
\(\kappa\geq0\). Substituer
\(b^\pm=(|b|\pm b)/2\) prouve
\eqref{eq:pdfv2-sol-jensen-plus} et
\eqref{eq:pdfv2-sol-jensen-minus} sans hypothèse supplémentaire.

Les définitions tensorielles
\eqref{eq:pdfv2-sol-pterm}--\eqref{eq:pdfv2-sol-qshell} prouvent que
\(P^{\rm term}\), \(Q^{\rm micro}\) et \(Q^{\rm shell}\) sont orthogonaux et
ont pour somme l'identité. Par conséquent, la ligne \(L_j\) compte chaque secteur une
seule fois.

La définition \eqref{eq:pdfv2-sol-g-column} sépare orthogonalement le terminal
et les pertes de chaque niveau. Séparer le premier terme de pertes produit
\eqref{eq:pdfv2-sol-g-recursion}. Prendre l'adjoint et composer prouve l'identité
de Stein \eqref{eq:pdfv2-sol-stein} ; l'itérer prouve
\eqref{eq:pdfv2-sol-telescopia-operatorial} et
\eqref{eq:pdfv2-sol-telescopia-normas}. Tout terme est une norme au carré,
de sorte que le signe opposé ne peut apparaître.

La compatibilité de la troncature des cylindres, de l'espérance, des projections et des
extensions \(\Gamma\) donne la naturalité
\eqref{eq:pdfv2-sol-naturality-d}. Enfin,
\(\mathcal G_{\rm sol}\), \(\mathfrak M_{\rm sol}\) et
\(\mathcal C_{\rm sol}^{\rm HMT}\) sont des graphes, et
\(X_{\chi,{\rm sol}}\) est un objet dépendant qui conserve explicitement
son antécédent. La première projection de chaque étape récupère l'étape
précédente. Ainsi aucune composition n'efface l'histoire, la mémoire, la multisection,
le terminal ni le lecteur.
\end{proof}

\section{Provenance probatoire et falsificateurs}
\label{sec:pdfv2-sol-provenance-falsifiers}

\paragraph{Résultat récupéré.}
Appartiennent à l'architecture déjà construite l'accroissement signé avec retenue,
la séparation \(b=b^+-b^-\), les mémoires \(M^\pm,D^0,H^0\), le couple
\(E_9,J_9\), la retenue de Jensen \(\kappa\), l'extension des cylindres, la
ligne de pertes, le terminal, le Gram, Stein et le télescopage.

\paragraph{Formalisation nouvelle.}
Sont des formalisations explicites de cette architecture le typage préalable à toute
évaluation extérieure, la factorisation
\(\mathscr H^{\rm depth}\widehat\otimes\mathscr H^{\rm res}\), la séparation
orthogonale qui empêche le double comptage, le terminal libre sur des histoires à
treize coordonnées et les quatre mises en paquet non ablatives par graphes.

\begin{criterion}[Falsificateurs internes de la branche \(\mathrm{sol}\)]
\label{crit:pdfv2-sol-falsifiers}
La construction échoue si l'un quelconque des faits suivants se produit :
\begin{enumerate}
 \item \(b\neq b^+-b^-\) o \(|b|\neq b^++b^-\);
 \item \(E_9J_9\neq I\), la mesure n'est pas projective ou \(P_9\) n'est pas un
       projecteur ;
 \item \(\kappa<0\) ou
       \(E_9b_f^\pm=b_c^\pm+\kappa\) échouent ;
 \item \(Q^{\rm micro}Q^{\rm shell}\neq0\), ou un même secteur apparaît dans
       les deux termes de la ligne \(L\) ;
 \item l'une de \(b^+,b^-,\kappa,Q^{\rm micro},Q^{\rm shell}\) est omise ou
       comptée deux fois ;
 \item le terminal conserve seulement la valeur visible et non l'histoire avec ses
       treize coordonnées ;
 \item l'identité de Stein échoue ou le télescopage contient un signe
       contraire à une somme de carrés ;
 \item une extension ne commute pas avec \(E_9,J_9,L,G,U\) ;
 \item la multisection des continuations est remplacée par un prolongement
       choisi rétrospectivement ;
 \item l'une quelconque des treize coordonnées est omise.
\end{enumerate}
Dans le dernier cas, le diagnostic contraignant est : \emph{objet substitué ;
conclusion non transférable}.
\end{criterion}

\begin{warningbox}
Ce chapitre conclut uniquement la construction interne de
\(\mathcal C_{\rm sol}^{\rm HMT}\). Il ne contient pas d'équation différentielle
classique, de données initiales classiques, d'application à partir de données extérieures ni de
conclusion sur un problème extérieur. Ces entrelaceurs appartiennent à un autre
document et ne gouvernent pas l'existence de cette macrostructure.
\end{warningbox}
\fi

\endgroup

\section{Orientation, retenue et composition simultanée des trajectoires}

L'inversion du parcours agit une seule fois sur \(\mathfrak e_\alpha\). Soit \(\mathsf J(v_\Sigma,v_\Pi):=(v_\Pi,v_\Sigma)\) l'échange des feuillets, et soit \(P_{\rm tr}\) la transposition des matrices de ports. Ses cinq expressions locales sont
\begin{equation}
\label{eq:pdfv2-cor-orientacion-cruzada}
 \begin{gathered}
 \mathsf J\sigma_k=-\sigma_k\mathsf J,
 \qquad b_{\bar\alpha}=-b_\alpha,
 \qquad (b_{\bar\alpha}^+,b_{\bar\alpha}^-)
       =(b_\alpha^-,b_\alpha^+),\\
 \vartheta_L(K,\epsilon):=
 \begin{cases}
  (L,-\epsilon),&K=L,\\
  (K,\epsilon),&K\neq L,
 \end{cases}
 \qquad
 \kappa_{J,I}(v,u)=-P_{\rm tr}\kappa_{I,J}(u,v),\\
 \kappa(\bar\alpha)
 =-\mathsf C_\alpha^{-1}\kappa(\alpha),
 \qquad v_{\bar\alpha}=-v_\alpha.
 \end{gathered}
\end{equation}
Ici, \(\mathsf J\) échange les deux feuillets, \(P_{\rm tr}\) transpose la matrice d'incidence et \(\vartheta_L\) échange exactement les deux faces de la direction \(L\) en fixant les six autres. Dans la base libre de \(\mathscr B=\mathscr I\times\{+,-\}\), les huit vecteurs sont distincts, \(\vartheta_L^2=I\) et \(\vartheta_L\neq\vartheta_{L'}\) pour \(L\neq L'\). La même orientation explique tous les signes de \eqref{eq:pdfv2-cor-orientacion-cruzada}.

Pour des extensions composables \(\alpha:x\to y\) et \(\beta:y\to z\), la retenue native satisfait
\begin{equation}
\label{eq:pdfv2-cor-carry-nativo}
 \kappa(\beta\alpha)=
 \kappa(\beta)+\mathsf C(\beta)\kappa(\alpha).
\end{equation}
Le relèvement matriciel ne se réduit pas à un nom. Dans la base ordonnée \((f,r,e,b,g)\), soit \(\widetilde L_N(\alpha)\in\operatorname{Mat}_5(\mathbb Z)\) le représentant entier centré de la matrice de \(\Psi_\alpha\) de \eqref{eq:pdfv2-cor-psi-generadores} modulo \(3^N\), et définissons
\begin{equation}
\label{eq:pdfv2-cor-carry-matricial-def}
 c_N(\beta,\alpha):=
 \frac{\widetilde L_N(\beta)\widetilde L_N(\alpha)
       -\widetilde L_N(\beta\alpha)}{3^N}
 \in\operatorname{Mat}_5(\mathbb Z).
\end{equation}
La divisibilité provient de l'égalité des deux compositions modulo \(3^N\). En écrivant \(L_N=[\widetilde L_N]_{3^N}\), les relèvements entiers satisfont
\begin{equation}
\label{eq:pdfv2-cor-carry-matricial}
 \widetilde L_N(\beta)\widetilde L_N(\alpha)
 =\widetilde L_N(\beta\alpha)+3^Nc_N(\beta,\alpha)
\end{equation}
et l'associativité produit le cocycle
\begin{equation}
\label{eq:pdfv2-cor-cociclo-matricial}
 c_N(\delta,\beta\alpha)+\widetilde L_N(\delta)c_N(\beta,\alpha)
 =c_N(\delta\beta,\alpha)
  +c_N(\delta,\beta)\widetilde L_N(\alpha).
\end{equation}

Le même accroissement de retenue agit sur le complexe local au moyen de
\begin{equation}
\label{eq:pdfv2-cor-psi-generadores}
 \begin{gathered}
 \Psi_\alpha(r_x)=r_y,\quad
 \Psi_\alpha(e_x)=e_y,\quad
 \Psi_\alpha(g_x)=g_y,\quad
 \Psi_\alpha(f_x)=f_y,\\
 \Psi_\alpha(b_x)=b_y+3(c_y-c_x)e_y.
 \end{gathered}
\end{equation}
Le changement de frontière orientée est
\begin{equation}
\label{eq:pdfv2-cor-k-alpha}
 K_\alpha:=(v_y-v_x)f_y,
 \qquad
 K_{\beta\alpha}=K_\beta+\Psi_\beta K_\alpha.
\end{equation}
Les équations \eqref{eq:pdfv2-cor-carry-nativo}, \eqref{eq:pdfv2-cor-cociclo-matricial} et \eqref{eq:pdfv2-cor-k-alpha} sont trois degrés du même registre de composition. Aucune n'autorise à effacer les deux autres.

\subsection{Nerf des prolongements et diagonale d'Alexander--Whitney}
Pour une histoire du registre commun, on écrit \(A_n:=\mathbf Z/3^n\mathbf Z\) et \(\mathsf T_m(x)\) pour la catégorie finie de \emph{tous} ses préfixes TPK de profondeur \(m\), munis des morphismes d'extension déjà construits. La diagonale et la counité qui suivent s'appliquent à ce même nerf ; elles ne définissent pas une cinquième branche autonome.
\begingroup
\let\section\subsection
% Extracción quirúrgica del propietario V2 05e: líneas 68--110.
% El resto permanece en este archivo como procedencia, pero no se publica.
\iffalse
\chapter[Macrostructure déterminantielle interne]{Macrostructure déterminantielle interne engendrée par le continu discret}
\label{chap:detint-macroestructura}

\begin{thesisbox}
La branche \({\rm det}\) part exclusivement d'histoires survivantes de
\(\mathcal C_{\rm cont}^{\rm disc}\). Le nerf des extensions, la diagonale
Alexander--Whitney, le complexe local, l'unité, les deux publications, le
filler, la réduction terminale et la tour \(3\)-adique se construisent avant
toute évaluation ultérieure. Les conditions de bilan et d'acyclicité sont
incorporées au domaine typé et ne sont jamais tacitement présupposées.
\end{thesisbox}

\section{Domaine déterminantiel survivant}

Soit
\begin{equation}\label{eq:detint-an}
 A_n:=\mathbf Z/3^n\mathbf Z,
 \qquad
 \rho_{n+1,n}:A_{n+1}\longrightarrow A_n.
\end{equation}
Pour \(x\in\mathcal C_{\rm cont}^{\rm disc}\), soit
\(\mathsf T_m(x)\) la catégorie finie des préfixes TPK de profondeur \(m\)
compatibles avec \(x\). Chaque objet conserve feuillet, orientation, résidu,
quotient, retenue entière, mémoire, frontière, chemin, multisection, survie
et terminal ; ses morphismes sont les extensions admissibles qui conservent ces
coordonnées.

De chaque objet \(y\in\mathsf T_m(x)\), on extrait les entiers internes
\begin{equation}\label{eq:detint-c-v-lifts}
 \widetilde c_y,\widetilde v_y\in\mathbf Z,
 \qquad
 c_{y,n}:=\widetilde c_y\bmod3^n,
 \qquad
 v_{y,n}:=\widetilde v_y\bmod3^n.
\end{equation}
La première coordonnée est la retenue non orientée et la seconde, le coefficient
orienté de frontière. Le renversement satisfait
\begin{equation}\label{eq:detint-or-cv}
 c_{\bar y,n}=c_{y,n},
 \qquad
 v_{\bar y,n}=-v_{y,n}.
\end{equation}

On écrit
\begin{equation}\label{eq:detint-z3-k3}
 \mathbf Z_3:=\varprojlim_n A_n,
 \qquad
 K_3:=\operatorname{Frac}(\mathbf Z_3).
\end{equation}
La restriction déterminantielle est définie par
\begin{equation}\label{eq:detint-dominio}
\begin{split}
 \mathcal C_{\rm cont}^{\rm det}:=
 \bigl\{x\in\mathcal C_{\rm cont}^{\rm disc}:{}&
 x\text{ survit à toute profondeur};\\
 &\widehat P_x^{\rm red}\text{ est fini et basé};\\
 &\chi(\widehat P_x^{\rm red})=0;\\
 &H_*(\widehat P_x^{\rm red}\otimes_{\mathbf Z_3}K_3)=0;\\
 &\text{la réduction canonique produit un bloc carré}\\
 &\text{stable par raffinement}\bigr\}.
\end{split}
\end{equation}
Les objets \(\widehat P_x^{\rm red}\) et le bloc carré seront construits
dans \cref{sec:detint-terminal}. Par conséquent,
\eqref{eq:detint-dominio} est une condition d'appartenance vérifiable et non
une conclusion insérée dans la définition des applications.

\fi
\section{Nerf des chemins, diagonale d'Alexander--Whitney et counité}

On définit le complexe normalisé
\begin{equation}\label{eq:detint-gmn}
 G_{m,n}(x):=
 N_*^{\rm norm}\!\left(N\mathsf T_m(x);A_n\right).
\end{equation}
Pour un simplexe non dégénéré
\(\sigma=[y_0\to y_1\to\cdots\to y_q]\), la diagonale est
\begin{equation}\label{eq:detint-aw}
 \Delta_{\rm AW}(\sigma)
 =\sum_{j=0}^q
 [y_0\to\cdots\to y_j]\otimes
 [y_j\to\cdots\to y_q].
\end{equation}
La counité est donnée par
\begin{equation}\label{eq:detint-counidad-g}
 \epsilon_G([y])=1,
 \qquad
 \epsilon_G(\sigma)=0\quad(q>0).
\end{equation}
En particulier,
\begin{equation}\label{eq:detint-vertice-grouplike}
 \Delta_{\rm AW}[y]=[y]\otimes[y],
 \qquad
 \epsilon_G([y])=1.
\end{equation}

Si \(m'\ge m\) et \(n'\ge n\), la troncature et la réduction des coefficients
induisent
\begin{equation}\label{eq:detint-rmn}
 R_{m',m}^{n',n}:G_{m',n'}(x)\longrightarrow G_{m,n}(x),
\end{equation}
et l'on vérifie
\begin{equation}\label{eq:detint-aw-natural}
 R\partial=\partial R,
 \qquad
 (R\otimes R)\Delta_{\rm AW}=\Delta_{\rm AW}R,
 \qquad
 \epsilon_GR=\rho_{n',n}\epsilon_G.
\end{equation}
Le chemin ne se réduit pas à ses extrémités : le simplexe entier reste dans
\(G_{m,n}(x)\).

\iffalse
\section{Complexe local interne}

Pour \(c\in A_n\), on définit
\begin{equation}\label{eq:detint-p0-grados}
 P^0_{n,1}(c)=A_nf,
 \qquad
 P^0_{n,0}(c)=A_nr\oplus A_ne\oplus A_nb,
 \qquad
 P^0_{n,-1}(c)=A_ng,
\end{equation}
avec différentielle
\begin{equation}\label{eq:detint-p0-diferencial}
 \partial f=3ce+b,
 \qquad
 \partial r=0,
 \qquad
 \partial e=g,
 \qquad
 \partial b=-3cg.
\end{equation}
L'identité de chaîne est littérale :
\begin{equation}\label{eq:detint-d2}
 \partial^2f=3c\,\partial e+\partial b=3cg-3cg=0,
 \qquad
 \partial^2e=\partial^2b=\partial^2r=0.
\end{equation}

L'augmentation
\begin{equation}\label{eq:detint-epsilon-p}
 \epsilon_P:P_n^0(c)\longrightarrow A_n[0]
\end{equation}
est fixée par
\[
 \epsilon_P(r)=1,
 \qquad
 \epsilon_P(e)=\epsilon_P(b)=\epsilon_P(f)=\epsilon_P(g)=0.
\]
Le complexe possède la coalgèbre différentielle
\begin{equation}\label{eq:detint-coalgebra-p}
 \Delta_P(r)=r\otimes r,
 \qquad
 \Delta_P(q)=q\otimes r+r\otimes q
 \quad(q\in\{e,b,f,g\}).
\end{equation}
Ainsi, \(r\) est group-like et les quatre autres générateurs sont primitifs
relativement à \(r\). Les équations
\eqref{eq:detint-p0-diferencial} prouvent que \(\Delta_P\) commute avec la
différentielle tensorielle.

\section{Système local d'extensions et ledger de mémoire}

Pour une extension \(\alpha:y\to y'\), nous définissons
\begin{equation}\label{eq:detint-psi}
 \Psi_\alpha:P_n^0(c_y)\longrightarrow P_n^0(c_{y'})
\end{equation}
par
\begin{equation}\label{eq:detint-psi-generadores}
\begin{gathered}
 \Psi_\alpha(r_y)=r_{y'},\quad
 \Psi_\alpha(e_y)=e_{y'},\quad
 \Psi_\alpha(g_y)=g_{y'},\quad
 \Psi_\alpha(f_y)=f_{y'},\\
 \Psi_\alpha(b_y)=b_{y'}+3(c_{y'}-c_y)e_{y'}.
\end{gathered}
\end{equation}
C'est une application de complexes parce que
\begin{equation}\label{eq:detint-psi-b-chain}
 \partial\Psi_\alpha(b_y)
 =-3c_{y'}g+3(c_{y'}-c_y)g
 =-3c_yg
 =\Psi_\alpha(\partial b_y),
\end{equation}
et
\begin{equation}\label{eq:detint-psi-f-chain}
 \Psi_\alpha(\partial f_y)
 =3c_ye+b+3(c_{y'}-c_y)e
 =3c_{y'}e+b
 =\partial\Psi_\alpha(f_y).
\end{equation}
De plus,
\begin{equation}\label{eq:detint-psi-functorial}
 \Psi_\beta\Psi_\alpha=\Psi_{\beta\alpha}.
\end{equation}

Le changement de la coordonnée orientée n'est pas effacé. Il est enregistré par
\begin{equation}\label{eq:detint-k-alpha}
 K_\alpha:=(v_{y'}-v_y)f_{y'}.
\end{equation}
Pour des extensions composables,
\begin{equation}\label{eq:detint-k-coherencia}
 K_{\beta\alpha}=K_\beta+\Psi_\beta K_\alpha.
\end{equation}
Cette égalité est le ledger exact de l'accroissement de frontière et rend visible
la mémoire qu'une composition purement nominale perdrait.

Le diagramme local complet est la construction de Grothendieck
\begin{equation}\label{eq:detint-grothendieck-local}
 \int_{y\in\mathsf T_m(x)}P_n^0(c_y).
\end{equation}
Dans celle-ci, une flèche au-dessus de \(\alpha:y\to y'\) porte comme composante linéaire
\(\Psi_\alpha\) et comme cohérence \(K_\alpha\). Lorsqu'elle est appliquée à un
élément supporté en un sommet, on utilise la notation
\begin{equation}\label{eq:detint-psi-soporte-vertice}
 \Psi_\alpha(p,[y]):=(\Psi_\alpha p,[y']).
\end{equation}
L'arête \([y\to y']\) reste dans la coordonnée de chemin du nerf. Cette
convention sera celle utilisée dans les formules de naturalité de
\(z_\pm\) et \(h_*\) ; on n'affirme pas une application globale qui efface les autres
simplexes de \(G_{m,n}(x)\).

\section{Pullback, unité, racine, publications et filler}

Le pullback de complexes est
\begin{equation}\label{eq:detint-pgen}
 P^{\rm gen}_{y;m,n}
 :=P_n^0(c_y)\times_{A_n[0]}G_{m,n}(x)
 =\{(p,\gamma):\epsilon_P(p)=\epsilon_G(\gamma)\}.
\end{equation}
Sa différentielle est
\begin{equation}\label{eq:detint-pgen-d}
 \partial(p,\gamma)=(\partial p,\partial\gamma).
\end{equation}
L'unité commune associée au sommet \(y\) est
\begin{equation}\label{eq:detint-unidad}
 \mathbf1_y:=(r,[y]).
\end{equation}
Ses deux projections sont group-like par
\eqref{eq:detint-vertice-grouplike} et
\eqref{eq:detint-coalgebra-p} ; c'est le sens précis d'unité
group-like compatible du pullback.

La projection racine est l'application de complexes
\begin{equation}\label{eq:detint-proyeccion-root}
 \varpi_{\rm root}:P^{\rm gen}_{y;m,n}\longrightarrow
 \operatorname{Root}_{y;m,n},
 \qquad
 \varpi_{\rm root}(p,\gamma):=(\epsilon_P(p)r,\gamma).
\end{equation}
On définit les deux publications internes
\begin{equation}\label{eq:detint-zpm}
 z_-(y):=(r,[y]),
 \qquad
 z_+(y):=(r+3c_yv_ye+v_yb,[y]).
\end{equation}
Toutes deux sont des cycles :
\begin{equation}\label{eq:detint-zpm-ciclos}
 \partial z_-(y)=0,
 \qquad
 \partial z_+(y)=3c_yv_yg-3c_yv_yg=0.
\end{equation}
Leurs racines coïncident :
\begin{equation}\label{eq:detint-raiz-comun}
 \varpi_{\rm root}z_-(y)
 =(r,[y])
 =\varpi_{\rm root}z_+(y).
\end{equation}

Le filler se construit directement :
\begin{equation}\label{eq:detint-hstar}
 h_*(y):=(-v_yf,0)\in P^{\rm gen}_{y;m,n,1}.
\end{equation}
Alors
\begin{equation}\label{eq:detint-filler-identidad}
 \partial h_*(y)
 =(-3c_yv_ye-v_yb,0)
 =z_-(y)-z_+(y).
\end{equation}
En particulier,
\begin{equation}\label{eq:detint-clases-iguales}
 [z_-(y)]=[z_+(y)]
 \quad\text{dans}\quad H_0(P^{\rm gen}_{y;m,n}),
\end{equation}
et l'égalité conserve sa chaîne témoin.

Sous une extension \(\alpha:y\to y'\), on a
\begin{equation}\label{eq:detint-z-natural}
 z_-(y')=\Psi_\alpha z_-(y),
 \qquad
 z_+(y')-\Psi_\alpha z_+(y)=\partial K_\alpha,
\end{equation}
et
\begin{equation}\label{eq:detint-h-natural}
 h_*(y')-\Psi_\alpha h_*(y)=-K_\alpha.
\end{equation}
Par conséquent, la naturalité n'identifie pas des valeurs distinctes de \(v\) : elle enregistre
leur différence au moyen de l'homotopie cohérente \(K_\alpha\).

\section{Orientation interne}

Sur \(P_n^0(c)\), on définit
\begin{equation}\label{eq:detint-omega-p}
 \Omega(r)=r,
 \qquad
 \Omega(q)=-q\quad(q\in\{e,b,f,g\}).
\end{equation}
C'est un automorphisme de complexes et de la coalgèbre relative. Dans le nerf,
l'inversion orientée est
\begin{equation}\label{eq:detint-omega-g}
 \mathfrak o_G[y_0\to\cdots\to y_q]
 :=(-1)^{q(q+1)/2}
 [\bar y_q\to\cdots\to\bar y_0].
\end{equation}
Avec \eqref{eq:detint-or-cv}, ces actions donnent
\begin{equation}\label{eq:detint-or-z-h}
 \Omega z_-(y)=z_-(\bar y),
 \qquad
 \Omega z_+(c_y,v_y)=z_+(c_{\bar y},v_{\bar y}),
 \qquad
 \Omega h_*(y)=h_*(\bar y).
\end{equation}
L'orientation agit sur le chemin, le complexe, les publications, le filler et la droite
déterminantielle ; ce n'est pas une étiquette sans action.

\section{Réduction terminale et condition de survie}
\label{sec:detint-terminal}

Seule la racine commune est supprimée :
\begin{equation}\label{eq:detint-p-reducido}
 \overline P_{y;m,n}
 :=P^{\rm gen}_{y;m,n}/A_n\mathbf1_y,
 \qquad
 \widehat P_{y;m}^{\rm red}:=\varprojlim_n\overline P_{y;m,n}.
\end{equation}
La base est ordonnée d'abord selon l'ordre TPK des chemins, puis selon
\begin{equation}\label{eq:detint-orden-base}
 f\prec e\prec b\prec g.
\end{equation}
On exécute l'algorithme interne suivant.
\begin{enumerate}
 \item Chercher les coefficients qui sont des unités de \(\mathbf Z_3\).
 \item Choisir le premier selon \eqref{eq:detint-orden-base} et l'ordre des
 chemins.
 \item Annuler le couple correspondant et enregistrer les opérations
 élémentaires, leur signe et leur contribution à l'orientation.
 \item Répéter jusqu'à ce qu'il ne reste aucun pivot unitaire annulable.
\end{enumerate}
La condition de survie de \eqref{eq:detint-dominio} exige que,
cofinalement en \(m\), le résultat soit un bloc à deux termes
\begin{equation}\label{eq:detint-sx}
 \mathbf Z_3^{r_x}\xrightarrow{S_x}\mathbf Z_3^{r_x},
 \qquad
 \det S_x\ne0\text{ dans }K_3,
\end{equation}
et qu'un raffinement n'ajoute que des blocs contractiles unitaires et des changements
de base enregistrés. Si ces conditions échouent, l'histoire n'appartient pas à
\(\mathcal C_{\rm cont}^{\rm det}\) ; on ne force pas un terminal inexistant.

Le chemin identité fournit un témoin élémentaire. Après le retrait de la racine,
le complexe local est
\begin{equation}\label{eq:detint-testigo-contractil}
 A_nf\xrightarrow{\binom{3c}{1}}
 A_ne\oplus A_nb\xrightarrow{(1,-3c)}A_ng.
\end{equation}
Il est contractile par
\begin{equation}\label{eq:detint-contraccion-local}
 s(g)=e,
 \qquad
 s(ae+\beta b)=\beta f,
 \qquad
 s(f)=0.
\end{equation}
Ainsi, les conditions terminales admettent au moins le témoin identité
interne et ne décrivent pas un sous-domaine formellement vide.

\section{Smith, ordre \texorpdfstring{\(3\)}{3}-adique et unité initiale}

Dans des bases orientées, une diagonalisation par valuation a la forme
\begin{equation}\label{eq:detint-smith}
 U_xS_xV_x
 =\operatorname{diag}(3^{a_1}u_1,\ldots,3^{a_r}u_r),
 \qquad
 a_i\in\mathbf N,quad u_i\in\mathbf Z_3^\times.
\end{equation}
On définit
\begin{equation}\label{eq:detint-orden-d}
 d_x:=v_3(\det S_x)=\sum_{i=1}^r a_i
\end{equation}
et l'unité initiale
\begin{equation}\label{eq:detint-leading-unit}
 \lambda_x:=3^{-d_x}\det S_x\in\mathbf Z_3^\times.
\end{equation}
La coordonnée d'orientation trivialise la droite déterminantielle. Si une
transition produit
\begin{equation}\label{eq:detint-estabilidad-bloques}
 S_x\oplus I_q=U S_{x'}V,
\end{equation}
l'orientation est transportée par le facteur inverse de \(\det U\det V\).
L'unité orientée \(\lambda_x^{\rm or}\) est alors naturelle et ne change pas
lors de l'ajout de blocs contractiles.

En précision finie,
\begin{equation}\label{eq:detint-smith-finito}
 S_{x,n}=S_x\bmod3^n,
 \qquad
 a_i^{(n)}=\min(a_i,n).
\end{equation}
Le terminal est donc la famille compatible de toutes les précisions et
non un entier isolé.

\section{Tour de Bockstein et retenue d'ordre supérieur}

Soit
\[
 C_x^1\xrightarrow{S_x}C_x^0
\]
le bloc terminal. Si une classe \([\bar y]\) admet un relèvement
\(y\in C_x^1\) avec \(S_xy\in3^qC_x^0\), nous définissons
\begin{equation}\label{eq:detint-beta-q}
 \beta_q[\bar y]
 :=\left[3^{-q}S_xy\bmod3\right].
\end{equation}
Le changement de relèvement modifie le représentant par un bord de la page
correspondante, de sorte que \(\beta_q\) est bien défini. La retenue
supérieure est
\begin{equation}\label{eq:detint-carry-superior}
 \operatorname{Carr}_q(y):=3^{-q}S_xy.
\end{equation}
Elle ne remplace pas la retenue TPK de \eqref{eq:detint-c-v-lifts} ; elle enregistre son
prolongement dans la tour des coefficients.

Pour un bloc \(3^{a_i}u_i\), la classe persiste jusqu'à la page \(a_i\),
et à cette page, la différentielle non nulle est la multiplication par
\(\bar u_i\in\mathbf F_3^\times\). Par conséquent,
\begin{equation}\label{eq:detint-bock-longitudes}
 \{\text{longueurs de survie}\}=\{a_1,\ldots,a_r\},
 \qquad
 d_x=\sum_i a_i.
\end{equation}
Le produit orienté des trivialisations de page satisfait
\begin{equation}\label{eq:detint-bock-leading}
 \Theta_{\rm Bock}(x)=\overline{\lambda_x^{\rm or}}
 \in\mathbf F_3^\times.
\end{equation}
La preuve se fait dans chaque bloc diagonal et se transporte par les changements
unimodulaires enregistrés. Smith, Bockstein et l'unité initiale sont ainsi trois
lectures du même terminal.

\section{Les treize coordonnées de la restriction déterminantielle}

Pour chaque \((m,n)\), on définit
\begin{equation}\label{eq:detint-d-trece}
 D_{{\rm det},m,n}(x):=
 (\mathcal F,\operatorname{Ext}_\Gamma,\operatorname{Rel},
 \mathcal N,\operatorname{or},\operatorname{Carr},\operatorname{Mem},
 \partial,\gamma,\operatorname{Msect},\operatorname{Surv},
 \operatorname{Term},\mathcal L)_{{\rm det},m,n}(x),
\end{equation}
avec le contenu suivant.
\begin{enumerate}
 \item La fibre est
 \[
  \mathcal F=(A_n,\mathsf T_m(x),
  \{P_n^0(c_y)\}_y,\{P^{\rm gen}_{y;m,n}\}_y).
 \]
 \item \(\operatorname{Ext}_\Gamma\) contient tous les morphismes
 \(\alpha\), les applications \(\Psi_\alpha\), les homotopies \(K_\alpha\) et
 les applications \(R_{m',m}^{n',n}\).
 \item \(\operatorname{Rel}\) contient \(\partial^2=0\), l'équation de
 pullback, AW, la counité, \(\Psi_\beta\Psi_\alpha=\Psi_{\beta\alpha}\) et
 \(K_{\beta\alpha}=K_\beta+\Psi_\beta K_\alpha\).
 \item
 \(\mathcal N=(m,n,\widetilde c_y,\widetilde v_y,c_{y,n},v_{y,n},
 r_x,a_1,\ldots,a_{r_x},d_x)\).
 \item \(\operatorname{or}\) contient \(y\mapsto\bar y\),
 \(\mathfrak o_G\), \(\Omega\) et l'orientation de la droite
 déterminantielle.
 \item \(\operatorname{Carr}\) conserve séparément la retenue entière,
 sa réduction \(c_{y,n}\) et toutes les retenues supérieures
 \(\operatorname{Carr}_q\).
 \item \(\operatorname{Mem}\) contiene
 \[
  y_0\to\cdots\to y_m,
  \quad\{(\widetilde c_{y_j},\widetilde v_{y_j},\mathbf1_{y_j})\}_{j=0}^m,
  \quad\{K_{\alpha_j}\}.
 \]
 \item
 \(\partial=(\partial_G,\partial_P,\partial_{P^{\rm gen}},
 \partial_{\overline P},h_*)\).
 \item \(\gamma=[y_0\to\cdots\to y_q]\) conserve le simplexe orienté
 complet.
 \item \(\operatorname{Msect}\) contient toutes les extensions
 compatibles, tous les relèvements \(3\)-adiques et toutes les présentations
 unimodulaires du même terminal orienté ; il n'en choisit aucune a posteriori.
 \item \(\operatorname{Surv}\) enregistre la profondeur arbitraire,
 la compatibilité des coefficients, le bilan, l'acyclicité sur \(K_3\),
 la stabilité par blocs unitaires et la permanence de la racine.
 \item
 \[
  \operatorname{Term}=(\widehat P_x^{\rm red},S_x,
  \{a_i,u_i\},d_x,\lambda_x^{\rm or},
  \{\beta_q\},\Theta_{\rm Bock}).
 \]
 \item Le lecteur interne est
 \[
  \mathcal L_{\rm det}^{\rm int}
  =(\mathbf1_y,
  \varpi_{\rm root}z_-=\varpi_{\rm root}z_+,
  [z_-]=[z_+],d_x,\lambda_x^{\rm or},\Theta_{\rm Bock}).
 \]
\end{enumerate}

Les transitions satisfont, coordonnée par coordonnée,
\begin{equation}\label{eq:detint-naturalidad-trece}
 u^{\rm det}_{(m',n'),(m,n)}D_{{\rm det},m',n'}(x')
 =D_{{\rm det},m,n}(\tau_{m',m}x').
\end{equation}

\section{Assembleur, publicateur et chaîne non ablative}

On définit
\begin{equation}\label{eq:detint-graph-d}
 \mathcal G_{\rm det}:=\operatorname{Graph}(D_{\rm det}),
 \qquad
 \operatorname{Res}_{\rm det}(x):=(x,D_{\rm det}x).
\end{equation}
Soit \(\chi_{\rm det}(x)\) la multisection complète de germes
\[
 (y,\widetilde c_y,\widetilde v_y,\operatorname{or}_y,\gamma_y)
\]
compatibles à toute profondeur. L'objet dépendant est
\begin{equation}\label{eq:detint-x-dependiente}
 X_{\chi,{\rm det}}
 :=\{(\chi_{\rm det}(x),x,D_{\rm det}x):
 x\in\mathcal C_{\rm cont}^{\rm det}\},
\end{equation}
et
\begin{equation}\label{eq:detint-prx}
 \operatorname{pr}_{X,{\rm det}}(x,D_{\rm det}x)
 :=(\chi_{\rm det}(x),x,D_{\rm det}x).
\end{equation}

L'assembleur brut
\begin{equation}\label{eq:detint-ensamblador-crudo}
 a_{\rm det}:X_{\chi,{\rm det}}\longrightarrow
 \mathscr M_{\rm det}^{\rm amb}
\end{equation}
produit
\begin{equation}\label{eq:detint-ensamblador-salida}
\begin{split}
 a_{\rm det}(z)=\bigl(&
 \{G_{m,n},\Delta_{\rm AW},\epsilon_G\}_{m,n};\\
 &\{P_n^0(c),\Psi_\alpha,K_\alpha,P^{\rm gen},
 \mathbf1,z_-,z_+,h_*\}_{m,n};\\
 &\{\widehat P^{\rm red},S,d,\lambda^{\rm or},
 \beta_q,\Theta_{\rm Bock}\}\bigr).
\end{split}
\end{equation}
Le macro-objet conserve son entrée :
\begin{equation}\label{eq:detint-graph-a}
 \mathfrak M_{\rm det}:=\operatorname{Graph}(a_{\rm det}),
 \qquad
 \mathcal A_{\rm det}(z):=(z,a_{\rm det}z).
\end{equation}

Le publicateur brut
\begin{equation}\label{eq:detint-publicador-crudo}
 p_{\rm det}:\mathfrak M_{\rm det}\longrightarrow\mathscr Y_{\rm det}
\end{equation}
publie conjointement le certificat
\begin{equation}\label{eq:detint-certificado-publicado}
\begin{gathered}
 \partial^2=0,quad \text{AW et counité},quad
 \mathbf1_y\text{ group-like compatible},\\
 \partial z_\pm=0,quad
 \varpi_{\rm root}z_-=\varpi_{\rm root}z_+,quad
 \partial h_*=z_--z_+,\\
 \Psi_\beta\Psi_\alpha=\Psi_{\beta\alpha},quad
 K_{\beta\alpha}=K_\beta+\Psi_\beta K_\alpha,\\
 d=v_3(\det S)=\sum_i a_i,quad
 \Theta_{\rm Bock}=\overline{\lambda^{\rm or}},quad
 \text{naturalité complète par raffinement}.
\end{gathered}
\end{equation}
Enfin,
\begin{equation}\label{eq:detint-graph-p}
 \mathcal C_{\rm det}^{\rm HMT}:=\operatorname{Graph}(p_{\rm det}),
 \qquad
 \mathcal P_{\rm det}(M):=(M,p_{\rm det}M).
\end{equation}
La chaîne intégrale est
\begin{equation}\label{eq:detint-cadena-completa}
 \mathcal C_{\rm cont}^{\rm disc}\supset
 \mathcal C_{\rm cont}^{\rm det}
 \xrightarrow{\operatorname{Res}_{\rm det}}\mathcal G_{\rm det}
 \xrightarrow{\operatorname{pr}_{X,{\rm det}}}X_{\chi,{\rm det}}
 \xrightarrow{\mathcal A_{\rm det}}\mathfrak M_{\rm det}
 \xrightarrow{\mathcal P_{\rm det}}\mathcal C_{\rm det}^{\rm HMT}.
\end{equation}
Chaque flèche possède un inverse sur son image donné par une projection de
coordonnées. On ne remplace pas l'histoire par une valeur, la multisection par un sélecteur,
le complexe par un déterminant ni la preuve par un nom.

\section{Théorème de réalisation déterminantielle interne}

\begin{theorem}[Réalisation déterminantielle interne complète]
\label{thm:detint-realizacion}
Pour tout \(x\in\mathcal C_{\rm cont}^{\rm det}\), la chaîne
\eqref{eq:detint-cadena-completa} construit un macro-objet naturel en
profondeur et en précision. Il contient deux cycles à racine commune, un filler
explicite qui les identifie, une unité group-like compatible et un terminal
\(3\)-adique dont l'ordre, la forme de Smith et la tour de Bockstein produisent le même
lecteur orienté. Les treize coordonnées de \eqref{eq:detint-d-trece}
restent récupérables en sortie.
\end{theorem}

\begin{proof}
Les équations \eqref{eq:detint-p0-diferencial}--
\eqref{eq:detint-d2} prouvent que \(P_n^0(c)\) est un complexe. Les deux
augmentations sont des applications de complexes, de sorte que le pullback est typé.
Alexander--Whitney est naturel et rend chaque sommet group-like ; avec
\(r\), cela donne l'unité commune. Le calcul direct
\eqref{eq:detint-zpm-ciclos}--\eqref{eq:detint-filler-identidad} prouve que
les publications sont des cycles, ont la même racine et sont reliées par
\(h_*\). Les formules de \(\Psi_\alpha\) prouvent la fonctorialité, et
\(K_\alpha\) prouve la naturalité homotopique en conservant l'accroissement de
mémoire. La définition de \(\mathcal C_{\rm cont}^{\rm det}\) garantit que
la réduction terminale produit un bloc carré sans le postuler hors du
domaine. Smith sur \(\mathbf Z_3\) donne \(d=\sum_i a_i\) ; le calcul de chaque
bloc \(3^{a_i}u_i\) donne
\(\Theta_{\rm Bock}=\overline{\lambda^{\rm or}}\). Troncature, réduction et
ajout de blocs unitaires conservent ces identités. Enfin, les
trois constructions par graphe retiennent l'entrée comme première coordonnée,
de sorte que la composition complète est non ablative.
\end{proof}

\section{Provenance et falsificateurs}

\paragraph{Provenance.}
Le nerf TPK, Alexander--Whitney, la counité, l'unité commune, la forme
normale \(\partial f=3ce+b\), \(\partial e=g\),
\(\partial b=-3cg\), la racine commune, le filler et l'architecture de tour
ont le statut \texttt{ARCHITECTURE\_AUTORALE\_PRÉEXISTANTE}. Le calcul
littéral des deux cycles, de leur racine et de l'identité
\(\partial h_*=z_--z_+\) a le statut \texttt{RÉSULTAT\_RÉCUPÉRÉ}. Le
système local \(\Psi_\alpha/K_\alpha\), le sous-domaine survivant
équilibré, le terminal exclusivement \(3\)-adique et la mise en paquet par
graphes ont le statut \texttt{FORMALISATION\_NOUVELLE} ; ils formalisent sans changer
la provenance conceptuelle de l'appareil.

\begin{ablation}[Falsificateurs internes de la branche déterminantielle]
\label{abl:detint-falsadores}
La construction est falsifiée au niveau correspondant si l'un des
faits suivants se produit.
\begin{enumerate}
 \item \(\partial^2=0\) échoue ; le complexe local n'existe alors pas.
 \item Alexander--Whitney ne commute pas avec le raffinement, la counité échoue ou
 le sommet cesse d'être group-like.
 \item \(z_\pm\) ne sont pas des cycles, leurs racines diffèrent ou
 \(\partial h_*=z_--z_+\) échoue.
 \item \(\Psi_\alpha\) n'est pas une application de complexes ou
 \(K_{\beta\alpha}=K_\beta+\Psi_\beta K_\alpha\) échoue ; dans ce cas, on a
 perdu de la retenue ou de la mémoire.
 \item Le complexe réduit n'est pas équilibré ou n'est pas acyclique sur
 \(K_3\) ; l'histoire reste hors du domaine terminal typé.
 \item Les longueurs de Bockstein ne coïncident pas avec les \(a_i\), ou
 \(\Theta_{\rm Bock}\ne\overline{\lambda^{\rm or}}\).
 \item Le raffinement modifie \(d\) ou \(\lambda^{\rm or}\) en dehors de
 blocs unitaires et de changements d'orientation enregistrés.
 \item L'une des treize coordonnées est omise : \emph{objet substitué ;
 conclusion non transférable}.
\end{enumerate}
Ces falsificateurs vérifient l'objet interne construit ; ils ne le remplacent pas
par un contrôle ultérieur.
\end{ablation}
\fi

\endgroup

\section{Incidence cellulaire et remplissage corrélatif d'une même trajectoire}

La suite \eqref{eq:pdfv2-cor-sucesion-celda} associe à chaque préfixe le complexe à deux termes
\begin{equation}
\label{eq:pdfv2-cor-cono-celda}
 C_{k,N}(\widehat x):=
 \left[A_N^{I_k}\oplus A_N^{J_k}
 \xrightarrow{\kappa_{k,N}}A_N^{I_k\times J_k}\right].
\end{equation}
Un raffinement agit sur les ports par troncature de l'extension et sur les coefficients par \(A_{N'}\twoheadrightarrow A_N\). Lorsque \(\alpha_u:I_{k'}\to I_k\) et \(\beta_u:J_{k'}\to J_k\) sont ces applications, leur action sur les générateurs est
\begin{equation}
\label{eq:pdfv2-cor-refinamiento-celda}
 (p,q)\longmapsto(p\circ\alpha_u,q\circ\beta_u),
 \qquad
 w\longmapsto w\circ(\alpha_u\times\beta_u),
\end{equation}
et vérifie
\begin{equation}
\label{eq:pdfv2-cor-refinamiento-cadena}
 \kappa_{k',N}\Gamma(u)_1=\Gamma(u)_0\kappa_{k,N}.
\end{equation}

L'incidence n'est pas fabriquée en étiquetant quatre copies d'une émission neutre. Elle est engendrée par la complétion cellulaire libre du module des deux feuillets que porte déjà le même état. Pour
\[
 V_{\widehat x}:=\mathbf F_3e_{\Sigma,\widehat x}
 \oplus\mathbf F_3e_{\Pi,\widehat x},
\]
les quatre droites internes sont
\begin{equation}
\label{eq:pdfv2-cor-cuatro-rectas-tpk}
 \mathscr I_{\widehat x}:=
 \{L_0,L_\infty,L_+,L_-\}
 =\mathbf P^1(V_{\widehat x}),\qquad
 \begin{aligned}
 L_0&=\langle e_\Sigma\rangle,&
 L_\infty&=\langle e_\Pi\rangle,\\
 L_+&=\langle e_\Sigma+e_\Pi\rangle,&
 L_-&=\langle e_\Sigma-e_\Pi\rangle.
 \end{aligned}
\end{equation}
Si \(v_L\) est respectivement \((1,0)^t,(0,1)^t,(1,1)^t,(1,-1)^t\), l'action de la direction \(L\) sur les générateurs de feuillet est le projecteur
\begin{equation}
\label{eq:pdfv2-cor-proyectores-cuatro-direcciones}
 P_L:=v_L(v_L^tv_L)^{-1}v_L^t,
\end{equation}
c'est-à-dire, sur \(\mathbf F_3\),
\[
 P_0=\begin{pmatrix}1&0\\0&0\end{pmatrix},\quad
 P_\infty=\begin{pmatrix}0&0\\0&1\end{pmatrix},\quad
 P_+=\begin{pmatrix}2&2\\2&2\end{pmatrix},\quad
 P_-=\begin{pmatrix}2&1\\1&2\end{pmatrix}.
\]
Le calcul direct donne
\begin{equation}
\label{eq:pdfv2-cor-proyectores-no-colapso}
 P_L^2=P_L,\qquad \operatorname{rk}P_L=1,\qquad
 \operatorname{Im}P_L=L,\qquad
 P_L=P_{L'}\Longrightarrow L=L'.
\end{equation}
Par conséquent, les quatre directions sont des sorties distinctes du module \(V_{\widehat x}\) ; elles ne sont pas quatre noms attribués à la même action.

\begin{definition}[Complétion cellulaire libre des feuillets TPK]
\label{def:pdfv2-cor-completacion-celular-tpk}
Soit \(\operatorname{Cell}_{\rm TPK}(\widehat x)\) le complexe libre
\[
 C_2(\widehat x)\xrightarrow{\partial_2}
 C_1(\widehat x)\xrightarrow{\partial_1}C_0(\widehat x)
\]
avec
\begin{align*}
 C_0(\widehat x)
 &:=
 \mathbf F_3[v_\star]\oplus
 \bigoplus_{L\in\mathscr I}\mathbf F_3[v_L]\oplus
 \bigoplus_{(L,\epsilon)\in\mathscr B}
       \mathbf F_3[v_{L,\epsilon}]\oplus
 \bigoplus_{E\in\mathscr E}\mathbf F_3[v_E],\\
 C_1(\widehat x)
 &:=
 \bigoplus_{L\in\mathscr I}\mathbf F_3[u_L]\oplus
 \bigoplus_{(L,\epsilon)\in\mathscr B}
       \mathbf F_3[u_{L,\epsilon}]\oplus
 \bigoplus_{E=\{L,L'\}\in\mathscr E}
 \bigl(\mathbf F_3[u_{L'\mid L}]
       \oplus\mathbf F_3[u_{L\mid L'}]\bigr),\\
 C_2(\widehat x)
 &:=
 \bigoplus_{E\in\mathscr E}\mathbf F_3[\omega_E].
\end{align*}
Son action sur les générateurs est
\begin{align}
 \partial u_L&=v_L-v_\star,
 &
 \partial u_{L'\mid L}&=v_E-v_L,
 \label{eq:pdfv2-cor-cell-frontera-aristas}\\
 \partial u_{L,\epsilon}&=v_{L,\epsilon}-v_L,
 &
 \partial\omega_{\{L,L'\}}
 &=u_L+u_{L'\mid L}-u_{L'}-u_{L\mid L'}.
 \label{eq:pdfv2-cor-cell-frontera-celda}
\end{align}
\end{definition}

Les huit générateurs \(v_{L,\epsilon}\) réalisent concrètement l'écran \(\mathscr B\). L'involution précédente s'étend au complexe par
\begin{equation}
\label{eq:pdfv2-cor-pantalla-ocho-generadores}
 \vartheta_Lv_{K,\epsilon}:=v_{\vartheta_L(K,\epsilon)},\qquad
 \vartheta_Lu_{K,\epsilon}:=u_{\vartheta_L(K,\epsilon)},
\end{equation}
et fixe les autres générateurs. La base étant libre,
\begin{equation}
\label{eq:pdfv2-cor-pantalla-ocho-no-colapso}
 |\{v_{L,\epsilon}:(L,\epsilon)\in\mathscr B\}|=8,\qquad
 \vartheta_L^2=I,\qquad
 \vartheta_L\vartheta_{L'}=\vartheta_{L'}\vartheta_L.
\end{equation}
En outre, \(\partial\vartheta_L=\vartheta_L\partial\), car les deux parcours envoient \(u_{L,\epsilon}\) sur \(v_{L,-\epsilon}-v_L\) et fixent les frontières restantes.

Les deux mots
\[
 p_{L,L'}:=u_{L'\mid L}u_L,\qquad
 p_{L',L}:=u_{L\mid L'}u_{L'}
\]
sont des chemins différents du complexe libre. Leur différence est la frontière de \(\omega_{\{L,L'\}}\), et
\begin{equation}
\label{eq:pdfv2-cor-cell-dos-cero}
 \partial^2\omega_E
 =(v_L-v_\star)+(v_E-v_L)
 -(v_{L'}-v_\star)-(v_E-v_{L'})=0.
\end{equation}
La représentation du TPK sur les feuillets est fixée par
\[
 \varrho(u_L)=P_L,\qquad
 \varrho(u_{L'\mid L})=P_{L'},\qquad
 \varrho(p_{L,L'})=P_{L'}P_L.
\]
Même si deux matrices exprimées coïncidaient dans une spécialisation, les deux chemins ne se confondraient pas : ils restent des générateurs libres distincts de \(C_1(\widehat x)\), reliés par une 2--cellule explicite.

La coordonnée complète d'incidence est le groupe des 2--cochaînes
\begin{equation}
\label{eq:pdfv2-cor-q-incidencia-total}
 Q_{\rm inc}(\widehat x):=
 C^2(\operatorname{Cell}_{\rm TPK}(\widehat x);\mathbf F_3)
 =\bigoplus_{E\in\mathscr E}\mathbf F_3\omega_E^*
 \simeq\mathbf F_3^6.
\end{equation}
L'opération interne renvoie tout \(Q_{\rm inc}(\widehat x)\), non un élément sélectionné. Chaque \(q\) possède une expansion unique \(q=\sum_Eq_E\omega_E^*\). Pour \(a\in\mathbf F_3^{\mathscr I}\), son changement de potentiel est
\begin{equation}
\label{eq:pdfv2-cor-collar-variacion}
 q\longmapsto q+Ba,\qquad (Ba)_{\{L,L'\}}=a_L+a_{L'}.
\end{equation}
Puisque chaque droite appartient à trois paires,
\[
 s(Ba)=3\sum_{L\in\mathscr I}a_L=0,
\]
et donc
\begin{equation}
\label{eq:pdfv2-cor-wc-invariante-celular}
 W_c(q+Ba)=W_c(q)\qquad(q\in Q_{\rm inc}(\widehat x)).
\end{equation}
L'unité de la complétion ne fait qu'ancrer le complexe dans l'état :
\begin{equation}
\label{eq:pdfv2-cor-cell-unidad}
 \eta^{\rm cell}_{\widehat x}:V_{\widehat x}\longrightarrow
 V_{\widehat x}\otimes C_0(\operatorname{Cell}_{\rm TPK}(\widehat x)),\qquad
 e_{\diamond,\widehat x}\longmapsto
 e_{\diamond,\widehat x}\otimes v_\star,
\end{equation}
et marque l'origine \(q=0\) ; elle n'affirme pas que l'orbite de cette origine soit tout \(Q_{\rm inc}\).

La construction est fonctorielle relativement à la troncature commune. Si \(\tau_{\ell,k}\widehat x_\ell=\widehat x_k\), nous définissons
\begin{align*}
 \tau_V&:V_{\widehat x_\ell}\longrightarrow V_{\widehat x_k},&
 \tau_Ve_{\diamond,\widehat x_\ell}&=e_{\diamond,\widehat x_k},\\
 \tau_a&:\mathbf F_3^{\mathscr I_{\widehat x_\ell}}
          \longrightarrow\mathbf F_3^{\mathscr I_{\widehat x_k}},&
 (\tau_aa)_{\tau_{\mathscr I}L}&=a_L,\\
 \tau_Q&:Q_{\rm inc}(\widehat x_\ell)
          \longrightarrow Q_{\rm inc}(\widehat x_k),&
 (\tau_Qq)_{\tau_{\mathscr E}E}&=q_E,
\end{align*}
où \(\tau_{\mathscr I}[v]=[\tau_Vv]\) et \(\tau_{\mathscr E}\{L,L'\}=\{\tau_{\mathscr I}L,
\tau_{\mathscr I}L'\}\). L'application de chaînes \(\operatorname{Cell}(\tau):C_\bullet(\widehat x_\ell)
\to C_\bullet(\widehat x_k)\) agit sur les générateurs par
\begin{equation}
\label{eq:pdfv2-cor-cell-truncamiento-generadores}
 v_\bullet(\widehat x_\ell)\mapsto v_\bullet(\widehat x_k),\quad
 u_\bullet(\widehat x_\ell)\mapsto u_\bullet(\widehat x_k),\quad
 \omega_E(\widehat x_\ell)\mapsto\omega_E(\widehat x_k).
\end{equation}
De \eqref{eq:pdfv2-cor-cell-frontera-aristas}--\eqref{eq:pdfv2-cor-cell-frontera-celda}, on obtient
\begin{equation}
\label{eq:pdfv2-cor-collar-naturalidad}
\begin{aligned}
 \operatorname{Cell}(\tau)\partial
 &=\partial\operatorname{Cell}(\tau),\\
 \tau_VP_L&=P_{\tau_{\mathscr I}L}\tau_V,\qquad
 \tau_QB=B\tau_a,\qquad W_c\circ\tau_Q=W_c,\\
 \operatorname{Cell}(\tau)\vartheta_L
 &=\vartheta_{\tau_{\mathscr I}L}\operatorname{Cell}(\tau),\\
 (\tau_V\otimes\operatorname{Cell}_0(\tau))
 \eta^{\rm cell}_{\widehat x_\ell}
 &=\eta^{\rm cell}_{\widehat x_k}\tau_V.
\end{aligned}
\end{equation}
L'identité et la composition se vérifient sur chaque générateur. Ainsi, \(\operatorname{Cell}_{\rm TPK}\) est la complétion cellulaire fonctorielle de la même fibre APP--TRIT--TPK, non un support ajouté à partir d'une expression ultérieure.

\subsection{Mesure d'incidence, courbure de retenue et naturalité adjointe}
Dans le développement suivant, on identifie, exclusivement comme abréviations de la complétion commune,
\[
 I_{\rm gau}:=\mathscr I_{\widehat x},\qquad
 \mathscr B_{\rm gau}:=\mathscr B_{\widehat x},\qquad
 Q_{\rm gau}:=Q_{\rm inc}(\widehat x).
\]
Avec la moyenne uniforme de comptage sur \(Q_{\rm gau}\), la fonction \(W_c\) de \eqref{eq:pdfv2-cor-wc} satisfait
\begin{equation}
\label{eq:pdfv2-cor-wc-momentos}
 \mathbb E_QW_c=0,\qquad
 \mathbb E_QW_c^2=\frac29,\qquad
 \mathbb E_QW_c^3=\frac{2}{27}.
\end{equation}
Ces identités sont des propriétés de la même cellule et de la même retenue, non des données d'une théorie extérieure.
\begingroup
\let\section\subsection
% Extracción quirúrgica del propietario V2 05c: líneas 99--335.
% El resto permanece en este archivo como procedencia, pero no se publica.
\iffalse
\chapter[Macrostructure de jauge interne]{Macrostructure de jauge interne engendrée par le continu discret}
\label{chap:gauint-macroestructura}

\begin{thesisbox}
La macrostructure de ce chapitre se construit uniquement à partir d'une
restriction survivante de \(\mathcal C_{\rm cont}^{\rm disc}\). Les nombres
\(4\), \(6\) et \(8\), l'écran, les quatre involutions, la mesure
projective, l'opérateur de transfert, le mode d'incidence et le terminal
cellulaire se déduisent au sein de cette restriction. Aucune évaluation ultérieure
ne fait partie du domaine ni des flèches qui suivent.
\end{thesisbox}

\section{Le germe projectif et la dérivation interne de \texorpdfstring{\(4/6/8\)}{4/6/8}}

Soit \(V_{\rm TPK}=\mathbf F_3^2\), avec la base ordonnée héritée des deux
feuillets de l'état TPK. Nous définissons l'ensemble des directions
\begin{equation}\label{eq:gauint-cuatro-direcciones}
 I_{\rm gau}:=\mathbf P^1(\mathbf F_3)
 =\{L_\infty,L_0,L_1,L_{-1}\},
\end{equation}
où
\[
 L_\infty=[1:0],\qquad L_0=[0:1],\qquad
 L_1=[1:1],\qquad L_{-1}=[1:-1].
\]
Par conséquent,
\begin{equation}\label{eq:gauint-cardinal-cuatro}
 |I_{\rm gau}|=\frac{3^2-1}{3-1}=4.
\end{equation}
On n'introduit pas un ensemble de quatre indices : il apparaît comme le quotient des
huit vecteurs non nuls par l'action de \(\mathbf F_3^\times\).

Les incidences non orientées sont
\begin{equation}\label{eq:gauint-seis-incidencias}
 E_{\rm gau}:=\binom{I_{\rm gau}}2,
 \qquad |E_{\rm gau}|=\binom42=6,
 \qquad Q_{\rm gau}:=\mathbf F_3^{E_{\rm gau}}.
\end{equation}
L'écran orienté double chaque direction, et non chaque incidence :
\begin{equation}\label{eq:gauint-ocho-caras}
 \mathscr B_{\rm gau}:=I_{\rm gau}\times\{+,-\},
 \qquad |\mathscr B_{\rm gau}|=2|I_{\rm gau}|=8.
\end{equation}
Ainsi, \(4\), \(6\) et \(8\) proviennent respectivement de directions, de couples de
directions et de faces orientées d'un même germe.

Pour chaque \(L\in I_{\rm gau}\), soit \(\vartheta_L\) l'involution de
l'écran qui échange \((L,+)\) et \((L,-)\) et laisse fixes les six autres
faces. Alors
\begin{equation}\label{eq:gauint-cuatro-involuciones}
 \vartheta_L^2=I,
 \qquad
 \vartheta_L\vartheta_{L'}=\vartheta_{L'}\vartheta_L
 \quad(L\ne L').
\end{equation}
Les quatre involutions ont été obtenues avant de définir une mesure ou un opérateur.

\section{Incidence, mode interne et témoin non constant}

L'application d'incidence est
\begin{equation}\label{eq:gauint-incidencia-b}
 B:\mathbf F_3^{I_{\rm gau}}\longrightarrow Q_{\rm gau},
 \qquad
 (Ba)_{\{L,L'\}}:=a_L+a_{L'}.
\end{equation}
Si
\[
 s(q):=\sum_{e\in E_{\rm gau}}q_e\in\mathbf F_3,
\]
chaque direction apparaît exactement trois fois dans la somme des six
incidences, et donc
\begin{equation}\label{eq:gauint-incidencia-invariante}
 s(Ba)=3\sum_{L\in I_{\rm gau}}a_L=0,
 \qquad
 s(q+Ba)=s(q).
\end{equation}

Sur \(Q_{\rm gau}\), on utilise exclusivement la moyenne de comptage
normalisée. Le mode d'incidence est la fonction rationnelle
\begin{equation}\label{eq:gauint-wc-definicion}
 W_c(q):=\mathbf 1_{\{s(q)=0\}}-\frac13.
\end{equation}
Comme \(s:Q_{\rm gau}\to\mathbf F_3\) est surjective, chaque fibre a
\(3^5\) éléments. On en obtient, par comptage exact,
\begin{equation}\label{eq:gauint-wc-momentos}
 \mathbb E_{Q}W_c=0,
 \qquad
 \mathbb E_{Q}W_c^2=\frac29,
 \qquad
 \mathbb E_{Q}W_c^3=\frac{2}{27}.
\end{equation}
En outre, \eqref{eq:gauint-incidencia-invariante} donne
\begin{equation}\label{eq:gauint-wc-invariancia}
 W_c(q+Ba)=W_c(q).
\end{equation}
Donc \(W_c\) est non nul, centré et invariant par l'incidence
engendrée en interne.

\fi
\section{Incidence, courbure de retenue et domaine commun de naturalité}

Pour \(x\in\mathcal C_{\rm cont}^{\rm disc}\), notons
\(Y_m(x)\) l'ensemble fini des préfixes de jauge de profondeur \(m\) qui
conservent simultanément
\begin{equation*}
\begin{gathered}
 \text{fibre, relations d’extension, nombres, orientation, retenue, mémoire,}\\
 \text{frontière, route, multisection, survie, terminal et lecteur}.
\end{gathered}
\end{equation*}
Il existe une application de collier
\begin{equation}\label{eq:gauint-collar}
 q_m:Y_m(x)\longrightarrow Q_{\rm gau}
\end{equation}
qui enregistre les six incidences sans effacer l'histoire qui les produit.
Pour \(\ell\ge m\), la troncature s'écrit
\[
 \tau_{\ell,m}:Y_\ell(x)\longrightarrow Y_m(x).
\]

Pour \(y,z\in Y_m(x)\), soit
\begin{equation}\label{eq:gauint-todas-extensiones}
 \operatorname{Ext}^{\rm adm}_{m}(y,z)
\end{equation}
l'ensemble de \emph{toutes} les extensions TPK admissibles de \(y\) à \(z\).
Une extension appartient à cet ensemble seulement si elle transporte la fibre entière,
les relations \(\operatorname{Ext}_\Gamma\), la retenue native et opératorielle,
l'orientation, la mémoire accumulée, la frontière, le chemin et la condition de
survie. Nous définissons
\begin{equation}\label{eq:gauint-conteo-kernel}
 a_m(y,z):=\#\operatorname{Ext}^{\rm adm}_m(y,z),
 \qquad
 d_m(y):=\sum_{z\in Y_m(x)}a_m(y,z),
\end{equation}
et, lorsque \(d_m(y)>0\),
\begin{equation}\label{eq:gauint-kernel-normalizado}
 k_m(y,z):=\frac{a_m(y,z)}{d_m(y)}.
\end{equation}
Ainsi, \(k_m\) est un comptage normalisé d'extensions complètes, non une
fonction introduite après la projection de l'état.

La restriction \(\mathcal C_{\rm cont}^{\rm gau}\) est formée des
histoires survivantes \(x\) pour lesquelles les conditions
suivantes sont satisfaites à toute profondeur.
\begin{enumerate}
 \item Les ensembles \(Y_m(x)\) sont non vides, \(d_m(y)>0\), et les
 troncatures \(\tau_{\ell,m}\) sont surjectives.
 \item Il existe une multisection non vide de poids rationnels strictement
 positifs \(\mu_m:Y_m(x)\to\mathbf Q_{>0}\) telle que
 \begin{equation}\label{eq:gauint-medida-proyectiva}
  \sum_{y\in Y_m}\mu_m(y)=1,
  \qquad
  \mu_m(y)=\sum_{\tau_{\ell,m}\widetilde y=y}\mu_\ell(\widetilde y).
 \end{equation}
 \item Chaque \(\mu_m\) est stationnaire :
 \begin{equation}\label{eq:gauint-estacionariedad}
  \sum_{y\in Y_m}\mu_m(y)k_m(y,z)=\mu_m(z).
 \end{equation}
 \item La marginale de \((q_m)_*\mu_m\) est le comptage uniforme de
 \(Q_{\rm gau}\). Cette condition rend littéralement applicables les calculs
 de \eqref{eq:pdfv2-cor-wc-momentos} au relèvement \(W_c\circ q_m\).
 \item La lumpabilité directe est satisfaite
 \begin{equation}\label{eq:gauint-lump-directa}
  \sum_{\tau_{\ell,m}\widetilde z=z}
   k_\ell(\widetilde y,\widetilde z)
  =k_m(\tau_{\ell,m}\widetilde y,z)
 \end{equation}
 pour tout \(\widetilde y\in Y_\ell\) et \(z\in Y_m\).
 \item Pour le noyau adjoint
 \begin{equation}\label{eq:gauint-kernel-adjunto}
  k_m^*(z,y):=\frac{\mu_m(y)k_m(y,z)}{\mu_m(z)},
 \end{equation}
 la lumpabilité adjointe est satisfaite
 \begin{equation}\label{eq:gauint-lump-adjunta}
  \sum_{\tau_{\ell,m}\widetilde y=y}
   k_\ell^*(\widetilde z,\widetilde y)
  =k_m^*(\tau_{\ell,m}\widetilde z,y).
 \end{equation}
 \item Les catégories de préfixes possèdent la racine TPK comme objet initial,
 de sorte que leur nerf possède la contraction terminale explicite de
 \eqref{eq:gauint-sdr} ci-dessous.
\end{enumerate}
Ces conditions définissent le domaine typé ; en dehors de celui-ci, on ne publie pas une
mesure ou un opérateur naturel par simple déclaration.

\section{Mesure projective, noyau et naturalité directe et adjointe}

Une branche de la multisection \(\{\mu_m\}_m\) étant fixée, soit
\begin{equation}\label{eq:gauint-hm}
 H_m:=\operatorname{Fun}(Y_m,\mathbf Q),
 \qquad
 \langle f,g\rangle_m:=\sum_{y\in Y_m}\mu_m(y)f(y)g(y).
\end{equation}
L'opérateur engendré par toutes les extensions est
\begin{equation}\label{eq:gauint-km}
 (K_mf)(y):=\sum_{z\in Y_m}k_m(y,z)f(z).
\end{equation}
Les équations de Markov et de stationnarité donnent
\begin{equation}\label{eq:gauint-markov-estacionario}
 K_m\mathbf1=\mathbf1,
 \qquad
 K_m^*\mathbf1=\mathbf1,
 \qquad
 \|K_mf\|_m\le\|f\|_m.
\end{equation}

Le relèvement et l'espérance de troncature sont
\begin{equation}\label{eq:gauint-lift-esperanza}
 J_{\ell,m}f:=f\circ\tau_{\ell,m},
 \qquad
 E_{m,\ell}:=J_{\ell,m}^*.
\end{equation}
La projectivité de \(\mu\) implique
\begin{equation}\label{eq:gauint-ej-identidad}
 E_{m,\ell}J_{\ell,m}=I_{H_m}.
\end{equation}
Les deux conditions de lumpabilité établissent, respectivement,
\begin{equation}\label{eq:gauint-naturalidad-directa-adjunta}
 K_\ell J_{\ell,m}=J_{\ell,m}K_m,
 \qquad
 K_\ell^*J_{\ell,m}=J_{\ell,m}K_m^*.
\end{equation}
La seconde égalité ne se déduit pas de la première sans la condition adjointe :
les deux sont enregistrées dans le domaine.

On définit l'opérateur interne positif
\begin{equation}\label{eq:gauint-transferencia-t}
 T_m:=K_m^*K_m.
\end{equation}
Alors
\begin{equation}\label{eq:gauint-t-positivo-natural}
 \langle f,T_mf\rangle_m=\|K_mf\|_m^2\ge0,
 \qquad
 T_\ell J_{\ell,m}=J_{\ell,m}T_m.
\end{equation}

\section{Loi commune de \(8\) faces et \(4\) carrés de Fubini}

Sur \(Y_m^{\mathscr B_{\rm gau}}\), on définit la masse rationnelle
\begin{equation}\label{eq:gauint-ley-ocho-caras}
 \Omega_m^{(8)}(\mathbf y)
 :=\sum_{z\in Y_m}\mu_m(z)
 \prod_{(L,\varepsilon)\in\mathscr B_{\rm gau}}
 k_m(z,y_{L,\varepsilon}).
\end{equation}
C'est une probabilité parce que chaque facteur est markovien. Toutes les faces proviennent
du même état latent \(z\) et du même noyau ; ce ne sont pas huit mesures
indépendantes ajoutées.

Pour des fonctions \(f_{L,+},f_{L,-}\in H_m\), Fubini fini donne
\begin{equation}\label{eq:gauint-fubini-ocho}
\begin{aligned}
 &\sum_{\mathbf y}\Omega_m^{(8)}(\mathbf y)
  \prod_{L\in I_{\rm gau}}
  f_{L,+}(y_{L,+})f_{L,-}(y_{L,-})\\
 &\qquad=
 \sum_{z\in Y_m}\mu_m(z)
 \prod_{L\in I_{\rm gau}}
 (K_mf_{L,+})(z)(K_mf_{L,-})(z).
\end{aligned}
\end{equation}
En prenant \(f_{L,+}=f_{L,-}=f_L\), les quatre
carrés apparaissent simultanément
\begin{equation}\label{eq:gauint-cuatro-cuadrados}
 \sum_{z\in Y_m}\mu_m(z)
 \prod_{L\in I_{\rm gau}}(K_mf_L(z))^2.
\end{equation}
La marginale de tout couple opposé satisfait
\begin{equation}\label{eq:gauint-marginal-t}
 \sum_{y_+,y_-}\Omega_{m,L}^{(2)}(y_+,y_-)
 f(y_+)g(y_-)
 =\langle K_mf,K_mg\rangle_m
 =\langle f,T_mg\rangle_m.
\end{equation}
La commutativité du produit dans \eqref{eq:gauint-ley-ocho-caras} prouve
\begin{equation}\label{eq:gauint-involuciones-ley}
 (\vartheta_L)_*\Omega_m^{(8)}=\Omega_m^{(8)}
 \qquad(L\in I_{\rm gau}).
\end{equation}

\section{Opérateur interne exact et témoin propre}

La projection sur le secteur constant et son complément sont
\begin{equation}\label{eq:gauint-pq}
 P_{\Omega,m}f:=\langle\mathbf1,f\rangle_m\mathbf1,
 \qquad
 Q_m:=I-P_{\Omega,m}.
\end{equation}
L'opérateur cellulaire interne est défini par
\begin{equation}\label{eq:gauint-d-hmt}
 D_m^{\rm HMT}:=3Q_m.
\end{equation}
Le coefficient \(3\) est le cardinal du corps des résidus utilisé dans
\eqref{eq:pdfv2-cor-proyectores-cuatro-direcciones} ; il ne représente pas une
échelle ajoutée.
Pour le relèvement
\[
 W_{c,m}:=W_c\circ q_m
\]
on a, par \eqref{eq:pdfv2-cor-wc-momentos},
\begin{equation}\label{eq:gauint-w-eigen}
 P_{\Omega,m}W_{c,m}=0,
 \qquad
 D_m^{\rm HMT}W_{c,m}=3W_{c,m},
 \qquad
 \|W_{c,m}\|_m^2=\frac29.
\end{equation}
Ainsi, le secteur complémentaire possède un témoin non nul construit par
l'incidence des six arêtes.

\section{Terminal cellulaire et rétracte fort}

Soit \(C_*(N\mathsf Y_m;A_n)\) le complexe normalisé du nerf de la
catégorie des préfixes, et soit \(y_*\) sa racine initiale. L'inclusion et
l'augmentation sont
\[
 \iota_m:A_n[0]\longrightarrow C_*(N\mathsf Y_m;A_n),
 \quad \iota_m(1)=[y_*],
 \qquad
 \epsilon_m:C_*(N\mathsf Y_m;A_n)\longrightarrow A_n[0].
\]
L'objet initial fournit la dégénérescence supplémentaire
\begin{equation}\label{eq:gauint-homotopia-celular}
 H_m[y_0\to\cdots\to y_q]
 :=[y_*\to y_0\to\cdots\to y_q],
\end{equation}
avec la convention normalisée lorsqu'une dégénérescence apparaît. On vérifie
\begin{equation}\label{eq:gauint-sdr}
 \epsilon_m\iota_m=I,
 \qquad
 \partial H_m+H_m\partial=I-\iota_m\epsilon_m,
 \qquad
 H_m^2=H_m\iota_m=\epsilon_mH_m=0.
\end{equation}
Ce terminal cellulaire est distinct de la projection constante
\(P_{\Omega,m}\) : le premier contracte le nerf et la seconde sépare les
secteurs de l'espace des fonctions. Les deux sont conservés.

\iffalse
\section{Les treize coordonnées de la restriction de jauge}

Pour chaque \((m,n)\), on définit
\begin{equation}\label{eq:gauint-d-trece}
 D_{{\rm gau},m,n}(x):=
 (\mathcal F,\operatorname{Ext}_\Gamma,\operatorname{Rel},
 \mathcal N,\operatorname{or},\operatorname{Carr},\operatorname{Mem},
 \partial,\gamma,\operatorname{Msect},\operatorname{Surv},
 \operatorname{Term},\mathcal L)_{{\rm gau},m,n}(x),
\end{equation}
avec le contenu suivant.
\begin{enumerate}
 \item \(\mathcal F=(V_{\rm TPK},I_{\rm gau},E_{\rm gau},
 Q_{\rm gau},\mathscr B_{\rm gau},Y_m,H_m)\).
 \item \(\operatorname{Ext}_\Gamma\) contient tous les ensembles
 \(\operatorname{Ext}^{\rm adm}_m(y,z)\), les troncatures, les relèvements,
 les espérances et leurs compositions.
 \item \(\operatorname{Rel}\) contient \(B\), \(sB=0\), les quatre
 involutions, Markov, la stationnarité et les deux lumpabilités.
 \item \(\mathcal N=(3,4,6,8,m,n,|Y_m|,a_m,d_m)\) ; ces nombres ne
 remplacent aucune des autres coordonnées.
 \item \(\operatorname{or}\) conserve les faces \(+/-\), les quatre
 inversions \(\vartheta_L\) et le renversement de chaque chemin.
 \item \(\operatorname{Carr}\) contient la retenue complète de chaque
 extension comptée dans \(a_m(y,z)\), y compris la retenue déjà propagée.
 \item \(\operatorname{Mem}\) conserve le préfixe complet, l'état
 commun \(z\) de \(\Omega_m^{(8)}\) et la mémoire de ses huit extensions.
 \item \(\partial=(q_m,\partial_{N\mathsf Y_m},Q_m,H_m)\) conserve
 la frontière de collier, la frontière cellulaire et le complément.
 \item \(\gamma\) est le simplexe orienté du nerf qui enregistre le chemin
 et non seulement ses extrémités.
 \item \(\operatorname{Msect}\) est la famille de toutes les mesures
 stationnaires projectives compatibles et de toutes les extensions qui
 interviennent dans les comptages ; aucune branche rétrospective n'est sélectionnée.
 \item \(\operatorname{Surv}\) enregistre la non-vacuité, la profondeur
 arbitraire, la projectivité, les deux lumpabilités et la persistance de la
 marginale uniforme de collier.
 \item \(\operatorname{Term}=(\iota_m,\epsilon_m,H_m;
 P_{\Omega,m},Q_m)\) conserve séparément les deux terminaux.
 \item \(\mathcal L=(\Omega_m^{(8)},K_m,T_m,D_m^{\rm HMT},W_{c,m})\)
 est le lecteur interne postérieur à la construction de l'histoire.
\end{enumerate}

Pour \(\ell\ge m\) et \(n'\ge n\), toutes les coordonnées satisfont
\begin{equation}\label{eq:gauint-naturalidad-trece}
 u^{\rm gau}_{(\ell,n'),(m,n)}D_{{\rm gau},\ell,n'}(x')
 =D_{{\rm gau},m,n}(\tau_{\ell,m}x').
\end{equation}

\section{Assembleur, publicateur et chaîne non ablative}

Soit
\begin{equation}\label{eq:gauint-graph-d}
 \mathcal G_{\rm gau}:=\operatorname{Graph}(D_{\rm gau}),
 \qquad
 \operatorname{Res}_{\rm gau}(x):=(x,D_{\rm gau}x).
\end{equation}
La multisection dépendante \(\chi_{\rm gau}(x)\) contient toutes les branches
compatibles de collier, de mesure et d'extension. Nous définissons
\begin{equation}\label{eq:gauint-x-dependiente}
 X_{\chi,{\rm gau}}
 :=\{(\chi_{\rm gau}(x),x,D_{\rm gau}x):
 x\in\mathcal C_{\rm cont}^{\rm gau}\},
\end{equation}
\begin{equation}\label{eq:gauint-prx}
 \operatorname{pr}_{X,{\rm gau}}(x,D_{\rm gau}x)
 :=(\chi_{\rm gau}(x),x,D_{\rm gau}x).
\end{equation}

L'assembleur brut
\begin{equation}\label{eq:gauint-ensamblador-crudo}
 a_{\rm gau}:X_{\chi,{\rm gau}}\longrightarrow
 \mathscr M_{\rm gau}^{\rm amb}
\end{equation}
produit la famille complète
\[
 \bigl(I,E,\mathscr B,B,W_c;
 Y_m,\mu_m,k_m,J_{\ell,m},E_{m,\ell};
 K_m,T_m,\Omega_m^{(8)},P_{\Omega,m},Q_m,D_m^{\rm HMT};
 \iota_m,\epsilon_m,H_m\bigr)_{m,n}.
\]
Le macro-objet conserve son entrée par
\begin{equation}\label{eq:gauint-graph-a}
 \mathfrak M_{\rm gau}:=\operatorname{Graph}(a_{\rm gau}),
 \qquad
 \mathcal A_{\rm gau}(z):=(z,a_{\rm gau}z).
\end{equation}

Le publicateur brut
\begin{equation}\label{eq:gauint-publicador-crudo}
 p_{\rm gau}:\mathfrak M_{\rm gau}\longrightarrow
 \mathscr Y_{\rm gau}
\end{equation}
publie conjointement les identités
\[
 \begin{gathered}
 |I|=4,quad |E|=6,quad |\mathscr B|=8,quad sB=0,
 \quad W_c(q+Ba)=W_c(q),\\
 K_\ell J=JK_m,quad K_\ell^*J=JK_m^*,
 \quad T_m=K_m^*K_m,\\
 (\vartheta_L)_*\Omega_m^{(8)}=\Omega_m^{(8)},
 \quad \text{les quatre factorisations quadratiques de Fubini},\\
 D_m^{\rm HMT}=3Q_m,quad D_m^{\rm HMT}W_{c,m}=3W_{c,m},
 \quad \partial H_m+H_m\partial=I-\iota_m\epsilon_m.
 \end{gathered}
\]
Enfin,
\begin{equation}\label{eq:gauint-graph-p}
 \mathcal C_{\rm gau}^{\rm HMT}:=\operatorname{Graph}(p_{\rm gau}),
 \qquad
 \mathcal P_{\rm gau}(M):=(M,p_{\rm gau}M).
\end{equation}
La chaîne complète est
\begin{equation}\label{eq:gauint-cadena-completa}
 \mathcal C_{\rm cont}^{\rm disc}\supset
 \mathcal C_{\rm cont}^{\rm gau}
 \xrightarrow{\operatorname{Res}_{\rm gau}}\mathcal G_{\rm gau}
 \xrightarrow{\operatorname{pr}_{X,{\rm gau}}}X_{\chi,{\rm gau}}
 \xrightarrow{\mathcal A_{\rm gau}}\mathfrak M_{\rm gau}
 \xrightarrow{\mathcal P_{\rm gau}}\mathcal C_{\rm gau}^{\rm HMT}.
\end{equation}
Chaque flèche possède un inverse sur son image donné par une projection de
coordonnées. Aucune sortie ne remplace son histoire génératrice.

\section{Théorème de réalisation de jauge interne}

\begin{theorem}[Réalisation interne complète et naturelle]
\label{thm:gauint-realizacion}
Pour tout \(x\in\mathcal C_{\rm cont}^{\rm gau}\), la chaîne
\eqref{eq:gauint-cadena-completa} construit, à partir de la structure
discrète du continu, une macrostructure naturelle avec quatre directions,
six incidences, huit faces, quatre involutions, une mesure projective, un
noyau de comptage complet, son transfert positif, une loi commune de huit
faces, un témoin d'incidence non constant et un terminal cellulaire. Les
treize coordonnées de \eqref{eq:gauint-d-trece} restent récupérables en
sortie.
\end{theorem}

\begin{proof}
Les équations \eqref{eq:gauint-cuatro-direcciones}--
\eqref{eq:gauint-ocho-caras} déduisent \(4/6/8\). La somme des incidences
\eqref{eq:gauint-incidencia-invariante} prouve l'invariance de \(W_c\), et
le comptage uniforme prouve \eqref{eq:gauint-wc-momentos}. Les extensions
complètes produisent \(k_m\) par \eqref{eq:gauint-kernel-normalizado} ; les
conditions explicites du domaine prouvent Markov, la stationnarité et les deux
égalités de naturalité. Donc \(T_m=K_m^*K_m\) est positif et naturel.
Fubini appliqué une seule fois à \eqref{eq:gauint-ley-ocho-caras} donne les quatre
formes quadratiques sans changer de mesure. La marginale uniforme centre
\(W_{c,m}\), d'où \eqref{eq:gauint-w-eigen}. L'objet initial du nerf
donne la dégénérescence supplémentaire et l'identité SDR \eqref{eq:gauint-sdr}. Enfin,
les trois constructions par graphe retiennent leurs entrées comme premières
coordonnées, donc la composition n'oublie aucune des treize entrées.
\end{proof}

\section{Provenance et falsificateurs}

\paragraph{Provenance.}
Le germe TPK, l'écran orienté, les extensions survivantes, la
mesure commune et la contraction terminale ont le statut
\texttt{ARCHITECTURE\_AUTORALE\_PRÉEXISTANTE}. Le transfert par comptage,
le mode d'incidence et le secteur exact \(3Q\) sont présentés comme
\texttt{RÉSULTAT\_RÉCUPÉRÉ}. La dérivation explicite
\(\mathbf P^1(\mathbf F_3)\to4/6/8\), les deux conditions de lumpabilité et
la mise en paquet non ablative par graphes ont le statut
\texttt{FORMALISATION\_NOUVELLE} ; elles formalisent l'architecture antérieure sans
revendiquer une origine conceptuelle nouvelle.

\begin{ablation}[Falsificateurs internes de la branche de jauge]
\label{abl:gauint-falsadores}
La construction est falsifiée au niveau indiqué si l'un quelconque des
faits suivants se produit.
\begin{enumerate}
 \item On remplace \(I_{\rm gau}\), \(E_{\rm gau}\) ou
 \(\mathscr B_{\rm gau}\) par une liste sans les applications qui les déduisent.
 \item \(sB=0\) échoue ou \(W_c\) cesse d'être invariant et non constant.
 \item Le comptage \(a_m(y,z)\) omet la retenue, l'orientation, la mémoire, la frontière ou
 le chemin : dans ce cas, on a compté un autre objet.
 \item La projectivité, la stationnarité ou l'une des deux
 lumpabilités échoue ; on ne peut alors publier la naturalité correspondante.
 \item Les huit faces ne proviennent pas d'une unique \(\mu_m\) et d'un unique
 \(k_m\), ou l'une des \(\vartheta_L\) ne préserve pas \(\Omega_m^{(8)}\).
 \item \(P_{\Omega,m}W_{c,m}\ne0\) o
 \(D_m^{\rm HMT}W_{c,m}\ne3W_{c,m}\).
 \item L'identité SDR \(\partial H+H\partial=I-\iota\epsilon\) échoue.
 \item L'une des treize coordonnées est omise : \emph{objet substitué ;
 conclusion non transférable}.
\end{enumerate}
Ces falsificateurs contrôlent l'appartenance et la naturalité ; ils ne dégradent pas une
histoire qui satisfait toutes les identités.
\end{ablation}
\fi

\endgroup

La même règle définit son action sur une émission élémentaire \(\alpha:x\to y\). Puisque \(U_\alpha\) transporte la paire \emph{ordonnée} de feuillets, son application sur les étiquettes est \(\overline U_\alpha e_{\Sigma,x}=e_{\Sigma,y}\) et \(\overline U_\alpha e_{\Pi,x}=e_{\Pi,y}\). Sa projectivisation \(\alpha_{\mathscr I}[v]:=[\overline U_\alpha v]\) agit sur \(\mathbf P^1(V_x)\) et, par conséquent,
\begin{equation}
\label{eq:pdfv2-cor-cell-accion-emision}
 \begin{gathered}
 \operatorname{Cell}_{\rm TPK}(\alpha)v_\star(x)=v_\star(y),\qquad
 \operatorname{Cell}_{\rm TPK}(\alpha)v_L(x)
   =v_{\alpha_{\mathscr I}L}(y),\\
 \operatorname{Cell}_{\rm TPK}(\alpha)v_{\{L,L'\}}(x)
   =v_{\{\alpha_{\mathscr I}L,\alpha_{\mathscr I}L'\}}(y),\qquad
 \operatorname{Cell}_{\rm TPK}(\alpha)v_{L,\epsilon}(x)
   =v_{\alpha_{\mathscr I}L,\epsilon}(y),\\
 \operatorname{Cell}_{\rm TPK}(\alpha)u_L(x)
   =u_{\alpha_{\mathscr I}L}(y),\qquad
 \operatorname{Cell}_{\rm TPK}(\alpha)u_{L,\epsilon}(x)
   =u_{\alpha_{\mathscr I}L,\epsilon}(y),\\
 \operatorname{Cell}_{\rm TPK}(\alpha)u_{L'\mid L}(x)
   =u_{\alpha_{\mathscr I}L'\mid\alpha_{\mathscr I}L}(y),\qquad
 \operatorname{Cell}_{\rm TPK}(\alpha)\omega_{\{L,L'\}}(x)
   =\omega_{\{\alpha_{\mathscr I}L,
                 \alpha_{\mathscr I}L'\}}(y).
 \end{gathered}
\end{equation}
Dans la fibre des feuillets et sur son écran, on a en outre
\begin{equation}
\label{eq:pdfv2-cor-cell-entrelaza-proyectores-pantalla}
 \overline U_\alpha P_L=P_{\alpha_{\mathscr I}L}\overline U_\alpha,
 \qquad
 \operatorname{Cell}_{\rm TPK}(\alpha)\vartheta_L
 =\vartheta_{\alpha_{\mathscr I}L}
  \operatorname{Cell}_{\rm TPK}(\alpha).
\end{equation}
Les formules de frontière donnent
\begin{equation}
\label{eq:pdfv2-cor-cell-accion-emision-natural}
 \partial\operatorname{Cell}_{\rm TPK}(\alpha)
 =\operatorname{Cell}_{\rm TPK}(\alpha)\partial,\qquad
 \operatorname{Cell}_{\rm TPK}(\beta\alpha)
 =\operatorname{Cell}_{\rm TPK}(\beta)
  \operatorname{Cell}_{\rm TPK}(\alpha).
\end{equation}
Ainsi, l'application de complexes utilisée dans le registre terminal est l'action libre induite par la même émission, non un symbole ajouté.

\subsection{Module des parcours, totalisation et cône du retour nonadique}
Soit \(J_d(x)\) la catégorie complète des préfixes de profondeur \(d\) de l'arbre commun, \(\operatorname{Term}_d(x)\) son ensemble terminal et \(\Gamma_{d,N}(\nu):=C_{d,N}(\nu)\) le complexe cellulaire de \eqref{eq:pdfv2-cor-cono-celda}. Le module, la totalisation et le cône qui suivent conservent simultanément toutes les continuations ; ils ne sélectionnent pas une histoire et ne séparent pas une construction cohomologique du continu commun.
\begingroup
\let\section\subsection
% Extracción quirúrgica del propietario V2 05d: líneas 286--512.
% El resto permanece en este archivo como procedencia, pero no se publica.
\iffalse
% Macroestructura de coherencia estrictamente interna del continuo discreto.
% No usa geometrías, clases ni conclusiones clásicas.

\chapter{Macrostructure interne de cellules, cohérence et complétude}
\label{chap:pdfv2-coh-interna}

\begin{statusbox}
L'unique entrée de ce chapitre est la structure discrète du continu. La
branche \(\mathrm{coh}\) se construit sur les anneaux
\(A_N=\mathbb Z/3^N\mathbb Z\), les cellules APP, leurs histoires et leurs
terminaux. Sa chaîne non ablative est
\[
 \mathcal C_{\rm cont}^{\rm disc}
 \supset\mathcal C_{\rm cont}^{\rm coh}
 \xrightarrow{\operatorname{Res}_{\rm coh}}
 \mathcal G_{\rm coh}
 \xrightarrow{\operatorname{pr}_{X,{\rm coh}}}
 X_{\chi,{\rm coh}}
 \xrightarrow{\mathcal A_{\rm coh}}
 \mathfrak M_{\rm coh}
 \xrightarrow{\mathcal P_{\rm coh}}
 \mathcal C_{\rm coh}^{\rm HMT}.
\]
Aucun objet classique n'apparaît en entrée ni en sortie. Chaque image
conserve littéralement son histoire génératrice.
\end{statusbox}

\section{Domaine typé et séparation des indices}
\label{sec:pdfv2-coh-domain}

Pour ne pas identifier profondeur et précision, on utilise l'ensemble dirigé
\begin{equation}\label{eq:pdfv2-coh-lambda}
 \Lambda:=\mathbb N_{\rm prof}\times\mathbb N_{\rm prec},
 \qquad
 \lambda=(d,N),
 \qquad
 A_N:=\mathbb Z/3^N\mathbb Z.
\end{equation}
Si \(\lambda\leq\lambda'=(d',N')\), la transition
\begin{equation}\label{eq:pdfv2-coh-rho-lambda}
 \rho_{\lambda',\lambda}:
 \mathcal C_{{\rm cont},\lambda'}^{\rm disc}
 \longrightarrow
 \mathcal C_{{\rm cont},\lambda}^{\rm disc}
\end{equation}
tronque l'histoire après la profondeur \(d\) et réduit les coefficients par
\(A_{N'}\twoheadrightarrow A_N\), en conservant feuillet, orientation, résidu,
quotient, retenue, mémoire, frontière, chemin, survie et terminal.

La sous-structure \(\mathcal C_{\rm cont}^{\rm coh}\) se compose des histoires
\(x\in\mathcal C_{\rm cont}^{\rm disc}\) qui satisfont, à tout niveau :
\begin{enumerate}
 \item le préfixe de \(x\) porte les cellules orientées et basées définies dans
       la section suivante ;
 \item chaque raffinement induit une application du complexe de la cellule ;
 \item troncature, réduction des coefficients, orientation et extension
       nonadique commutent ;
 \item les multisections d'extension survivante sont non vides à toute
       profondeur ;
 \item le retour nonadique conserve la différence terminale entre l'histoire
       et sa continuation.
\end{enumerate}
Ces conditions définissent le domaine de la branche. Elles ne sont pas attribuées à une
histoire arbitraire sans être vérifiées.

\section{Cellule APP sur \texorpdfstring{\(A_N\)}{AN}}
\label{sec:pdfv2-coh-cell}

Une cellule orientée est
\[
 \nu=(h_\nu,I_\nu,J_\nu,\epsilon_\nu),
\]
où \(h_\nu\) est une histoire finie, \(I_\nu,J_\nu\) sont des ensembles
ordonnés non vides de ports,
\[
 a_\nu:=|I_\nu|\geq1,
 \qquad
 b_\nu:=|J_\nu|\geq1,
\]
et \(\epsilon_\nu\in\{+1,-1\}\) est son orientation.

On définit
\begin{equation}\label{eq:pdfv2-coh-kappa}
 \begin{aligned}
 \kappa_{\nu,N}:
 A_N^{I_\nu}\oplus A_N^{J_\nu}
 &\longrightarrow A_N^{I_\nu\times J_\nu},\\
 \kappa_{\nu,N}(u,v)_{ij}&:=u_i-v_j.
 \end{aligned}
\end{equation}
Les ports de base \(i_0\in I_\nu\), \(j_0\in J_\nu\) étant choisis, soit
\begin{equation}\label{eq:pdfv2-coh-delta}
 \begin{aligned}
 \delta_{\nu,N}:A_N^{I_\nu\times J_\nu}
 &\longrightarrow
 A_N^{(I_\nu\setminus i_0)\times(J_\nu\setminus j_0)},\\
 \delta_{\nu,N}(w)_{ij}
 &:={}
 w_{ij}-w_{i j_0}-w_{i_0j}+w_{i_0j_0}.
 \end{aligned}
\end{equation}

\begin{proposition}[Suite exacte de la cellule]
\label{prop:pdfv2-coh-cell-exact}
Pour toute cellule \(\nu\) et toute précision \(N\), la suite
\begin{equation}\label{eq:pdfv2-coh-exact-sequence}
 0\longrightarrow A_N
 \xrightarrow{\ \Delta\ }
 A_N^{I_\nu}\oplus A_N^{J_\nu}
 \xrightarrow{\ \kappa_{\nu,N}\ }
 A_N^{I_\nu\times J_\nu}
 \xrightarrow{\ \delta_{\nu,N}\ }
 A_N^{(a_\nu-1)(b_\nu-1)}
 \longrightarrow0,
\end{equation}
où
\[
 \Delta(t)=\bigl((t)_{i\in I_\nu},(t)_{j\in J_\nu}\bigr),
\]
est exacte et scindée sur \(A_N\).
\end{proposition}

\begin{proof}
Si \(\kappa_{\nu,N}(u,v)=0\), alors \(u_i=v_j\) pour tout \(i,j\),
de sorte que \((u,v)\) est diagonale. Ainsi,
\(\ker\kappa_{\nu,N}=\operatorname{im}\Delta\). La substitution directe donne
\(\delta_{\nu,N}\kappa_{\nu,N}=0\).

Si \(\delta_{\nu,N}(w)=0\), posons
\[
 u_i=w_{i j_0},
 \qquad
 v_j=w_{i_0j_0}-w_{i_0j}.
\]
Alors
\[
 u_i-v_j=w_{i j_0}+w_{i_0j}-w_{i_0j_0}=w_{ij},
\]
de sorte que \(\ker\delta_{\nu,N}=\operatorname{im}\kappa_{\nu,N}\). Pour
prouver la surjectivité de \(\delta_{\nu,N}\), on annule la ligne et la
colonne de base et l'on place arbitrairement les coordonnées souhaitées aux
positions restantes. Ces mêmes formules fournissent des scindements.
\end{proof}

La cellule complexe, aux degrés \(1\to0\), est le cône homologique de
\(\kappa_{\nu,N}\) avec la convention qui place sa source au degré \(1\) :
\begin{equation}\label{eq:pdfv2-coh-cell-complex}
 C_{\nu,N}:=\operatorname{Cone}(\kappa_{\nu,N})
 :=\left[
 A_N^{I_\nu}\oplus A_N^{J_\nu}
 \xrightarrow{\ \kappa_{\nu,N}\ }
 A_N^{I_\nu\times J_\nu}
 \right].
\end{equation}
La proposition précédente donne
\begin{equation}\label{eq:pdfv2-coh-cell-homology}
 H_1(C_{\nu,N})\simeq A_N,
 \qquad
 H_0(C_{\nu,N})\simeq A_N^{(a_\nu-1)(b_\nu-1)}.
\end{equation}
Le défaut de cellule est ainsi construit, non ajouté comme un cardinal sans
application.

\section{Orientation et retenue de la cellule}
\label{sec:pdfv2-coh-orientation-carry}

Soient
\[
 \tau_1(u,v)=(v,u),
 \qquad
 P(w)_{ji}=w_{ij}.
\]
L'échange orienté est l'application de complexes
\begin{equation}\label{eq:pdfv2-coh-theta}
 \Theta_{\nu,N}:=(\tau_1,-P):
 C_{a_\nu,b_\nu,N}\longrightarrow C_{b_\nu,a_\nu,N},
\end{equation}
parce que
\begin{equation}\label{eq:pdfv2-coh-orientation}
 \kappa_{b_\nu,a_\nu,N}\tau_1
 =-P\kappa_{a_\nu,b_\nu,N}.
\end{equation}
En outre,
\(\Theta_{\nu^{\rm op},N}\Theta_{\nu,N}=I\). Le chemin opposé ne
s'identifie pas à l'original : il change le signe de degré zéro et conserve l'acte
d'inversion.

Pour \(r\in A_N\), soit
\(\langle r\rangle_N\in\{0,\ldots,3^N-1\}\) son représentant canonique.
L'emprunt/retenue local est
\begin{equation}\label{eq:pdfv2-coh-beta}
 \beta_N(r,s):=
 \frac{\langle r-s\rangle_N-\langle r\rangle_N+\langle s\rangle_N}
 {3^N}\in\{0,1\}.
\end{equation}
Ainsi,
\begin{equation}\label{eq:pdfv2-coh-carry-cell}
 \langle u_i-v_j\rangle_N
 =\langle u_i\rangle_N-\langle v_j\rangle_N
  +3^N\beta_N(u_i,v_j).
\end{equation}
La matrice
\[
 \operatorname{Carr}_{\nu,N}(u,v)
 :=\bigl(\beta_N(u_i,v_j)\bigr)_{ij}
\]
est conservé avec \(\kappa_{\nu,N}(u,v)\) : résidu et retenue sont deux
coordonnées distinctes.

Lorsque les tailles des ports proviennent des deux feuillets APP, l'admissibilité
de la cellule inclut l'identité entière
\begin{equation}\label{eq:pdfv2-coh-app-defect}
 (a_\nu-1)(b_\nu-1)
 =\rho_{\Pi,\nu}-\rho_{\Sigma,\nu}+1
  +9(c_{\Pi,\nu}-c_{\Sigma,\nu}).
\end{equation}
L'égalité est exigée uniquement pour une cellule ayant cette généalogie ; elle n'est pas affirmée
pour des entiers dépourvus de feuillets APP.

Si \(L_N(u)\) est le relèvement entier canonique de la matrice d'une flèche
\(u\), on définit la retenue de composition par
\begin{equation}\label{eq:pdfv2-coh-carry-composition}
 L_N(v)L_N(u)=L_N(vu)+3^N c_N(v,u).
\end{equation}
L'associativité de trois flèches donne le cocycle exact
\begin{equation}\label{eq:pdfv2-coh-carry-cocycle}
 c_N(w,vu)+L_N(w)c_N(v,u)
 =c_N(wv,u)+c_N(w,v)L_N(u).
\end{equation}

\section{Catégorie d'histoires et foncteur cellulaire}
\label{sec:pdfv2-coh-category}

Pour un préfixe \(x_{\leq d}\), soit \(J_d(x)\) la catégorie finie dont les
objets sont toutes ses cellules \(\nu\). Ses flèches sont engendrées par :
\begin{enumerate}
 \item des raffinements \(u:\nu\to\nu'\) avec des applications surjectives de
       ports
       \[
        \alpha_u:I_{\nu'}\twoheadrightarrow I_\nu,
        \qquad
        \beta_u:J_{\nu'}\twoheadrightarrow J_\nu;
       \]
 \item les échanges d'orientation \(\Theta_\nu\) ;
 \item les extensions nonadiques \(\Gamma_9\nu\) ;
 \item l'identité et la composition, soumises aux relations de frontière,
       de chemin, d'orientation et de retenue de l'histoire.
\end{enumerate}
La finitude se rapporte au préfixe ; le système de tous les préfixes n'est pas
déclaré fini.

Le foncteur cellulaire est
\begin{equation}\label{eq:pdfv2-coh-gamma-functor}
 \Gamma_{d,N}:J_d(x)\longrightarrow
 \operatorname{Ch}^{[0,1]}(A_N),
 \qquad
 \Gamma_{d,N}(\nu)=C_{\nu,N}.
\end{equation}
Pour une flèche de raffinement \(u:\nu\to\nu'\), il agit par
\begin{align}
 \Gamma_{d,N}(u)_1(p,q)
 &:=(p\circ\alpha_u,q\circ\beta_u),
 \label{eq:pdfv2-coh-gamma-one}\\
 \Gamma_{d,N}(u)_0(w)
 &:=w\circ(\alpha_u\times\beta_u).
 \label{eq:pdfv2-coh-gamma-zero}
\end{align}
On vérifie coordonnée par coordonnée
\begin{equation}\label{eq:pdfv2-coh-gamma-chain-map}
 \kappa_{\nu',N}\Gamma_{d,N}(u)_1
 =\Gamma_{d,N}(u)_0\kappa_{\nu,N}.
\end{equation}
Sur une inversion d'orientation, il agit par
\(\Theta_{\nu,N}\), et sur les compositions, il agit par composition, en gardant
séparément \(c_N(v,u)\).

La réduction \(r_{N+1,N}:A_{N+1}\twoheadrightarrow A_N\) satisfait
\begin{equation}\label{eq:pdfv2-coh-kappa-reduction}
 r_{N+1,N}^{I\times J}\kappa_{\nu,N+1}
 =\kappa_{\nu,N}
  \bigl(r_{N+1,N}^{I}\oplus r_{N+1,N}^{J}\bigr).
\end{equation}
Par conséquent, le diagramme des cellules est naturel simultanément en histoire et en
précision.

\fi
\section{Module des chemins, totalisation cellulaire et terminal nonadique}
\label{sec:pdfv2-coh-route-module}

Soit \(\operatorname{Term}_d(x)\) l'ensemble des terminaux de mémoire du
préfixe. Pour \(\nu\in J_d(x)\), définissons
\begin{equation}\label{eq:pdfv2-coh-w-module}
 W_{d,N}(\nu):=
 A_N\bigl[
 \operatorname{Path}_{J_d(x)}(\nu,\operatorname{Term}_d(x))
 \bigr].
\end{equation}
C'est le module libre sur tous les chemins terminaux ; il n'en sélectionne aucun. Si
\(u:\nu\to\nu'\), l'action contravariante est
\begin{equation}\label{eq:pdfv2-coh-w-action}
 W_{d,N}(u):W_{d,N}(\nu')\longrightarrow W_{d,N}(\nu),
 \qquad
 [\lambda]\longmapsto[\lambda\circ u].
\end{equation}
Chaque générateur conserve orientation, retenue, frontière, chemin et mémoire.

\section{Complexe bar bilatéral et totalisation de la bichaîne}
\label{sec:pdfv2-coh-bar}

Pour \(q\geq0\) et \(r\in\{0,1\}\), soit
\begin{equation}\label{eq:pdfv2-coh-bar-bigraded}
 B_{q,r}^{d,N}(x):=
 \bigoplus_{\nu_0\xrightarrow{u_1}\cdots\xrightarrow{u_q}\nu_q}
 W_{d,N}(\nu_q)\otimes_{A_N}\Gamma_{d,N}(\nu_0)_r.
\end{equation}
La différentielle bar est définie sans abréviation par
\begin{align}
 d_{\rm bar}(w;u_1,\ldots,u_q;c)
 :={}&(w;u_2,\ldots,u_q;\Gamma(u_1)c)
 \nonumber\\
 &+\sum_{i=1}^{q-1}(-1)^i
 (w;u_1,\ldots,u_{i+1}u_i,\ldots,u_q;c)
 \nonumber\\
 &+(-1)^q
 (W(u_q)w;u_1,\ldots,u_{q-1};c).
 \label{eq:pdfv2-coh-bar-differential}
\end{align}
Pour \(q=1\), c'est exactement la relation
\begin{equation}\label{eq:pdfv2-coh-coend-relation}
 d_{\rm bar}(w;u;c)
 =w\otimes\Gamma(u)c-W(u)w\otimes c.
\end{equation}
On identifie ainsi une incidence transportée à la même incidence écrite
dans la carte raffinée ; la relation ne reste pas cachée sous le nom du
totalisateur.

Le complexe total est
\begin{equation}\label{eq:pdfv2-coh-total-complex}
 K_{d,N}(x):=\operatorname{Tot}B_{\bullet,\bullet}^{d,N}(x),
 \qquad
 D_{\rm tot}:=d_{\rm bar}+(-1)^q\partial_\Gamma
 \quad\text{sur }B_{q,r}^{d,N}.
\end{equation}
La fonctorialité de \(W\) et \(\Gamma\), avec
\(\partial_\Gamma^2=0\), donne
\begin{equation}\label{eq:pdfv2-coh-total-square-zero}
 D_{\rm tot}^2=0.
\end{equation}
En notation compacte, mais désormais établie par
\eqref{eq:pdfv2-coh-bar-bigraded}--
\eqref{eq:pdfv2-coh-bar-differential},
\begin{equation}\label{eq:pdfv2-coh-derived-coend}
 K_{d,N}(x)
 \simeq
 W_{d,N}\otimes^{\mathbf L}_{A_N[J_d(x)]}\Gamma_{d,N}.
\end{equation}

\section{Colimite de profondeur et terminal nonadique}
\label{sec:pdfv2-coh-terminal}

Les préfixes forment la colimite
\begin{equation}\label{eq:pdfv2-coh-k-infinity}
 K_{\infty,N}(x):=\varinjlim_d K_{d,N}(x).
\end{equation}
L'extension de neuf phases fait passer de la profondeur \(d\) à la profondeur \(d+9\).
Ce n'est qu'après la formation de \eqref{eq:pdfv2-coh-k-infinity} qu'elle induit un
endomorphisme bien typé
\begin{equation}\label{eq:pdfv2-coh-u-gamma9}
 U_{\Gamma_9,N}:K_{\infty,N}(x)
 \longrightarrow K_{\infty,N}(x).
\end{equation}
Le terminal est
\begin{equation}\label{eq:pdfv2-coh-terminal-cone}
 \mathcal T_{{\rm coh},N}(x)
 :=\operatorname{Cone}\bigl(I-U_{\Gamma_9,N}\bigr).
\end{equation}
Avec la convention homologique
\[
 \operatorname{Cone}(f)_q
 =K_{\infty,N,q-1}\oplus K_{\infty,N,q},
\]
sa différentielle est
\begin{equation}\label{eq:pdfv2-coh-cone-differential}
 \partial_{\rm Cone}(a,b)
 =\bigl(-\partial_Ka,\partial_Kb+f(a)\bigr).
\end{equation}

Bien que la phase visible de \(h\) et \(\Gamma_9h\) coïncide, les générateurs
\([h]\) et \([\Gamma_9h]\) sont différents parce que leurs ledgers et leurs profondeurs
sont différents. Le terminal conserve
\begin{equation}\label{eq:pdfv2-coh-terminal-difference}
 (I-U_{\Gamma_9,N})[h]=[h]-[\Gamma_9h].
\end{equation}
Les réductions de précision commutent avec le terminal :
\begin{equation}\label{eq:pdfv2-coh-u-reduction}
 r_{N+1,N}U_{\Gamma_9,N+1}
 =U_{\Gamma_9,N}r_{N+1,N}.
\end{equation}

\section{Lecteur cohomologique terminal : cycle d'incidence et classe dans \(H_0(\operatorname{Cone}(I-U_{\Gamma_9,N}))\)}
\label{sec:pdfv2-coh-reader}

Chaque cellule de \(x_{\leq d}\) porte un vecteur d'incidence
\[
 \omega_{\nu,N}(x)
 \in A_N^{I_\nu\times J_\nu}=C_{\nu,N,0}
\]
et le chemin terminal complet \(\lambda_{\nu,x}\). Le cycle brut est
\begin{equation}\label{eq:pdfv2-coh-z-cycle}
 z_{d,N}(x):=
 \sum_{\nu\in\operatorname{Cell}_d(x)}
 \epsilon_\nu[\lambda_{\nu,x}]
 \otimes\omega_{\nu,N}(x)
 \in B_{0,0}^{d,N}(x).
\end{equation}
Comme le complexe total est homologique non négatif,
\begin{equation}\label{eq:pdfv2-coh-z-closed}
 D_{\rm tot}z_{d,N}(x)=0.
\end{equation}
On conserve aussi bien le cycle que sa classe
\[
 [z_{d,N}(x)]\in H_0(K_{d,N}(x)).
\]
Après son inclusion dans \(K_{\infty,N}\), sa classe terminale est
\[
 [(0,z_{d,N}(x))]
 \in H_0(\mathcal T_{{\rm coh},N}(x)).
\]
Le lecteur complet est
\begin{equation}\label{eq:pdfv2-coh-reader-full}
 \mathcal L_{{\rm coh},d,N}(x)
 :=\bigl(
 z_{d,N}(x),[z_{d,N}(x)],
 U_{\Gamma_9,N}z_{d,N}(x),
 [(0,z_{d,N}(x))]
 \bigr).
\end{equation}
On ne publie pas uniquement une classe quotient : le représentant, le
chemin qui l'a produit, sa continuation et le terminal subsistent.

\section{Fibres survivantes et corrections sans sélecteur}
\label{sec:pdfv2-coh-surviving-fibers}

Soit \(\mathscr S_\lambda^{\rm coh}\) le \(A_N\)-module libre engendré par les
histoires finies admissibles au niveau \(\lambda=(d,N)\). Le lecteur se
prolonge linéairement à
\begin{equation}\label{eq:pdfv2-coh-e-lambda}
 e_\lambda:\mathscr S_\lambda^{\rm coh}
 \longrightarrow\mathscr O_\lambda^{\rm coh},
 \qquad
 \mathscr O_\lambda^{\rm coh}
 :=H_0(K_{d,N})\oplus H_0(\mathcal T_{{\rm coh},N}).
\end{equation}
Le codomaine survivant est
\begin{equation}\label{eq:pdfv2-coh-surviving-output}
 \mathscr O_\infty^{\rm surv}
 :=\operatorname{Im}\!\left(
 e_\infty:\mathcal C_{\rm cont}^{\rm coh}
 \longrightarrow\varprojlim_\lambda\mathscr O_\lambda^{\rm coh}
 \right).
\end{equation}
Pour \(a=(a_\lambda)_\lambda\in\mathscr O_\infty^{\rm surv}\), on définit
\begin{equation}\label{eq:pdfv2-coh-fiber}
 \mathscr F_\lambda(a):=
 \left\{s_\lambda\in\mathscr S_\lambda^{\rm coh}:
 \begin{array}{l}
 e_\lambda(s_\lambda)=a_\lambda,\\
 \exists y\in\mathcal C_{\rm cont}^{\rm coh}\text{ avec }
 y|_\lambda=s_\lambda\text{ et }e_\infty(y)=a
 \end{array}
 \right\}.
\end{equation}
Ces fibres sont non vides par la définition du codomaine survivant. Si
\(\lambda'\geq\lambda\), la restriction est surjective :
\begin{equation}\label{eq:pdfv2-coh-fiber-surjection}
 \rho_{\lambda',\lambda}:\mathscr F_{\lambda'}(a)
 \twoheadrightarrow\mathscr F_\lambda(a).
\end{equation}
En effet, une histoire complète qui témoigne de
\(s_\lambda\in\mathscr F_\lambda(a)\) fournit par troncature un
antécédent à tout niveau ultérieur. Cet antécédent n'est pas choisi comme
partie de l'objet.

La multisection complète est
\begin{equation}\label{eq:pdfv2-coh-extension-multisection}
 \Sigma_{\lambda',\lambda}^{\rm coh}(s_\lambda;a)
 :=\{s_{\lambda'}\in\mathscr F_{\lambda'}(a):
 \rho_{\lambda',\lambda}s_{\lambda'}=s_\lambda\}.
\end{equation}
Pour un prolongement provisoire \(\widetilde s_{\lambda'}\), soit
\begin{equation}\label{eq:pdfv2-coh-discrepancy}
 \delta_{\lambda'}:=
 a_{\lambda'}-e_{\lambda'}(\widetilde s_{\lambda'}).
\end{equation}
La relation de toutes les corrections est
\begin{equation}\label{eq:pdfv2-coh-correction-relation}
 \operatorname{Corr}_{\lambda',\lambda}
 (\widetilde s_{\lambda'},\delta_{\lambda'})
 :=\left\{\eta\in\mathscr S_{\lambda'}^{\rm coh}:
 \begin{array}{l}
 \rho_{\lambda',\lambda}\eta=0,\\
 e_{\lambda'}(\eta)=\delta_{\lambda'},\\
 \eta\text{ est supportée sur de nouvelles cellules},\\
 \eta\text{ respecte l’orientation, la retenue, la mémoire et le terminal}
 \end{array}
 \right\}.
\end{equation}
Si deux relèvements sont comparables, leur différence
\(\eta=s_{\lambda'}-\widetilde s_{\lambda'}\) appartient à cette relation
exactement lorsqu'elle satisfait les quatre conditions précédentes. On n'affirme pas
que \eqref{eq:pdfv2-coh-correction-relation} soit non vide pour un
écart arbitraire : la non-vacuité déterminante est celle de
\eqref{eq:pdfv2-coh-extension-multisection} sur le domaine survivant.

\iffalse
\section{Les treize coordonnées de la restriction}
\label{sec:pdfv2-coh-thirteen}

Pour \(x_\lambda\in\mathcal C_{{\rm cont},\lambda}^{\rm coh}\), soit
\begin{equation}\label{eq:pdfv2-coh-d-package}
 \begin{aligned}
 D_{{\rm coh},\lambda}(x_\lambda):=\bigl(&
 \mathcal F_{{\rm coh},\lambda},
 \operatorname{Ext}_{\Gamma,{\rm coh},\lambda},
 \operatorname{Rel}_{{\rm coh},\lambda},
 \mathcal N_{{\rm coh},\lambda},
 \operatorname{or}_{{\rm coh},\lambda},
 \operatorname{Carr}_{{\rm coh},\lambda},
 \operatorname{Mem}_{{\rm coh},\lambda},\\
 &\partial_{{\rm coh},\lambda},
 \gamma_{{\rm coh},\lambda},
 \operatorname{Msect}_{{\rm coh},\lambda},
 \operatorname{Surv}_{{\rm coh},\lambda},
 \operatorname{Term}_{{\rm coh},\lambda},
 \mathcal L_{{\rm coh},\lambda}\bigr).
 \end{aligned}
\end{equation}
Ses coordonnées sont :
\begin{enumerate}
 \item \(\mathcal F_{{\rm coh},\lambda}\) : tous les modules
       \(A_N^{I_\nu}\), \(A_N^{J_\nu}\),
       \(A_N^{I_\nu\times J_\nu}\), les complexes \(C_{\nu,N}\), les
       modules \(W_{d,N}(\nu)\) et \(\mathscr S_\lambda^{\rm coh}\).
 \item \(\operatorname{Ext}_{\Gamma,{\rm coh},\lambda}\) : toutes les applications
       \(\Gamma(u)\), les réductions de coefficients, les troncatures, les extensions
       \(\Gamma_9\) et les relations \(\Sigma_{\lambda',\lambda}\).
 \item \(\operatorname{Rel}_{{\rm coh},\lambda}\) :
       \(\kappa\), \(\delta\), la suite exacte, les faces bar, la
       relation de cofin, l'orientation et les lois de composition.
 \item \(\mathcal N_{{\rm coh},\lambda}\):
       \[
       (d,N,3^N,a_\nu,b_\nu,(a_\nu-1)(b_\nu-1),
       \rho_{\Sigma,\nu},\rho_{\Pi,\nu},c_{\Sigma,\nu},c_{\Pi,\nu}).
       \]
 \item \(\operatorname{or}_{{\rm coh},\lambda}\) : les signes
       \(\epsilon_\nu\), les ordres des ports et les applications
       \(\Theta_{\nu,N}\).
 \item \(\operatorname{Carr}_{{\rm coh},\lambda}\) : les matrices
       \(\beta_N\), les cocycles \(c_N(v,u)\), les résidus et les quotients,
       conservés séparément.
 \item \(\operatorname{Mem}_{{\rm coh},\lambda}\) : le ledger hérité et le
       mot ordonné des raffinements, retenues, corrections et avancées
       \(\Gamma_9\), non seulement leur valeur agrégée.
 \item \(\partial_{{\rm coh},\lambda}\) :
       \(\kappa_{\nu,N}\), \(D_{\rm tot}\),
       \(\partial_{\rm Cone}\) et les frontières des cellules.
 \item \(\gamma_{{\rm coh},\lambda}\) : toutes les chaînes composables
       \(\nu_0\to\cdots\to\nu_q\) et leurs continuations terminales.
 \item \(\operatorname{Msect}_{{\rm coh},\lambda}\) : toutes les cellules,
       les chemins terminaux, les extensions \(\Sigma\) et les corrections
       \(\operatorname{Corr}\), sans sélecteur.
 \item \(\operatorname{Surv}_{{\rm coh},\lambda}\) : le système des fibres
       \(\mathscr F_\lambda(a)\) et des restrictions surjectives.
 \item \(\operatorname{Term}_{{\rm coh},\lambda}\) :
       \(\operatorname{Cone}(I-U_{\Gamma_9,N})\), avec \(z\), \(Uz\)
       et leur différence.
 \item \(\mathcal L_{{\rm coh},\lambda}\) :
       \((z,[z],Uz,[(0,z)])\) et toutes ses compatibilités de réduction.
\end{enumerate}

Pour \(\lambda'\geq\lambda\), le transport coordonné satisfait
\begin{equation}\label{eq:pdfv2-coh-d-naturality}
 u_{\lambda',\lambda}^{\rm coh}
 D_{{\rm coh},\lambda'}(x_{\lambda'})
 =D_{{\rm coh},\lambda}
  (\rho_{\lambda',\lambda}x_{\lambda'}).
\end{equation}

\section{Graphe de restriction et objet dépendant}
\label{sec:pdfv2-coh-graphs}

La restriction totale est le graphe
\begin{equation}\label{eq:pdfv2-coh-g-graph}
 \mathcal G_{\rm coh}:=\operatorname{Graph}(D_{\rm coh}),
 \qquad
 \operatorname{Res}_{\rm coh}(x):=(x,D_{\rm coh}x).
\end{equation}
Ainsi,
\begin{equation}\label{eq:pdfv2-coh-res-inverse}
 \pi_1\operatorname{Res}_{\rm coh}=I,
 \qquad
 \operatorname{Res}_{\rm coh}\pi_1=I_{\mathcal G_{\rm coh}}.
\end{equation}

L'extracteur complet est
\begin{equation}\label{eq:pdfv2-coh-chi}
 \chi_{\rm coh}(x):=
 \bigl(
 \{C_{\nu,N}\},
 \{z_\lambda,[z_\lambda]\},
 \{\mathscr F_\lambda(e_\infty x)\},
 \{\Sigma_{\lambda',\lambda}\},
 \{\operatorname{Corr}_{\lambda',\lambda}\}
 \bigr).
\end{equation}
On définit
\begin{equation}\label{eq:pdfv2-coh-xchi}
 X_{\chi,{\rm coh}}
 :=\{(\chi_{\rm coh}(x),x,D_{\rm coh}x):
 x\in\mathcal C_{\rm cont}^{\rm coh}\}
\end{equation}
et
\begin{equation}\label{eq:pdfv2-coh-prx}
 \operatorname{pr}_{X,{\rm coh}}(x,D_{\rm coh}x)
 :=(\chi_{\rm coh}(x),x,D_{\rm coh}x).
\end{equation}
La projection sur \((x,D_{\rm coh}x)\) est inverse sur cet objet
dépendant.

\section{Assembleur et publicateur internes}
\label{sec:pdfv2-coh-assembler-publisher}

L'espace ambiant \(\mathscr M_{\rm coh}^{\rm amb}\) se compose de systèmes
compatibles de suites exactes de cellules, de catégories, de foncteurs, de modules de
chemins, de totalisateurs, de terminaux, de lecteurs, de fibres et de multisections.
L'assembleur brut est
\begin{equation}\label{eq:pdfv2-coh-assembler}
 \begin{aligned}
 a_{\rm coh}:X_{\chi,{\rm coh}}&\longrightarrow
 \mathscr M_{\rm coh}^{\rm amb},\\
 a_{\rm coh}(\chi(x),x,Dx)
 &:={}
 \bigl(
 \{C_{\nu,N},\kappa_{\nu,N},\delta_{\nu,N}\},
 \{J_d,\Gamma,W,K\},\\
 &\hspace{31mm}
 \{U_{\Gamma_9,N},\mathcal T_N\},
 \{e_\lambda,\mathscr F_\lambda,\Sigma,\operatorname{Corr}\}
 \bigr)_x.
 \end{aligned}
\end{equation}
La macrostructure est un autre graphe :
\begin{equation}\label{eq:pdfv2-coh-macro-graph}
 \mathfrak M_{\rm coh}:=\operatorname{Graph}(a_{\rm coh}),
 \qquad
 \mathcal A_{\rm coh}(z):=(z,a_{\rm coh}z).
\end{equation}

Soit \(\mathscr Y_{\rm coh}^{\rm int}\) le type de certificats qui contient
\begin{equation}\label{eq:pdfv2-coh-certificate-list}
 \begin{gathered}
 \text{la sucesión exacta \eqref{eq:pdfv2-coh-exact-sequence}},\\
 \kappa_{b,a}\tau=-P\kappa_{a,b},
 \qquad D_{\rm tot}^2=0,
 \qquad\partial_{\rm Cone}^2=0,\\
 rU_{\Gamma_9}=U_{\Gamma_9}r,
 \qquad
 e_\lambda\rho_{\lambda',\lambda}
 =r_{\lambda',\lambda}e_{\lambda'},\\
 \rho_{\lambda',\lambda}:\mathscr F_{\lambda'}(a)
 \twoheadrightarrow\mathscr F_\lambda(a),
 \qquad
 \Sigma_{\lambda',\lambda}(s_\lambda;a)\neq\varnothing,
 \end{gathered}
\end{equation}
avec les cycles bruts, leurs classes et leurs terminaux. Le publicateur
\begin{equation}\label{eq:pdfv2-coh-publisher}
 p_{\rm coh}^{\rm raw}:\mathfrak M_{\rm coh}
 \longrightarrow\mathscr Y_{\rm coh}^{\rm int}
\end{equation}
produit ce certificat sans supprimer la macrostructure. Enfin,
\begin{equation}\label{eq:pdfv2-coh-output-graph}
 \mathcal C_{\rm coh}^{\rm HMT}
 :=\operatorname{Graph}(p_{\rm coh}^{\rm raw}),
 \qquad
 \mathcal P_{\rm coh}(m):=(m,p_{\rm coh}^{\rm raw}m).
\end{equation}

\section{Théorème interne de génération cohérente}
\label{sec:pdfv2-coh-theorem}

\begin{theorem}[Génération interne de la macrostructure de cohérence]
\label{thm:pdfv2-coh-internal-generation}
Sur le domaine explicite \(\mathcal C_{\rm cont}^{\rm coh}\), la chaîne
\[
 \mathcal C_{\rm cont}^{\rm coh}
 \xrightarrow{\operatorname{Res}_{\rm coh}}
 \mathcal G_{\rm coh}
 \xrightarrow{\operatorname{pr}_{X,{\rm coh}}}
 X_{\chi,{\rm coh}}
 \xrightarrow{\mathcal A_{\rm coh}}
 \mathfrak M_{\rm coh}
 \xrightarrow{\mathcal P_{\rm coh}}
 \mathcal C_{\rm coh}^{\rm HMT}
\]
construit naturellement les cellules, leurs cônes, le diagramme d'histoires, le
bar bilatéral, la totalisation, le terminal nonadique et les fibres
survivantes. Elle conserve les treize coordonnées et toutes les corrections comme
multisections. Chaque étape admet, sur son image, une projection qui récupère
l'étape précédente.
\end{theorem}

\begin{proof}
La \cref{prop:pdfv2-coh-cell-exact} construit chaque cellule sur \(A_N\), avec
noyau diagonal et défaut \((a_\nu-1)(b_\nu-1)\). L'identité
\eqref{eq:pdfv2-coh-orientation} fait de l'échange orienté une application de
complexes involutive. Les formules
\eqref{eq:pdfv2-coh-gamma-one}--\eqref{eq:pdfv2-coh-gamma-zero} prouvent que
chaque raffinement est une application de complexes, tandis que
\eqref{eq:pdfv2-coh-kappa-reduction} prouve la compatibilité en précision.

La précomposition des chemins rend \(W\) contravariant. Les identités
fonctorielles de \(W\) et \(\Gamma\) impliquent
\(d_{\rm bar}^2=0\) ; leur commutation avec la différentielle cellulaire donne
\(D_{\rm tot}^2=0\). Ainsi, \(K_{d,N}\) est construit par les formules bar,
non par une invocation nominale du cofin.

L'extension \(\Gamma_9\) augmente la profondeur. Elle n'induit donc
l'endomorphisme \(U_{\Gamma_9,N}\) qu'après le passage à la colimite
\(K_{\infty,N}\). Le cône
\(\operatorname{Cone}(I-U_{\Gamma_9,N})\) conserve alors la différence
\([h]-[\Gamma_9h]\) et n'identifie pas la répétition de phase à la répétition
d'état.

Le cycle \(z_{d,N}(x)\) se forme avec toutes les cellules, tous les signes, chemins et
incidences de l'histoire. Troncature et réduction commutent avec le lecteur. Si
\(s_\lambda\in\mathscr F_\lambda(a)\), l'histoire complète qui apparaît dans
la définition même de la fibre fournit un antécédent à tout niveau
\(\lambda'\geq\lambda\) ; cela prouve la surjectivité
\eqref{eq:pdfv2-coh-fiber-surjection}. L'objet conserve l'ensemble entier
\(\Sigma_{\lambda',\lambda}\), non un choix. Les différences entre
relèvements sont enregistrées par
\eqref{eq:pdfv2-coh-correction-relation}, avec support, orientation, retenue,
mémoire et terminal explicites.

Enfin, \(\mathcal G_{\rm coh}\), \(\mathfrak M_{\rm coh}\) et
\(\mathcal C_{\rm coh}^{\rm HMT}\) sont des graphes, tandis que
\(X_{\chi,{\rm coh}}\) conserve comme coordonnées \(x\) et \(D_{\rm coh}x\).
La première projection de chaque objet récupère le précédent. Par conséquent, la
composition complète ne remplace pas l'histoire par une classe, la multisection par un
sélecteur ni le terminal par la phase visible.
\end{proof}

\section{Provenance probatoire et falsificateurs}
\label{sec:pdfv2-coh-provenance-falsifiers}

\paragraph{Résultat récupéré.}
Appartiennent à l'architecture déjà construite la cellule \(\kappa_{a,b}\), son
noyau diagonal et son défaut, l'orientation
\(\kappa_{b,a}\tau=-P\kappa_{a,b}\), la retenue APP, le cône de la cellule, la
catégorie d'histoires, la totalisation par cofin, l'action \(\Gamma_9\), le
terminal \(\operatorname{Cone}(I-U_{\Gamma_9})\) et l'architecture
d'extension survivante.

\paragraph{Formalisation nouvelle.}
Sont des formalisations explicites la suite exacte
\eqref{eq:pdfv2-coh-exact-sequence} sur \(A_N\), la séparation entre
profondeur et précision, les faces complètes du bar, la formation du
terminal seulement après la colimite, la mise en paquet des treize coordonnées,
les corrections comme relation sans sélecteur et les quatre objets dépendants
ou graphes qui empêchent la perte de généalogie.

\begin{criterion}[Falsificateurs internes de la branche \(\mathrm{coh}\)]
\label{crit:pdfv2-coh-falsifiers}
La construction échoue si l'un quelconque des faits suivants se produit :
\begin{enumerate}
 \item la suite
       \eqref{eq:pdfv2-coh-exact-sequence} n'est pas exacte ;
 \item \(\kappa_{b,a}\tau=-P\kappa_{a,b}\) échoue ;
 \item un raffinement ne satisfait pas
       \(\kappa_{\nu'}\Gamma(u)_1=\Gamma(u)_0\kappa_\nu\) ;
 \item la retenue de composition ne satisfait pas
       \eqref{eq:pdfv2-coh-carry-cocycle} ;
 \item \(D_{\rm tot}^2\neq0\);
 \item la réduction, le lecteur ou le terminal ne commutent pas avec \(U_{\Gamma_9}\) ;
 \item on forme \(\operatorname{Cone}(I-U_{\Gamma_9})\) en identifiant d'abord
       les profondeurs \(d\) et \(d+9\) ;
 \item une fibre déclarée survivante est vide ou une restriction n'est pas
       surjective ;
 \item on remplace \(\Sigma_{\lambda',\lambda}\) par un unique
       relèvement ;
 \item on publie seulement \([z]\) et l'on supprime \(z\), le chemin, la retenue, la mémoire ou
       le terminal ;
 \item l'une quelconque des treize coordonnées est omise.
\end{enumerate}
Dans le dernier cas, le diagnostic contraignant est : \emph{objet substitué ;
conclusion non transférable}.
\end{criterion}

\begin{warningbox}
Ce chapitre conclut uniquement la construction interne de
\(\mathcal C_{\rm coh}^{\rm HMT}\). Il n'introduit pas d'espace classique, de
classe classique ni de publication classique. Tout entrelaceur ultérieur
appartient à un autre document et ne gouverne pas l'existence de cette
macrostructure.
\end{warningbox}
\fi

\endgroup

Le complexe local \eqref{eq:pdfv2-cor-complejo-local} et le sommet \(y_\alpha\), extrémité de cette même extension dans le nerf, forment le produit fibré augmenté. On y définit
\begin{equation}
\label{eq:pdfv2-cor-zpm-h}
 \mathbf1_\alpha=(r,[y_\alpha]),\qquad
 z_-(\alpha)=(r,[y_\alpha]),
\end{equation}
\begin{equation}
\label{eq:pdfv2-cor-zplus-h}
 z_+(\alpha)=
 (r+3c_\alpha v_\alpha e+v_\alpha b,[y_\alpha]),
 \qquad
 h_*(\alpha)=(-v_\alpha f,0).
\end{equation}
L'action sur les générateurs donne, sans affirmation d'existence abstraite supplémentaire,
\begin{equation}
\label{eq:pdfv2-cor-filler}
 \partial z_-=\partial z_+=0,\qquad
 \partial h_*=z_--z_+.
\end{equation}
Le remplissage corrélatif, la cellule et le collier conservent le même \((\alpha,\operatorname{or}(\alpha),\kappa(\alpha),\partial\alpha,
\gamma\alpha,m(\gamma\alpha;m_0))\) ; ils appliquent seulement des lecteurs internes distincts au registre commun.

Pour le même complexe fini muni de bases, on définit, sans exiger de bilan ni d'acyclicité, la droite déterminant totale
\begin{equation}
\label{eq:pdfv2-cor-linea-determinante-total}
 \operatorname{Det}(C_{k,N}):=
 \bigotimes_q
 \left(\bigwedge^{\rm top}C_{k,N,q}\right)^{(-1)^q}.
\end{equation}
Le raffinement induit le morphisme déterminantal de l'application de complexes et transporte simultanément les modules d'homologie, les formes normales de Smith et les pages de Bockstein comme champs du registre. Ni l'existence de \eqref{eq:pdfv2-cor-linea-determinante-total} ni son transport n'exigent l'annulation de l'homologie. Si le cône du raffinement est contractile et orienté, le morphisme se trivialise canoniquement ; dans le cas général, le cône est conservé sans être remplacé par un isomorphisme fictif.

\subsection{Forme normale de Smith de l'opérateur terminal et décomposition primaire de son conoyau}
Pour chaque préfixe \(y\), la retenue non orientée et le coefficient orienté de frontière du registre commun admettent les relèvements
\begin{equation}
\label{eq:detint-c-v-lifts}
 \widetilde c_y,\widetilde v_y\in\mathbf Z,\qquad
 c_{y,n}:=\widetilde c_y\bmod3^n,\qquad
 v_{y,n}:=\widetilde v_y\bmod3^n.
\end{equation}
La sortie déterminantale n'est ni une entrée ni une branche séparée : c'est le secteur produit par la même réduction terminale lorsque les conditions typées sont satisfaites
\begin{equation}
\label{eq:detint-dominio}
\begin{split}
 \mathcal C_{\rm cont}^{\rm det}:=\bigl\{x\in\mathcal C_{\rm cont}^{\rm disc}:{}&
 x\text{ survit à toute profondeur};\\
 &\widehat P_x^{\rm red}\text{ est fini et basé};\\
 &\chi(\widehat P_x^{\rm red})=0;\\
 &H_*(\widehat P_x^{\rm red}\otimes_{\mathbf Z_3}K_3)=0;\\
 &\text{la réduction produit un bloc carré stable par raffinement}
 \bigr\}.
\end{split}
\end{equation}
\begingroup
  \let\section\subsection
  % Extracción quirúrgica del propietario V2 05e: líneas 325--447.
% El resto permanece en este archivo como procedencia, pero no se publica.
\iffalse
\chapter[Macrostructure déterminantielle interne]{Macrostructure déterminantielle interne engendrée par le continu discret}
\label{chap:detint-macroestructura}

\begin{thesisbox}
La branche \({\rm det}\) part exclusivement d'histoires survivantes de
\(\mathcal C_{\rm cont}^{\rm disc}\). Le nerf des extensions, la diagonale
Alexander--Whitney, le complexe local, l'unité, les deux publications, le
filler, la réduction terminale et la tour \(3\)-adique se construisent avant
toute évaluation ultérieure. Les conditions de bilan et d'acyclicité sont
incorporées au domaine typé et ne sont jamais tacitement présupposées.
\end{thesisbox}

\section{Domaine déterminantiel survivant}

Soit
\begin{equation}\label{eq:detint-an}
 A_n:=\mathbf Z/3^n\mathbf Z,
 \qquad
 \rho_{n+1,n}:A_{n+1}\longrightarrow A_n.
\end{equation}
Pour \(x\in\mathcal C_{\rm cont}^{\rm disc}\), soit
\(\mathsf T_m(x)\) la catégorie finie des préfixes TPK de profondeur \(m\)
compatibles avec \(x\). Chaque objet conserve feuillet, orientation, résidu,
quotient, retenue entière, mémoire, frontière, chemin, multisection, survie
et terminal ; ses morphismes sont les extensions admissibles qui conservent ces
coordonnées.

De chaque objet \(y\in\mathsf T_m(x)\), on extrait les entiers internes
\begin{equation}\label{eq:detint-c-v-lifts}
 \widetilde c_y,\widetilde v_y\in\mathbf Z,
 \qquad
 c_{y,n}:=\widetilde c_y\bmod3^n,
 \qquad
 v_{y,n}:=\widetilde v_y\bmod3^n.
\end{equation}
La première coordonnée est la retenue non orientée et la seconde, le coefficient
orienté de frontière. Le renversement satisfait
\begin{equation}\label{eq:detint-or-cv}
 c_{\bar y,n}=c_{y,n},
 \qquad
 v_{\bar y,n}=-v_{y,n}.
\end{equation}

On écrit
\begin{equation}\label{eq:detint-z3-k3}
 \mathbf Z_3:=\varprojlim_n A_n,
 \qquad
 K_3:=\operatorname{Frac}(\mathbf Z_3).
\end{equation}
La restriction déterminantielle est définie par
\begin{equation}\label{eq:detint-dominio}
\begin{split}
 \mathcal C_{\rm cont}^{\rm det}:=
 \bigl\{x\in\mathcal C_{\rm cont}^{\rm disc}:{}&
 x\text{ survit à toute profondeur};\\
 &\widehat P_x^{\rm red}\text{ est fini et basé};\\
 &\chi(\widehat P_x^{\rm red})=0;\\
 &H_*(\widehat P_x^{\rm red}\otimes_{\mathbf Z_3}K_3)=0;\\
 &\text{la réduction canonique produit un bloc carré}\\
 &\text{stable par raffinement}\bigr\}.
\end{split}
\end{equation}
Les objets \(\widehat P_x^{\rm red}\) et le bloc carré seront construits
dans \cref{sec:detint-terminal}. Par conséquent,
\eqref{eq:detint-dominio} est une condition d'appartenance vérifiable et non
une conclusion insérée dans la définition des applications.

\section{Nerf, Alexander--Whitney et counité}

On définit le complexe normalisé
\begin{equation}\label{eq:detint-gmn}
 G_{m,n}(x):=
 N_*^{\rm norm}\!\left(N\mathsf T_m(x);A_n\right).
\end{equation}
Pour un simplexe non dégénéré
\(\sigma=[y_0\to y_1\to\cdots\to y_q]\), la diagonale est
\begin{equation}\label{eq:detint-aw}
 \Delta_{\rm AW}(\sigma)
 =\sum_{j=0}^q
 [y_0\to\cdots\to y_j]\otimes
 [y_j\to\cdots\to y_q].
\end{equation}
La counité est donnée par
\begin{equation}\label{eq:detint-counidad-g}
 \epsilon_G([y])=1,
 \qquad
 \epsilon_G(\sigma)=0\quad(q>0).
\end{equation}
En particulier,
\begin{equation}\label{eq:detint-vertice-grouplike}
 \Delta_{\rm AW}[y]=[y]\otimes[y],
 \qquad
 \epsilon_G([y])=1.
\end{equation}

Si \(m'\ge m\) et \(n'\ge n\), la troncature et la réduction des coefficients
induisent
\begin{equation}\label{eq:detint-rmn}
 R_{m',m}^{n',n}:G_{m',n'}(x)\longrightarrow G_{m,n}(x),
\end{equation}
et l'on vérifie
\begin{equation}\label{eq:detint-aw-natural}
 R\partial=\partial R,
 \qquad
 (R\otimes R)\Delta_{\rm AW}=\Delta_{\rm AW}R,
 \qquad
 \epsilon_GR=\rho_{n',n}\epsilon_G.
\end{equation}
Le chemin ne se réduit pas à ses extrémités : le simplexe entier reste dans
\(G_{m,n}(x)\).

\section{Complexe local interne}

Pour \(c\in A_n\), on définit
\begin{equation}\label{eq:detint-p0-grados}
 P^0_{n,1}(c)=A_nf,
 \qquad
 P^0_{n,0}(c)=A_nr\oplus A_ne\oplus A_nb,
 \qquad
 P^0_{n,-1}(c)=A_ng,
\end{equation}
avec différentielle
\begin{equation}\label{eq:detint-p0-diferencial}
 \partial f=3ce+b,
 \qquad
 \partial r=0,
 \qquad
 \partial e=g,
 \qquad
 \partial b=-3cg.
\end{equation}
L'identité de chaîne est littérale :
\begin{equation}\label{eq:detint-d2}
 \partial^2f=3c\,\partial e+\partial b=3cg-3cg=0,
 \qquad
 \partial^2e=\partial^2b=\partial^2r=0.
\end{equation}

L'augmentation
\begin{equation}\label{eq:detint-epsilon-p}
 \epsilon_P:P_n^0(c)\longrightarrow A_n[0]
\end{equation}
est fixée par
\[
 \epsilon_P(r)=1,
 \qquad
 \epsilon_P(e)=\epsilon_P(b)=\epsilon_P(f)=\epsilon_P(g)=0.
\]
Le complexe possède la coalgèbre différentielle
\begin{equation}\label{eq:detint-coalgebra-p}
 \Delta_P(r)=r\otimes r,
 \qquad
 \Delta_P(q)=q\otimes r+r\otimes q
 \quad(q\in\{e,b,f,g\}).
\end{equation}
Ainsi, \(r\) est group-like et les quatre autres générateurs sont primitifs
relativement à \(r\). Les équations
\eqref{eq:detint-p0-diferencial} prouvent que \(\Delta_P\) commute avec la
différentielle tensorielle.

\section{Système local d'extensions et ledger de mémoire}

Pour une extension \(\alpha:y\to y'\), nous définissons
\begin{equation}\label{eq:detint-psi}
 \Psi_\alpha:P_n^0(c_y)\longrightarrow P_n^0(c_{y'})
\end{equation}
par
\begin{equation}\label{eq:detint-psi-generadores}
\begin{gathered}
 \Psi_\alpha(r_y)=r_{y'},\quad
 \Psi_\alpha(e_y)=e_{y'},\quad
 \Psi_\alpha(g_y)=g_{y'},\quad
 \Psi_\alpha(f_y)=f_{y'},\\
 \Psi_\alpha(b_y)=b_{y'}+3(c_{y'}-c_y)e_{y'}.
\end{gathered}
\end{equation}
C'est une application de complexes parce que
\begin{equation}\label{eq:detint-psi-b-chain}
 \partial\Psi_\alpha(b_y)
 =-3c_{y'}g+3(c_{y'}-c_y)g
 =-3c_yg
 =\Psi_\alpha(\partial b_y),
\end{equation}
et
\begin{equation}\label{eq:detint-psi-f-chain}
 \Psi_\alpha(\partial f_y)
 =3c_ye+b+3(c_{y'}-c_y)e
 =3c_{y'}e+b
 =\partial\Psi_\alpha(f_y).
\end{equation}
De plus,
\begin{equation}\label{eq:detint-psi-functorial}
 \Psi_\beta\Psi_\alpha=\Psi_{\beta\alpha}.
\end{equation}

Le changement de la coordonnée orientée n'est pas effacé. Il est enregistré par
\begin{equation}\label{eq:detint-k-alpha}
 K_\alpha:=(v_{y'}-v_y)f_{y'}.
\end{equation}
Pour des extensions composables,
\begin{equation}\label{eq:detint-k-coherencia}
 K_{\beta\alpha}=K_\beta+\Psi_\beta K_\alpha.
\end{equation}
Cette égalité est le ledger exact de l'accroissement de frontière et rend visible
la mémoire qu'une composition purement nominale perdrait.

Le diagramme local complet est la construction de Grothendieck
\begin{equation}\label{eq:detint-grothendieck-local}
 \int_{y\in\mathsf T_m(x)}P_n^0(c_y).
\end{equation}
Dans celle-ci, une flèche au-dessus de \(\alpha:y\to y'\) porte comme composante linéaire
\(\Psi_\alpha\) et comme cohérence \(K_\alpha\). Lorsqu'elle est appliquée à un
élément supporté en un sommet, on utilise la notation
\begin{equation}\label{eq:detint-psi-soporte-vertice}
 \Psi_\alpha(p,[y]):=(\Psi_\alpha p,[y']).
\end{equation}
L'arête \([y\to y']\) reste dans la coordonnée de chemin du nerf. Cette
convention sera celle utilisée dans les formules de naturalité de
\(z_\pm\) et \(h_*\) ; on n'affirme pas une application globale qui efface les autres
simplexes de \(G_{m,n}(x)\).

\section{Pullback, unité, racine, publications et filler}

Le pullback de complexes est
\begin{equation}\label{eq:detint-pgen}
 P^{\rm gen}_{y;m,n}
 :=P_n^0(c_y)\times_{A_n[0]}G_{m,n}(x)
 =\{(p,\gamma):\epsilon_P(p)=\epsilon_G(\gamma)\}.
\end{equation}
Sa différentielle est
\begin{equation}\label{eq:detint-pgen-d}
 \partial(p,\gamma)=(\partial p,\partial\gamma).
\end{equation}
L'unité commune associée au sommet \(y\) est
\begin{equation}\label{eq:detint-unidad}
 \mathbf1_y:=(r,[y]).
\end{equation}
Ses deux projections sont group-like par
\eqref{eq:detint-vertice-grouplike} et
\eqref{eq:detint-coalgebra-p} ; c'est le sens précis d'unité
group-like compatible du pullback.

La projection racine est l'application de complexes
\begin{equation}\label{eq:detint-proyeccion-root}
 \varpi_{\rm root}:P^{\rm gen}_{y;m,n}\longrightarrow
 \operatorname{Root}_{y;m,n},
 \qquad
 \varpi_{\rm root}(p,\gamma):=(\epsilon_P(p)r,\gamma).
\end{equation}
On définit les deux publications internes
\begin{equation}\label{eq:detint-zpm}
 z_-(y):=(r,[y]),
 \qquad
 z_+(y):=(r+3c_yv_ye+v_yb,[y]).
\end{equation}
Toutes deux sont des cycles :
\begin{equation}\label{eq:detint-zpm-ciclos}
 \partial z_-(y)=0,
 \qquad
 \partial z_+(y)=3c_yv_yg-3c_yv_yg=0.
\end{equation}
Leurs racines coïncident :
\begin{equation}\label{eq:detint-raiz-comun}
 \varpi_{\rm root}z_-(y)
 =(r,[y])
 =\varpi_{\rm root}z_+(y).
\end{equation}

Le filler se construit directement :
\begin{equation}\label{eq:detint-hstar}
 h_*(y):=(-v_yf,0)\in P^{\rm gen}_{y;m,n,1}.
\end{equation}
Alors
\begin{equation}\label{eq:detint-filler-identidad}
 \partial h_*(y)
 =(-3c_yv_ye-v_yb,0)
 =z_-(y)-z_+(y).
\end{equation}
En particulier,
\begin{equation}\label{eq:detint-clases-iguales}
 [z_-(y)]=[z_+(y)]
 \quad\text{dans}\quad H_0(P^{\rm gen}_{y;m,n}),
\end{equation}
et l'égalité conserve sa chaîne témoin.

Sous une extension \(\alpha:y\to y'\), on a
\begin{equation}\label{eq:detint-z-natural}
 z_-(y')=\Psi_\alpha z_-(y),
 \qquad
 z_+(y')-\Psi_\alpha z_+(y)=\partial K_\alpha,
\end{equation}
et
\begin{equation}\label{eq:detint-h-natural}
 h_*(y')-\Psi_\alpha h_*(y)=-K_\alpha.
\end{equation}
Par conséquent, la naturalité n'identifie pas des valeurs distinctes de \(v\) : elle enregistre
leur différence au moyen de l'homotopie cohérente \(K_\alpha\).

\section{Orientation interne}

Sur \(P_n^0(c)\), on définit
\begin{equation}\label{eq:detint-omega-p}
 \Omega(r)=r,
 \qquad
 \Omega(q)=-q\quad(q\in\{e,b,f,g\}).
\end{equation}
C'est un automorphisme de complexes et de la coalgèbre relative. Dans le nerf,
l'inversion orientée est
\begin{equation}\label{eq:detint-omega-g}
 \mathfrak o_G[y_0\to\cdots\to y_q]
 :=(-1)^{q(q+1)/2}
 [\bar y_q\to\cdots\to\bar y_0].
\end{equation}
Avec \eqref{eq:detint-or-cv}, ces actions donnent
\begin{equation}\label{eq:detint-or-z-h}
 \Omega z_-(y)=z_-(\bar y),
 \qquad
 \Omega z_+(c_y,v_y)=z_+(c_{\bar y},v_{\bar y}),
 \qquad
 \Omega h_*(y)=h_*(\bar y).
\end{equation}
L'orientation agit sur le chemin, le complexe, les publications, le filler et la droite
déterminantielle ; ce n'est pas une étiquette sans action.

\fi
\section{Relèvement du conoyau au moyen de la tour de Bockstein et conservation de la retenue}
\label{sec:detint-terminal}

Seule la racine commune est supprimée :
\begin{equation}\label{eq:detint-p-reducido}
 \overline P_{y;m,n}
 :=P^{\rm gen}_{y;m,n}/A_n\mathbf1_y,
 \qquad
 \widehat P_{y;m}^{\rm red}:=\varprojlim_n\overline P_{y;m,n}.
\end{equation}
La base est ordonnée d'abord selon l'ordre TPK des chemins, puis selon
\begin{equation}\label{eq:detint-orden-base}
 f\prec e\prec b\prec g.
\end{equation}
On exécute l'algorithme interne suivant.
\begin{enumerate}
 \item Chercher les coefficients qui sont des unités de \(\mathbf Z_3\).
 \item Choisir le premier selon \eqref{eq:detint-orden-base} et l'ordre des
 chemins.
 \item Annuler le couple correspondant et enregistrer les opérations
 élémentaires, leur signe et leur contribution à l'orientation.
 \item Répéter jusqu'à ce qu'il ne reste aucun pivot unitaire annulable.
\end{enumerate}
La condition de survie de \eqref{eq:detint-dominio} exige que,
cofinalement en \(m\), le résultat soit un bloc à deux termes
\begin{equation}\label{eq:detint-sx}
 \mathbf Z_3^{r_x}\xrightarrow{S_x}\mathbf Z_3^{r_x},
 \qquad
 \det S_x\ne0\text{ dans }K_3,
\end{equation}
et qu'un raffinement n'ajoute que des blocs contractiles unitaires et des changements
de base enregistrés. Si ces conditions échouent, l'histoire n'appartient pas à
\(\mathcal C_{\rm cont}^{\rm det}\) ; on ne force pas un terminal inexistant.

Le chemin identité fournit un témoin élémentaire. Après le retrait de la racine,
le complexe local est
\begin{equation}\label{eq:detint-testigo-contractil}
 A_nf\xrightarrow{\binom{3c}{1}}
 A_ne\oplus A_nb\xrightarrow{(1,-3c)}A_ng.
\end{equation}
Il est contractile par
\begin{equation}\label{eq:detint-contraccion-local}
 s(g)=e,
 \qquad
 s(ae+\beta b)=\beta f,
 \qquad
 s(f)=0.
\end{equation}
Ainsi, les conditions terminales admettent au moins le témoin identité
interne et ne décrivent pas un sous-domaine formellement vide.

\section{Smith, ordre \texorpdfstring{\(3\)}{3}-adique et unité initiale}

Dans des bases orientées, une diagonalisation par valuation a la forme
\begin{equation}\label{eq:detint-smith}
 U_xS_xV_x
 =\operatorname{diag}(3^{a_1}u_1,\ldots,3^{a_r}u_r),
 \qquad
 a_i\in\mathbf N,quad u_i\in\mathbf Z_3^\times.
\end{equation}
On définit
\begin{equation}\label{eq:detint-orden-d}
 d_x:=v_3(\det S_x)=\sum_{i=1}^r a_i
\end{equation}
et l'unité initiale
\begin{equation}\label{eq:detint-leading-unit}
 \lambda_x:=3^{-d_x}\det S_x\in\mathbf Z_3^\times.
\end{equation}
La coordonnée d'orientation trivialise la droite déterminantielle. Si une
transition produit
\begin{equation}\label{eq:detint-estabilidad-bloques}
 S_x\oplus I_q=U S_{x'}V,
\end{equation}
l'orientation est transportée par le facteur inverse de \(\det U\det V\).
L'unité orientée \(\lambda_x^{\rm or}\) est alors naturelle et ne change pas
lors de l'ajout de blocs contractiles.

En précision finie,
\begin{equation}\label{eq:detint-smith-finito}
 S_{x,n}=S_x\bmod3^n,
 \qquad
 a_i^{(n)}=\min(a_i,n).
\end{equation}
Le terminal est donc la famille compatible de toutes les précisions et
non un entier isolé.

\section{Tour de Bockstein et retenue d'ordre supérieur}

Soit
\[
 C_x^1\xrightarrow{S_x}C_x^0
\]
le bloc terminal. Si une classe \([\bar y]\) admet un relèvement
\(y\in C_x^1\) avec \(S_xy\in3^qC_x^0\), nous définissons
\begin{equation}\label{eq:detint-beta-q}
 \beta_q[\bar y]
 :=\left[3^{-q}S_xy\bmod3\right].
\end{equation}
Le changement de relèvement modifie le représentant par un bord de la page
correspondante, de sorte que \(\beta_q\) est bien défini. La retenue
supérieure est
\begin{equation}\label{eq:detint-carry-superior}
 \operatorname{Carr}_q(y):=3^{-q}S_xy.
\end{equation}
Elle ne remplace pas la retenue TPK de \eqref{eq:detint-c-v-lifts} ; elle enregistre son
prolongement dans la tour des coefficients.

Pour un bloc \(3^{a_i}u_i\), la classe persiste jusqu'à la page \(a_i\),
et à cette page, la différentielle non nulle est la multiplication par
\(\bar u_i\in\mathbf F_3^\times\). Par conséquent,
\begin{equation}\label{eq:detint-bock-longitudes}
 \{\text{longueurs de survie}\}=\{a_1,\ldots,a_r\},
 \qquad
 d_x=\sum_i a_i.
\end{equation}
Le produit orienté des trivialisations de page satisfait
\begin{equation}\label{eq:detint-bock-leading}
 \Theta_{\rm Bock}(x)=\overline{\lambda_x^{\rm or}}
 \in\mathbf F_3^\times.
\end{equation}
La preuve se fait dans chaque bloc diagonal et se transporte par les changements
unimodulaires enregistrés. Smith, Bockstein et l'unité initiale sont ainsi trois
lectures du même terminal.

\iffalse
\section{Les treize coordonnées de la restriction déterminantielle}

Pour chaque \((m,n)\), on définit
\begin{equation}\label{eq:detint-d-trece}
 D_{{\rm det},m,n}(x):=
 (\mathcal F,\operatorname{Ext}_\Gamma,\operatorname{Rel},
 \mathcal N,\operatorname{or},\operatorname{Carr},\operatorname{Mem},
 \partial,\gamma,\operatorname{Msect},\operatorname{Surv},
 \operatorname{Term},\mathcal L)_{{\rm det},m,n}(x),
\end{equation}
avec le contenu suivant.
\begin{enumerate}
 \item La fibre est
 \[
  \mathcal F=(A_n,\mathsf T_m(x),
  \{P_n^0(c_y)\}_y,\{P^{\rm gen}_{y;m,n}\}_y).
 \]
 \item \(\operatorname{Ext}_\Gamma\) contient tous les morphismes
 \(\alpha\), les applications \(\Psi_\alpha\), les homotopies \(K_\alpha\) et
 les applications \(R_{m',m}^{n',n}\).
 \item \(\operatorname{Rel}\) contient \(\partial^2=0\), l'équation de
 pullback, AW, la counité, \(\Psi_\beta\Psi_\alpha=\Psi_{\beta\alpha}\) et
 \(K_{\beta\alpha}=K_\beta+\Psi_\beta K_\alpha\).
 \item
 \(\mathcal N=(m,n,\widetilde c_y,\widetilde v_y,c_{y,n},v_{y,n},
 r_x,a_1,\ldots,a_{r_x},d_x)\).
 \item \(\operatorname{or}\) contient \(y\mapsto\bar y\),
 \(\mathfrak o_G\), \(\Omega\) et l'orientation de la droite
 déterminantielle.
 \item \(\operatorname{Carr}\) conserve séparément la retenue entière,
 sa réduction \(c_{y,n}\) et toutes les retenues supérieures
 \(\operatorname{Carr}_q\).
 \item \(\operatorname{Mem}\) contiene
 \[
  y_0\to\cdots\to y_m,
  \quad\{(\widetilde c_{y_j},\widetilde v_{y_j},\mathbf1_{y_j})\}_{j=0}^m,
  \quad\{K_{\alpha_j}\}.
 \]
 \item
 \(\partial=(\partial_G,\partial_P,\partial_{P^{\rm gen}},
 \partial_{\overline P},h_*)\).
 \item \(\gamma=[y_0\to\cdots\to y_q]\) conserve le simplexe orienté
 complet.
 \item \(\operatorname{Msect}\) contient toutes les extensions
 compatibles, tous les relèvements \(3\)-adiques et toutes les présentations
 unimodulaires du même terminal orienté ; il n'en choisit aucune a posteriori.
 \item \(\operatorname{Surv}\) enregistre la profondeur arbitraire,
 la compatibilité des coefficients, le bilan, l'acyclicité sur \(K_3\),
 la stabilité par blocs unitaires et la permanence de la racine.
 \item
 \[
  \operatorname{Term}=(\widehat P_x^{\rm red},S_x,
  \{a_i,u_i\},d_x,\lambda_x^{\rm or},
  \{\beta_q\},\Theta_{\rm Bock}).
 \]
 \item Le lecteur interne est
 \[
  \mathcal L_{\rm det}^{\rm int}
  =(\mathbf1_y,
  \varpi_{\rm root}z_-=\varpi_{\rm root}z_+,
  [z_-]=[z_+],d_x,\lambda_x^{\rm or},\Theta_{\rm Bock}).
 \]
\end{enumerate}

Les transitions satisfont, coordonnée par coordonnée,
\begin{equation}\label{eq:detint-naturalidad-trece}
 u^{\rm det}_{(m',n'),(m,n)}D_{{\rm det},m',n'}(x')
 =D_{{\rm det},m,n}(\tau_{m',m}x').
\end{equation}

\section{Assembleur, publicateur et chaîne non ablative}

On définit
\begin{equation}\label{eq:detint-graph-d}
 \mathcal G_{\rm det}:=\operatorname{Graph}(D_{\rm det}),
 \qquad
 \operatorname{Res}_{\rm det}(x):=(x,D_{\rm det}x).
\end{equation}
Soit \(\chi_{\rm det}(x)\) la multisection complète de germes
\[
 (y,\widetilde c_y,\widetilde v_y,\operatorname{or}_y,\gamma_y)
\]
compatibles à toute profondeur. L'objet dépendant est
\begin{equation}\label{eq:detint-x-dependiente}
 X_{\chi,{\rm det}}
 :=\{(\chi_{\rm det}(x),x,D_{\rm det}x):
 x\in\mathcal C_{\rm cont}^{\rm det}\},
\end{equation}
et
\begin{equation}\label{eq:detint-prx}
 \operatorname{pr}_{X,{\rm det}}(x,D_{\rm det}x)
 :=(\chi_{\rm det}(x),x,D_{\rm det}x).
\end{equation}

L'assembleur brut
\begin{equation}\label{eq:detint-ensamblador-crudo}
 a_{\rm det}:X_{\chi,{\rm det}}\longrightarrow
 \mathscr M_{\rm det}^{\rm amb}
\end{equation}
produit
\begin{equation}\label{eq:detint-ensamblador-salida}
\begin{split}
 a_{\rm det}(z)=\bigl(&
 \{G_{m,n},\Delta_{\rm AW},\epsilon_G\}_{m,n};\\
 &\{P_n^0(c),\Psi_\alpha,K_\alpha,P^{\rm gen},
 \mathbf1,z_-,z_+,h_*\}_{m,n};\\
 &\{\widehat P^{\rm red},S,d,\lambda^{\rm or},
 \beta_q,\Theta_{\rm Bock}\}\bigr).
\end{split}
\end{equation}
Le macro-objet conserve son entrée :
\begin{equation}\label{eq:detint-graph-a}
 \mathfrak M_{\rm det}:=\operatorname{Graph}(a_{\rm det}),
 \qquad
 \mathcal A_{\rm det}(z):=(z,a_{\rm det}z).
\end{equation}

Le publicateur brut
\begin{equation}\label{eq:detint-publicador-crudo}
 p_{\rm det}:\mathfrak M_{\rm det}\longrightarrow\mathscr Y_{\rm det}
\end{equation}
publie conjointement le certificat
\begin{equation}\label{eq:detint-certificado-publicado}
\begin{gathered}
 \partial^2=0,quad \text{AW et counité},quad
 \mathbf1_y\text{ group-like compatible},\\
 \partial z_\pm=0,quad
 \varpi_{\rm root}z_-=\varpi_{\rm root}z_+,quad
 \partial h_*=z_--z_+,\\
 \Psi_\beta\Psi_\alpha=\Psi_{\beta\alpha},quad
 K_{\beta\alpha}=K_\beta+\Psi_\beta K_\alpha,\\
 d=v_3(\det S)=\sum_i a_i,quad
 \Theta_{\rm Bock}=\overline{\lambda^{\rm or}},quad
 \text{naturalité complète par raffinement}.
\end{gathered}
\end{equation}
Enfin,
\begin{equation}\label{eq:detint-graph-p}
 \mathcal C_{\rm det}^{\rm HMT}:=\operatorname{Graph}(p_{\rm det}),
 \qquad
 \mathcal P_{\rm det}(M):=(M,p_{\rm det}M).
\end{equation}
La chaîne intégrale est
\begin{equation}\label{eq:detint-cadena-completa}
 \mathcal C_{\rm cont}^{\rm disc}\supset
 \mathcal C_{\rm cont}^{\rm det}
 \xrightarrow{\operatorname{Res}_{\rm det}}\mathcal G_{\rm det}
 \xrightarrow{\operatorname{pr}_{X,{\rm det}}}X_{\chi,{\rm det}}
 \xrightarrow{\mathcal A_{\rm det}}\mathfrak M_{\rm det}
 \xrightarrow{\mathcal P_{\rm det}}\mathcal C_{\rm det}^{\rm HMT}.
\end{equation}
Chaque flèche possède un inverse sur son image donné par une projection de
coordonnées. On ne remplace pas l'histoire par une valeur, la multisection par un sélecteur,
le complexe par un déterminant ni la preuve par un nom.

\section{Théorème de réalisation déterminantielle interne}

\begin{theorem}[Réalisation déterminantielle interne complète]
\label{thm:detint-realizacion}
Pour tout \(x\in\mathcal C_{\rm cont}^{\rm det}\), la chaîne
\eqref{eq:detint-cadena-completa} construit un macro-objet naturel en
profondeur et en précision. Il contient deux cycles à racine commune, un filler
explicite qui les identifie, une unité group-like compatible et un terminal
\(3\)-adique dont l'ordre, la forme de Smith et la tour de Bockstein produisent le même
lecteur orienté. Les treize coordonnées de \eqref{eq:detint-d-trece}
restent récupérables en sortie.
\end{theorem}

\begin{proof}
Les équations \eqref{eq:detint-p0-diferencial}--
\eqref{eq:detint-d2} prouvent que \(P_n^0(c)\) est un complexe. Les deux
augmentations sont des applications de complexes, de sorte que le pullback est typé.
Alexander--Whitney est naturel et rend chaque sommet group-like ; avec
\(r\), cela donne l'unité commune. Le calcul direct
\eqref{eq:detint-zpm-ciclos}--\eqref{eq:detint-filler-identidad} prouve que
les publications sont des cycles, ont la même racine et sont reliées par
\(h_*\). Les formules de \(\Psi_\alpha\) prouvent la fonctorialité, et
\(K_\alpha\) prouve la naturalité homotopique en conservant l'accroissement de
mémoire. La définition de \(\mathcal C_{\rm cont}^{\rm det}\) garantit que
la réduction terminale produit un bloc carré sans le postuler hors du
domaine. Smith sur \(\mathbf Z_3\) donne \(d=\sum_i a_i\) ; le calcul de chaque
bloc \(3^{a_i}u_i\) donne
\(\Theta_{\rm Bock}=\overline{\lambda^{\rm or}}\). Troncature, réduction et
ajout de blocs unitaires conservent ces identités. Enfin, les
trois constructions par graphe retiennent l'entrée comme première coordonnée,
de sorte que la composition complète est non ablative.
\end{proof}

\section{Provenance et falsificateurs}

\paragraph{Provenance.}
Le nerf TPK, Alexander--Whitney, la counité, l'unité commune, la forme
normale \(\partial f=3ce+b\), \(\partial e=g\),
\(\partial b=-3cg\), la racine commune, le filler et l'architecture de tour
ont le statut \texttt{ARCHITECTURE\_AUTORALE\_PRÉEXISTANTE}. Le calcul
littéral des deux cycles, de leur racine et de l'identité
\(\partial h_*=z_--z_+\) a le statut \texttt{RÉSULTAT\_RÉCUPÉRÉ}. Le
système local \(\Psi_\alpha/K_\alpha\), le sous-domaine survivant
équilibré, le terminal exclusivement \(3\)-adique et la mise en paquet par
graphes ont le statut \texttt{FORMALISATION\_NOUVELLE} ; ils formalisent sans changer
la provenance conceptuelle de l'appareil.

\begin{ablation}[Falsificateurs internes de la branche déterminantielle]
\label{abl:detint-falsadores}
La construction est falsifiée au niveau correspondant si l'un des
faits suivants se produit.
\begin{enumerate}
 \item \(\partial^2=0\) échoue ; le complexe local n'existe alors pas.
 \item Alexander--Whitney ne commute pas avec le raffinement, la counité échoue ou
 le sommet cesse d'être group-like.
 \item \(z_\pm\) ne sont pas des cycles, leurs racines diffèrent ou
 \(\partial h_*=z_--z_+\) échoue.
 \item \(\Psi_\alpha\) n'est pas une application de complexes ou
 \(K_{\beta\alpha}=K_\beta+\Psi_\beta K_\alpha\) échoue ; dans ce cas, on a
 perdu de la retenue ou de la mémoire.
 \item Le complexe réduit n'est pas équilibré ou n'est pas acyclique sur
 \(K_3\) ; l'histoire reste hors du domaine terminal typé.
 \item Les longueurs de Bockstein ne coïncident pas avec les \(a_i\), ou
 \(\Theta_{\rm Bock}\ne\overline{\lambda^{\rm or}}\).
 \item Le raffinement modifie \(d\) ou \(\lambda^{\rm or}\) en dehors de
 blocs unitaires et de changements d'orientation enregistrés.
 \item L'une des treize coordonnées est omise : \emph{objet substitué ;
 conclusion non transférable}.
\end{enumerate}
Ces falsificateurs vérifient l'objet interne construit ; ils ne le remplacent pas
par un contrôle ultérieur.
\end{ablation}
\fi

\endgroup

\section{Naturalité de l'assemblage et compatibilité sous raffinement}

Avant de former le registre d'horizon, nous explicitons l'action totale des cinq opérations développées ici sur \emph{la même} arête \(\mathfrak e_\alpha\). Si \(P_x,Q_x\) désignent les restrictions de \eqref{eq:pdfv2-cor-sigma} à la fibre de \(x=s(\alpha)\), et de même dans \(y=t(\alpha)\), nous posons d'abord
\begin{equation}
\label{eq:pdfv2-cor-a-eta-emision}
 A_\alpha:=P_yU_\alpha P_x+Q_yU_\alpha Q_x,
 \qquad
 \eta_\alpha:=Q_yU_\alpha P_x+P_yU_\alpha Q_x.
\end{equation}
Les cinq lectures constituent des systèmes dépendants indexés par \(\mathfrak e_\alpha\) elle-même. L'indice fixe source, cible, parcours et registre, mais n'est pas recopié dans la valeur. Leurs types sont les registres dépendants \(\operatorname{LedgerEntry}:=\operatorname{cod}(\lambda)\) et
\begin{align*}
 \mathsf Y_{\rm sp}[\alpha]&:=
 \operatorname{Rec}[U:\mathsf H_x\to\mathsf H_y;
 \sigma_x\in\operatorname{End}\mathsf H_x;
 \sigma_y\in\operatorname{End}\mathsf H_y;A,\eta:\mathsf H_x\to\mathsf H_y],\\
 \mathsf Y_{\rm sol}[\alpha]&:=
 \operatorname{Rec}[b,b^+,b^-\in\operatorname{LedgerEntry};
 \operatorname{upd}:\mathsf{Mem}_x\to\mathsf{Mem}_y],\\
 \mathsf Y_{\rm gau}[\alpha]&:=
 \mathcal P_{\rm fin}\operatorname{Rec}[
 \operatorname{Cell}(\alpha);q\in Q_{\rm inc}(y);B;W_c(q)],\\
 \mathsf Y_{\rm coh}[\alpha]&:=
 \operatorname{Rec}[I_y,J_y,\kappa_{y,N},\delta_y,C_{y,N};
 F_\alpha:C_{y,N}\to C_{x,N}],\\
 \mathsf Y_{\rm det}[\alpha]&:=
 \operatorname{Rec}[\Psi_\alpha,K_\alpha,\mathbf1_\alpha,z_\pm,h_*;
 \operatorname{Det},H_*,\operatorname{Smith},\operatorname{Boc}].
\end{align*}
Par conséquent, le jugement commun est
\[
 \mathfrak e_\alpha:\operatorname{Ext}_k\ \vdash\
 \mathcal O_\bullet(\mathfrak e_\alpha)
 \in\mathsf Y_\bullet[\alpha],
 \qquad
 \bullet\in\{\mathrm{sp,sol,gau,coh,det}\}.
\]
\begin{align}
 \mathcal O_{\rm sp}(\mathfrak e_\alpha)
 &:=
 \bigl(U_\alpha,\sigma_x,\sigma_y,
 A_\alpha,\eta_\alpha\bigr),
 \label{eq:pdfv2-cor-accion-sp}\\
 \mathcal O_{\rm sol}(\mathfrak e_\alpha)
 &:=
 \bigl(b_\alpha,b_\alpha^+,b_\alpha^-;\,
 (M^+,M^-,D,H)\mapsto
 (M^++b_\alpha^+,M^-+b_\alpha^-,
  D+b_\alpha,H+|b_\alpha|)\bigr),
 \label{eq:pdfv2-cor-accion-sol}\\
 \mathcal O_{\rm gau}(\mathfrak e_\alpha)
 &:=
 \left\{
 \bigl(\operatorname{Cell}_{\rm TPK}(\alpha),
 q,B,W_c(q)\bigr):
 q\in Q_{\rm inc}(y)
 \right\},
 \label{eq:pdfv2-cor-accion-gau}\\
 \mathcal O_{\rm coh}(\mathfrak e_\alpha)
 &:=
 \bigl(I_y,J_y,\kappa_{y,N},\delta_y,C_{y,N},
 F_\alpha:C_{y,N}\longrightarrow C_{x,N}\bigr),
 \label{eq:pdfv2-cor-accion-coh}\\
 \mathcal O_{\rm det}(\mathfrak e_\alpha)
 &:=
 \bigl(\Psi_\alpha,K_\alpha,\mathbf1_\alpha,z_\pm(\alpha),
 h_*(\alpha),\operatorname{Det}(C_{y,N}),H_*(C_{y,N}),
 \operatorname{Smith}(C_{y,N}),\operatorname{Boc}(C_{y,N})\bigr).
 \label{eq:pdfv2-cor-accion-det}
\end{align}

La lecture archimédienne interne est la première composante de cette action simultanée ; elle n'introduit encore aucune forme classique :
\begin{equation}
\label{pII-full:eq:continuo-lector-arquimediano}
 \mathcal L_{\rm arq}^{\rm int}(\mathfrak e_\alpha)
 :=\mathcal O_{\rm sp}(\mathfrak e_\alpha).
\end{equation}
Pour toute valeur de ces systèmes, y compris chaque élément de la multisection \(\mathcal O_{\rm gau}\), son support est l'indice dépendant
\begin{equation}
\label{eq:pdfv2-cor-soporte-dependiente}
 \begin{aligned}
 \mathsf B_\alpha&:=
 \{(s(\alpha),t(\alpha),\gamma\alpha,
          \operatorname{Led}(\gamma\alpha))\},\\
 \operatorname{supp}_\alpha:\mathsf Y_\bullet[\alpha]&\longrightarrow
 \mathsf B_\alpha,\qquad
 o\longmapsto(s(\alpha),t(\alpha),\gamma\alpha,
          \operatorname{Led}(\gamma\alpha)).
 \end{aligned}
\end{equation}
Ainsi sont typés, sans projection reconstruisant l'arête,
\begin{equation}
\label{eq:pdfv2-cor-soporte-cinco-lecturas}
 \begin{gathered}
 s_\alpha(\mathcal O_\bullet):=s(\alpha),\qquad
 t_\alpha(\mathcal O_\bullet):=t(\alpha),\\
 \gamma_\alpha(\mathcal O_\bullet):=\gamma\alpha,
 \qquad
 \operatorname{Led}_\alpha(\mathcal O_\bullet)
 :=\operatorname{Led}(\gamma\alpha).
 \end{gathered}
\end{equation}
Ici, \(\operatorname{Cell}_{\rm TPK}(\alpha)\) est l'application de complexes \eqref{eq:pdfv2-cor-cell-accion-emision} induite par le transport de la même fibre de feuillets, et \(F_\alpha\) est \eqref{eq:pdfv2-cor-refinamiento-celda} pour cette extension. En particulier, \(\mathcal O_{\rm gau}\) conserve la multisection complète des \(q\) ; il ne choisit ni l'origine ni une orbite.

Les cinq formules précédentes constituent une seule action dépendante
\begin{equation}
\label{eq:pdfv2-cor-accion-unica-dependiente}
 \mathcal O(\mathfrak e_\alpha)
 :=\left\langle
 \mathcal O_{\rm sp},\mathcal O_{\rm sol},\mathcal O_{\rm gau},
 \mathcal O_{\rm coh},\mathcal O_{\rm det};
 \operatorname{Cpl}_\alpha
 \right\rangle,
\end{equation}
où \(\operatorname{Cpl}_\alpha\) est le certificat formé par les identités
\begin{equation}
\label{eq:pdfv2-cor-acoplamientos-unicos}
 \begin{gathered}
 \operatorname{supp}_\alpha(\mathcal O_{\rm sp})
 =\operatorname{supp}_\alpha(\mathcal O_{\rm sol})
 =\operatorname{supp}_\alpha(\mathcal O_{\rm gau})
 =\operatorname{supp}_\alpha(\mathcal O_{\rm coh})
 =\operatorname{supp}_\alpha(\mathcal O_{\rm det}),\\
 \mathsf Q(t(\alpha))=\sigma_9\mathsf Q(s(\alpha)),\quad
 g(t(\alpha))=g(s(\alpha))+[1]_9,\quad
 q^{(9)}(t(\alpha))=q^{(9)}(s(\alpha))+\chi_9(g(s(\alpha))),\\
 A_\alpha^*A_\alpha+\eta_\alpha^*\eta_\alpha
 =P_xU_\alpha^*U_\alpha P_x+Q_xU_\alpha^*U_\alpha Q_x,\\
 (M^+,M^-,D,H)'=(M^++b_\alpha^+,M^-+b_\alpha^-,
 D+b_\alpha,H+|b_\alpha|),\\
 Q_{\rm inc}=C^2(\operatorname{Cell}_{\rm TPK};\mathbf F_3),\qquad
 sB=0,\qquad W_c(q+Ba)=W_c(q),\\
 C_{y,N}\text{ est le même complexe dans }
 \mathcal O_{\rm coh}\text{ et }\mathcal O_{\rm det},\qquad
 \partial h_*=z_--z_+,\qquad
 K_{\beta\alpha}=K_\beta+\Psi_\beta K_\alpha.
 \end{gathered}
\end{equation}
\Needspace{18\baselineskip}
Le symbole \(\bullet\) parcourt les cinq lectures de \eqref{eq:pdfv2-cor-accion-unica-dependiente} ; toutes conservent les mêmes source, cible, parcours et registre.

Nous réunissons, sans séparer l'état, les opérations construites jusqu'ici :
\begin{equation}
\label{eq:pdfv2-cor-constructor-unico}
 \begin{aligned}
 \mathfrak D_{k,N}(\widehat x)
 &:=\bigl(\{\mathfrak e_\alpha\}_\alpha;\\
 &\quad U_k,\sigma_k,P_k,Q_k,A_k,\eta_k;\\
 &\quad b_\alpha,b_\alpha^\pm,M_k^\pm,D_k^0,H_k^0,
 \kappa_k^{\rm can};\\
 &\quad V_{\rm TPK},\mathscr I,\mathscr E,\mathscr B,B;\\
 &\quad \operatorname{Cell}_{\rm TPK}(\widehat x),Q_{\rm inc}(\widehat x),W_c;\\
 &\quad I_k,J_k,\kappa_{k,N},\delta,C_{k,N};\\
 &\quad c_\alpha,v_\alpha,\Psi_\alpha,K_\alpha,
 \mathbf1_\alpha,z_\pm(\alpha),h_*(\alpha);\\
 &\quad
 \operatorname{Det}(C_{k,N}),H_*(C_{k,N}),
 \operatorname{Smith}(C_{k,N}),\operatorname{Boc}(C_{k,N});\\
 &\quad \{\mathcal O(\mathfrak e_\alpha)\}_\alpha
 \bigr).
 \end{aligned}
\end{equation}
Le point-virgule ne sépare que des niveaux de lecture d'un même tuple dépendant ; il ne définit pas de structures indépendantes. L'état d'entrée n'est pas enregistré comme première coordonnée ni récupéré par une projection. Chaque opération se démontre par son action sur les générateurs \(\mathfrak e_\alpha\), qui conservent le même support, le même parcours et le même registre.

Pour \(\ell\geq k\) et \(N'\geq N\), le transport est défini composante par composante par troncature des histoires, réduction des coefficients et précomposition des ports. L'affirmation de naturalité à démontrer est unique :
\begin{equation}
\label{eq:pdfv2-cor-naturalidad-unica}
 u_{(\ell,N'),(k,N)}\mathfrak D_{\ell,N'}(\widehat x_\ell)
 =\mathfrak D_{k,N}(\tau_{\ell,k}\widehat x_\ell).
\end{equation}
La naturalité se formule sur l'histoire engendrée, sans prétendre qu'une arête fine isolée doive rester une arête après troncature de ses deux extrémités. Pour \(\gamma=(\varepsilon_1,\ldots,\varepsilon_\ell)\), nous définissons
\[
 \mathsf E_{\leq\ell}(\gamma)
 :=(\mathfrak e_{\varepsilon_i})_{1\leq i\leq\ell},
 \qquad
 \mathcal O_{\leq\ell}(\gamma)
 :=(\mathcal O(\mathfrak e_{\varepsilon_i}))_{1\leq i\leq\ell}.
\]
Les deux transports sont les applications explicites de préfixe
\begin{align}
 u^{\mathfrak e}_{\ell,k}:
 \mathsf E_{\leq\ell}(\gamma)&\longrightarrow
 \mathsf E_{\leq k}(\gamma),&
 (\mathfrak e_{\varepsilon_i})_{i\leq\ell}
 &\longmapsto(\mathfrak e_{\varepsilon_i})_{i\leq k},
 \label{eq:pdfv2-cor-transporte-arista}\\
 u^{\mathcal O}_{\ell,k}:
 \mathcal O_{\leq\ell}(\gamma)&\longrightarrow
 \mathcal O_{\leq k}(\gamma),&
 (\mathcal O(\mathfrak e_{\varepsilon_i}))_{i\leq\ell}
 &\longmapsto
 (\mathcal O(\mathfrak e_{\varepsilon_i}))_{i\leq k}.
 \label{eq:pdfv2-cor-transporte-accion}
\end{align}
Dans chaque entrée conservée, la composante \(\mathrm{sp}\) entrelace \(U,\sigma,P,Q,A,\eta\) ; la composante \(\mathrm{sol}\) applique l'espérance conditionnelle à \(b,b^\pm,\kappa\) ; la composante \(\mathrm{gau}\) applique \(\operatorname{Cell}(\tau)\oplus\tau_Q\) ; la composante \(\mathrm{coh}\) précompose les ports avec \(F\) ; et la composante \(\mathrm{det}\) applique \(\Psi,K,H_*,\operatorname{Smith},\operatorname{Boc}\) à la même application de complexes. Le préfixe agit simultanément sur ces cinq composantes. Ses composantes déjà établies sont
\begin{equation}
\label{eq:pdfv2-cor-naturalidad-componentes}
 \begin{gathered}
 U_{m,k}=U_{m,\ell}U_{\ell,k}\quad(m\geq\ell\geq k),\qquad
 E_kb_f^\pm=b_c^\pm+\kappa_k^{\rm can},\\
 r_{N',N}\kappa_{k,N'}=
 \kappa_{k,N}(r_{N',N}\oplus r_{N',N}),\\
 \Psi_\beta\Psi_\alpha=\Psi_{\beta\alpha},
 \qquad K_{\beta\alpha}=K_\beta+\Psi_\beta K_\alpha,\\
 u^{\mathcal O}_{\ell,k}\mathcal O_{\leq\ell}(\gamma)
 =\mathcal O_{\leq k}(\gamma)
 =\operatorname{map}(\mathcal O)\bigl(u^{\mathfrak e}_{\ell,k}
                    \mathsf E_{\leq\ell}(\gamma)\bigr).
 \end{gathered}
\end{equation}
La première égalité utilise la multiplication des poids conditionnels sur le même système de cylindres communs ; les autres sont les formules explicites précédentes. Aucune projection visible ne peut servir à reconstruire ou à sélectionner \(\widehat x\).

Pour compléter la dernière composante, notons \(F_{(\ell,N'),(k,N)}:C_{\ell,N'}\to C_{k,N}\) l'application de complexes de \eqref{eq:pdfv2-cor-refinamiento-celda}. Son cône est conservé comme partie de la transition. La fonctorialité de l'homologie et de la suite de Bockstein donne
\begin{equation}
\label{eq:pdfv2-cor-naturalidad-homologia-bockstein}
 H_*(F_{m,k})=H_*(F_{\ell,k})H_*(F_{m,\ell}),
 \qquad
 \operatorname{Boc}(F_{m,k})
 =\operatorname{Boc}(F_{\ell,k})
  \operatorname{Boc}(F_{m,\ell}).
\end{equation}
La droite déterminant est transportée au moyen de l'identité du triangle exact
\begin{equation}
\label{eq:pdfv2-cor-naturalidad-determinante}
 \operatorname{Det}(C_{k,N})
 \simeq
 \operatorname{Det}(C_{\ell,N'})\otimes
 \operatorname{Det}\!\left(\operatorname{Cone}
 F_{(\ell,N'),(k,N)}\right).
\end{equation}
Cette formule enregistre également les raffinements qui ajoutent de l'homologie. La lecture de Smith est conservée sans feindre l'invariance : à chaque transition est associé le certificat total
\begin{equation}
\label{eq:pdfv2-cor-transicion-smith}
 \mathcal S(F):=
 \bigl(\operatorname{Smith}(C_{\ell,N'}),
       \operatorname{Smith}(C_{k,N}),
       \operatorname{Smith}(\operatorname{Cone}F)\bigr),
\end{equation}
qui détermine littéralement quels facteurs se conservent, apparaissent ou disparaissent.

L'action d'une symétrie cohérente \(g\) sur l'arbre commun induit des isométries de permutation \(\widetilde g_k\) dans \(\mathscr H_k\). La bijection des successeurs et la conservation des poids de \eqref{eq:continuo-comun-naturalidad-medida} prouvent sur les vecteurs de base
\begin{equation}
\label{eq:pdfv2-cor-naturalidad-transporte-hojas}
 \widetilde g_{k+1}U_k=U_k\widetilde g_k,
 \qquad
 \widetilde g_k\sigma_k=\sigma_k\widetilde g_k,
 \qquad
 \widetilde g_kP_k=P_k\widetilde g_k,
 \qquad
 \widetilde g_kQ_k=Q_k\widetilde g_k.
\end{equation}

\begin{theorem}[Totalité et naturalité du constructeur corrélatif]
\label{thm:pdfv2-cor-totalidad-naturalidad}
Pour tout état total de \eqref{eq:continuo-comun-estado}, tout horizon fini et tout niveau de coefficients, l'expression \eqref{eq:pdfv2-cor-constructor-unico} est définie. Ses cinq caractéristiques sont des opérations simultanées sur le même registre d'extension et satisfont l'unique identité \eqref{eq:pdfv2-cor-naturalidad-unica}. Le collier, le complexe et la droite déterminant restent non vides même lorsque leurs lecteurs ultérieurs ne sont ni acycliques ni inversibles.
\end{theorem}

\begin{proof}
La non-vacuité de l'arbre et de ses ports est \cref{prop:continuo-comun-arbol-finito-no-vacio} ; la mesure et l'espérance sont totales par \cref{thm:continuo-comun-medida-canonica}. Les formules \eqref{eq:pdfv2-cor-u-comun}--\eqref{eq:pdfv2-cor-memorias} définissent transport, polarité, bilan et mémoire pour tous les générateurs. Les quatre directions de \eqref{eq:pdfv2-cor-cuatro-rectas-tpk} sont distinctes par \eqref{eq:pdfv2-cor-proyectores-no-colapso} ; la complétion cellulaire \cref{def:pdfv2-cor-completacion-celular-tpk} contient leurs deux ordres de chemin et six cellules, tandis que \eqref{eq:pdfv2-cor-collar-naturalidad} prouve sa naturalité. La suite cellulaire est exacte par \eqref{eq:pdfv2-cor-sucesion-celda}, le complexe local satisfait \(\partial^2=0\) et le remplissage corrélatif est explicite par \eqref{eq:pdfv2-cor-filler}. Toute construction finie munie de bases possède la droite \eqref{eq:pdfv2-cor-linea-determinante-total} ; homologie, Bockstein et certificat de Smith sont transportés par \eqref{eq:pdfv2-cor-naturalidad-homologia-bockstein}--\eqref{eq:pdfv2-cor-transicion-smith}. Enfin, les identités se vérifient sur chaque générateur \(\mathfrak e_\alpha\) au moyen de \eqref{eq:pdfv2-cor-naturalidad-componentes}, \eqref{eq:pdfv2-cor-collar-naturalidad} et \eqref{eq:pdfv2-cor-naturalidad-transporte-hojas} ; par linéarité et composition fonctorielle des parcours, on obtient \eqref{eq:pdfv2-cor-naturalidad-unica} pour le registre entier.
\end{proof}

\section{Distinction causale entre identités, déductions et réalisations du constructeur corrélatif}

\paragraph{Non-vacuité et identités exactes du constructeur corrélatif}
La non-vacuité importée de \eqref{eq:continuo-comun-out-finito} et les identités \eqref{eq:pdfv2-cor-b-identidades}, \eqref{eq:pdfv2-cor-sucesion-celda}, \eqref{eq:pdfv2-cor-sb-cero}, \eqref{eq:pdfv2-cor-u-gram}, \eqref{eq:pdfv2-cor-balance-hojas}, \eqref{eq:pdfv2-cor-jensen-hojas}, \eqref{eq:pdfv2-cor-dos-cero}, \eqref{eq:pdfv2-cor-cociclo-matricial} et \eqref{eq:pdfv2-cor-filler} se déduisent des définitions et des relations APP--TRIT--TPK exposées. Leur statut est \texttt{EXACT\_INTERNE}.

\paragraph{Hypothèses initiales pour l'isométrie de \(U_k\), la suite cellulaire et la composition \(\Psi/K\)}
La construction de \(U_k\) prend pour structure initiale la finitude des fibres, la survie à tout horizon, la troncature fonctionnelle et les transports des deux feuillets déjà produits par TPK\@. L'égalité \(U_k^*U_k=I\) équivaut à vérifier le Gram moyen normalisé de \eqref{eq:pdfv2-cor-gram-promedio-normalizado} ; il suffit, en particulier, de vérifier \(U_\alpha^*U_\alpha=I\) pour chaque émission. L'identité générale \eqref{eq:pdfv2-cor-u-gram} est conservée pour tous les transports de feuillet déclarés. La suite cellulaire utilise les deux ports APP ordonnés. La composition \(\Psi/K\) utilise la retenue entière, l'orientation et la frontière. Sous ces prémisses énumérées, les conclusions indiquées sont affirmatives.

\paragraph{Données supplémentaires pour les réalisations isométriques, les rétractes équilibrés et la trivialisation de la droite déterminant}
Le constructeur interne est complet avec le Gram positif de \eqref{eq:pdfv2-cor-u-gram}, l'espace total \(Q_{\rm inc}\), la suite cellulaire et la droite déterminant munie de son cône de transition. Une expression spécialisée peut exiger des données supplémentaires : l'isométrie de chaque \(U_\alpha\), une identification numérique concrète du défaut APP, un rétracte équilibré et acyclique ou la trivialisation d'une droite déterminant. Ces données sont des entrelaceurs du codomaine exprimé ; elles ne font partie ni du domaine ni de l'existence de \(\mathfrak D_{k,N}\). Leur absence empêche exclusivement la promotion qui les mentionne.

\begin{criterion}[Falsificateurs de la corrélation intrinsèque]
\label{crit:pdfv2-cor-falsadores}
La construction conjointe échoue au premier niveau où survient l'un des faits suivants :
\begin{enumerate}
 \item une opération utilise une extension, une orientation, une retenue, une mémoire, une frontière ou un parcours différent de celui enregistré dans \(\mathfrak e_\alpha\) ;
 \item un successeur de \(\operatorname{Ch}_k(\widehat x_k)\) est choisi et les autres sont supprimés ;
 \item les poids de \eqref{eq:pdfv2-cor-peso-conteo} ne somment pas à un, \eqref{eq:pdfv2-cor-u-gram} échoue, ou \(U_k^*U_k=I\) est déclaré sans vérifier le Gram moyen normalisé de \eqref{eq:pdfv2-cor-gram-promedio-normalizado}, par exemple au moyen de l'isométrie de chaque transport d'émission ;
 \item \(\kappa_k^{\rm can}\) est identifié à la retenue entière \(\kappa_\alpha\) ;
 \item \(sB=0\) échoue, \(Q_{\rm inc}(\widehat x)\) est contracté à son origine, l'action \(q\mapsto q+Ba\) est modifiée, ou les deux ordres de chemin sont identifiés avant l'application du lecteur ;
 \item un raffinement rompt \eqref{eq:pdfv2-cor-refinamiento-cadena} ou le cocycle \eqref{eq:pdfv2-cor-cociclo-matricial} ;
 \item \(\partial^2=0\) échoue, \(z_\pm\) ne sont pas des cycles ou \(\partial h_*\neq z_--z_+\) ;
 \item un scalaire déterminantal non nul est affirmé sans le rétracte, le bilan et l'acyclicité exigés par cette expression ;
 \item l'une quelconque des cinq opérations est présentée comme un objet séparé du constructeur \(\mathfrak D_{k,N}\).
\end{enumerate}
Dans tous ces cas, le diagnostic est : \emph{objet commun remplacé ; conclusion non transférable}.
\end{criterion}
